\documentclass{openreader}
\title{Πρὸς Ῥωμαίους}
\date{}
\setmainlanguage{english}
\setmainfont{Charis SIL}
\setotherlanguage{greek}
\newfontfamily\greekfont[Script=Greek]{SBL BibLit}
\begin{document}
\ORBselectlanguage{greek}
\maketitle
\raggedbottom 
\fontsize{16pt}{24pt}\selectfont

\ChapterHeader{1}{Chapter 1}

\MainTextVerseMark{1}{1}Παῦλος δοῦλος ⸂Χριστοῦ Ἰησοῦ⸃, κλητὸς\FnParseFormGloss{a}{nominative masculine singular}{κλητός}{called} ἀπόστολος, ἀφωρισμένος\FnParseFormGloss{b}{perfect passive participle nominative masculine singular}{ἀφορίζω}{to separate} εἰς εὐαγγέλιον θεοῦ 
\MainTextVerseMark{1}{2}ὃ προεπηγγείλατο\FnParseFormGloss{a}{aorist middle indicative 3rd singular}{προεπαγγέλλομαι}{to promise beforehand} διὰ τῶν προφητῶν αὐτοῦ ἐν γραφαῖς ἁγίαις 
\MainTextVerseMark{1}{3}περὶ τοῦ υἱοῦ αὐτοῦ, τοῦ γενομένου\FnParse{a}{aorist middle participle genitive masculine singular} ἐκ σπέρματος Δαυὶδ κατὰ σάρκα, 
\MainTextVerseMark{1}{4}τοῦ ὁρισθέντος\FnParseFormGloss{a}{aorist passive participle genitive masculine singular}{ὁρίζω}{to determine} υἱοῦ θεοῦ ἐν δυνάμει κατὰ πνεῦμα ἁγιωσύνης\FnParseFormGloss{b}{genitive feminine singular}{ἁγιωσύνη}{holiness} ἐξ ἀναστάσεως νεκρῶν, Ἰησοῦ Χριστοῦ τοῦ κυρίου ἡμῶν, 
\MainTextVerseMark{1}{5}δι’ οὗ ἐλάβομεν\FnParse{a}{aorist active indicative 1st plural} χάριν καὶ ἀποστολὴν\FnParseFormGloss{b}{accusative feminine singular}{ἀποστολή}{apostleship} εἰς ὑπακοὴν\FnParseFormGloss{c}{accusative feminine singular}{ὑπακοή}{obedience} πίστεως ἐν πᾶσιν τοῖς ἔθνεσιν ὑπὲρ τοῦ ὀνόματος αὐτοῦ, 
\MainTextVerseMark{1}{6}ἐν οἷς ἐστε\FnParse{a}{present active indicative 2nd plural} καὶ ὑμεῖς κλητοὶ\FnParseFormGloss{b}{nominative masculine plural}{κλητός}{called} Ἰησοῦ Χριστοῦ, 
\MainTextVerseMark{1}{7}πᾶσιν τοῖς οὖσιν\FnParse{a}{present active participle dative masculine plural} ἐν Ῥώμῃ\FnParseFormGloss{b}{dative feminine singular}{Ῥώμη}{Rome} ἀγαπητοῖς θεοῦ, κλητοῖς\FnParseFormGloss{c}{dative masculine plural}{κλητός}{called} ἁγίοις· χάρις ὑμῖν καὶ εἰρήνη ἀπὸ θεοῦ πατρὸς ἡμῶν καὶ κυρίου Ἰησοῦ Χριστοῦ. 
\MainTextVerseMark{1}{8}Πρῶτον μὲν εὐχαριστῶ\FnParse{a}{present active indicative 1st singular} τῷ θεῷ μου διὰ Ἰησοῦ Χριστοῦ ⸀περὶ πάντων ὑμῶν, ὅτι ἡ πίστις ὑμῶν καταγγέλλεται\FnParseFormGloss{b}{present passive indicative 3rd singular}{καταγγέλλω}{to preach} ἐν ὅλῳ τῷ κόσμῳ. 
\MainTextVerseMark{1}{9}μάρτυς γάρ μού ἐστιν\FnParse{a}{present active indicative 3rd singular} ὁ θεός, ᾧ λατρεύω\FnParseFormGloss{b}{present active indicative 1st singular}{λατρεύω}{to serve} ἐν τῷ πνεύματί μου ἐν τῷ εὐαγγελίῳ τοῦ υἱοῦ αὐτοῦ, ὡς ἀδιαλείπτως\FnFormGloss{c}{ἀδιαλείπτως}{unceasingly} μνείαν\FnParseFormGloss{d}{accusative feminine singular}{μνεία}{remembrance} ὑμῶν ποιοῦμαι\FnParse{e}{present middle indicative 1st singular} 
\MainTextVerseMark{1}{10}πάντοτε ἐπὶ τῶν προσευχῶν μου, δεόμενος\FnParse{a}{present middle participle nominative masculine singular} εἴ πως\FnFormGloss{b}{πώς}{in any way} ἤδη ποτὲ\FnFormGloss{c}{ποτέ}{once} εὐοδωθήσομαι\FnParseFormGloss{d}{future passive indicative 1st singular}{εὐοδόομαι}{to get along with} ἐν τῷ θελήματι τοῦ θεοῦ ἐλθεῖν\FnParse{e}{aorist active infinitive} πρὸς ὑμᾶς. 
\MainTextVerseMark{1}{11}ἐπιποθῶ\FnParseFormGloss{a}{present active indicative 1st singular}{ἐπιποθέω}{to long for} γὰρ ἰδεῖν\FnParse{b}{aorist active infinitive} ὑμᾶς, ἵνα τι μεταδῶ\FnParseFormGloss{c}{aorist active subjunctive 1st singular}{μεταδίδωμι}{to impart} χάρισμα\FnParseFormGloss{d}{accusative neuter singular}{χάρισμα}{gracious gift} ὑμῖν πνευματικὸν\FnParseFormGloss{e}{accusative neuter singular}{πνευματικός}{spiritual} εἰς τὸ στηριχθῆναι\FnParseFormGloss{f}{aorist passive infinitive}{στηρίζω}{to strengthen} ὑμᾶς, 
\MainTextVerseMark{1}{12}τοῦτο δέ ἐστιν\FnParse{a}{present active indicative 3rd singular} συμπαρακληθῆναι\FnParseFormGloss{b}{aorist passive infinitive}{συμπαρακαλέομαι}{to be mutually encouraged} ἐν ὑμῖν διὰ τῆς ἐν ἀλλήλοις πίστεως ὑμῶν τε καὶ ἐμοῦ. 
\MainTextVerseMark{1}{13}οὐ θέλω\FnParse{a}{present active indicative 1st singular} δὲ ὑμᾶς ἀγνοεῖν,\FnParseFormGloss{b}{present active infinitive}{ἀγνοέω}{to be ignorant} ἀδελφοί, ὅτι πολλάκις\FnFormGloss{c}{πολλάκις}{many times} προεθέμην\FnParseFormGloss{d}{aorist middle indicative 1st singular}{προτίθεμαι}{to plan} ἐλθεῖν\FnParse{e}{aorist active infinitive} πρὸς ὑμᾶς, καὶ ἐκωλύθην\FnParseFormGloss{f}{aorist passive indicative 1st singular}{κωλύω}{to hinder} ἄχρι τοῦ δεῦρο,\FnFormGloss{g}{δεῦρο}{come} ἵνα τινὰ καρπὸν σχῶ\FnParse{h}{aorist active subjunctive 1st singular} καὶ ἐν ὑμῖν καθὼς καὶ ἐν τοῖς λοιποῖς ἔθνεσιν. 
\MainTextVerseMark{1}{14}Ἕλλησίν\FnParseFormGloss{a}{dative masculine plural}{Ἕλλην}{Greek} τε καὶ βαρβάροις,\FnParseFormGloss{b}{dative masculine plural}{βάρβαρος}{non-Greek} σοφοῖς\FnParseFormGloss{c}{dative masculine plural}{σοφός}{wise} τε καὶ ἀνοήτοις\FnParseFormGloss{d}{dative masculine plural}{ἀνόητος}{foolish} ὀφειλέτης\FnParseFormGloss{e}{nominative masculine singular}{ὀφειλέτης}{debtor} εἰμί·\FnParse{f}{present active indicative 1st singular} 
\MainTextVerseMark{1}{15}οὕτως τὸ κατ’ ἐμὲ πρόθυμον\FnParseFormGloss{a}{nominative neuter singular}{πρόθυμος}{willing} καὶ ὑμῖν τοῖς ἐν Ῥώμῃ\FnParseFormGloss{b}{dative feminine singular}{Ῥώμη}{Rome} εὐαγγελίσασθαι.\FnParse{c}{aorist middle infinitive} 
\MainTextVerseMark{1}{16}Οὐ γὰρ ἐπαισχύνομαι\FnParseFormGloss{a}{present middle indicative 1st singular}{ἐπαισχύνομαι}{to be ashamed of} τὸ ⸀εὐαγγέλιον, δύναμις γὰρ θεοῦ ἐστιν\FnParse{b}{present active indicative 3rd singular} εἰς σωτηρίαν παντὶ τῷ πιστεύοντι,\FnParse{c}{present active participle dative masculine singular} Ἰουδαίῳ τε πρῶτον καὶ Ἕλληνι·\FnParseFormGloss{d}{dative masculine singular}{Ἕλλην}{Greek} 
\MainTextVerseMark{1}{17}δικαιοσύνη γὰρ θεοῦ ἐν αὐτῷ ἀποκαλύπτεται\FnParseFormGloss{a}{present passive indicative 3rd singular}{ἀποκαλύπτω}{to reveal} ἐκ πίστεως εἰς πίστιν, καθὼς γέγραπται·\FnParse{b}{perfect passive indicative 3rd singular} Ὁ δὲ δίκαιος ἐκ πίστεως ζήσεται.\FnParse{c}{future middle indicative 3rd singular} 
\MainTextVerseMark{1}{18}Ἀποκαλύπτεται\FnParseFormGloss{a}{present passive indicative 3rd singular}{ἀποκαλύπτω}{to reveal} γὰρ ὀργὴ θεοῦ ἀπ’ οὐρανοῦ ἐπὶ πᾶσαν ἀσέβειαν\FnParseFormGloss{b}{accusative feminine singular}{ἀσέβεια}{ungodliness} καὶ ἀδικίαν\FnParseFormGloss{c}{accusative feminine singular}{ἀδικία}{wickedness} ἀνθρώπων τῶν τὴν ἀλήθειαν ἐν ἀδικίᾳ\FnParseFormGloss{d}{dative feminine singular}{ἀδικία}{wickedness} κατεχόντων,\FnParseFormGloss{e}{present active participle genitive masculine plural}{κατέχω}{to hold back} 
\MainTextVerseMark{1}{19}διότι\FnFormGloss{a}{διότι}{therefore} τὸ γνωστὸν\FnParseFormGloss{b}{nominative neuter singular}{γνωστός}{known} τοῦ θεοῦ φανερόν\FnParseFormGloss{c}{nominative neuter singular}{φανερός}{visible} ἐστιν\FnParse{d}{present active indicative 3rd singular} ἐν αὐτοῖς, ὁ ⸂θεὸς γὰρ⸃ αὐτοῖς ἐφανέρωσεν.\FnParse{e}{aorist active indicative 3rd singular} 
\MainTextVerseMark{1}{20}τὰ γὰρ ἀόρατα\FnParseFormGloss{a}{nominative neuter plural}{ἀόρατος}{invisible} αὐτοῦ ἀπὸ κτίσεως\FnParseFormGloss{b}{genitive feminine singular}{κτίσις}{creation} κόσμου τοῖς ποιήμασιν\FnParseFormGloss{c}{dative neuter plural}{ποίημα}{what is made} νοούμενα\FnParseFormGloss{d}{present passive participle nominative neuter plural}{νοέω}{to understand} καθορᾶται,\FnParseFormGloss{e}{present passive indicative 3rd singular}{καθοράω}{to be clearly seen} ἥ τε ἀΐδιος\FnParseFormGloss{f}{nominative feminine singular}{ἀΐδιος}{eternal} αὐτοῦ δύναμις καὶ θειότης,\FnParseFormGloss{g}{nominative feminine singular}{θειότης}{divine nature} εἰς τὸ εἶναι\FnParse{h}{present active infinitive} αὐτοὺς ἀναπολογήτους,\FnParseFormGloss{i}{accusative masculine plural}{ἀναπολόγητος}{without excuse} 
\MainTextVerseMark{1}{21}διότι\FnFormGloss{a}{διότι}{therefore} γνόντες\FnParse{b}{aorist active participle nominative masculine plural} τὸν θεὸν οὐχ ὡς θεὸν ἐδόξασαν\FnParse{c}{aorist active indicative 3rd plural} ἢ ηὐχαρίστησαν,\FnParse{d}{aorist active indicative 3rd plural} ἀλλὰ ἐματαιώθησαν\FnParseFormGloss{e}{aorist passive indicative 3rd plural}{ματαιόομαι}{to become futile} ἐν τοῖς διαλογισμοῖς\FnParseFormGloss{f}{dative masculine plural}{διαλογισμός}{thought} αὐτῶν καὶ ἐσκοτίσθη\FnParseFormGloss{g}{aorist passive indicative 3rd singular}{σκοτίζομαι}{to be dark} ἡ ἀσύνετος\FnParseFormGloss{h}{nominative feminine singular}{ἀσύνετος}{senseless} αὐτῶν καρδία· 
\MainTextVerseMark{1}{22}φάσκοντες\FnParseFormGloss{a}{present active participle nominative masculine plural}{φάσκω}{to claim} εἶναι\FnParse{b}{present active infinitive} σοφοὶ\FnParseFormGloss{c}{nominative masculine plural}{σοφός}{wise} ἐμωράνθησαν,\FnParseFormGloss{d}{aorist passive indicative 3rd plural}{μωραίνω}{to make foolish} 
\MainTextVerseMark{1}{23}καὶ ἤλλαξαν\FnParseFormGloss{a}{aorist active indicative 3rd plural}{ἀλλάσσω}{to change} τὴν δόξαν τοῦ ἀφθάρτου\FnParseFormGloss{b}{genitive masculine singular}{ἄφθαρτος}{imperishable} θεοῦ ἐν ὁμοιώματι\FnParseFormGloss{c}{dative neuter singular}{ὁμοίωμα}{likeness} εἰκόνος\FnParseFormGloss{d}{genitive feminine singular}{εἰκών}{image} φθαρτοῦ\FnParseFormGloss{e}{genitive masculine singular}{φθαρτός}{perishable} ἀνθρώπου καὶ πετεινῶν\FnParseFormGloss{f}{genitive neuter plural}{πετεινόν}{bird} καὶ τετραπόδων\FnParseFormGloss{g}{genitive neuter plural}{τετράπους}{four-footed} καὶ ἑρπετῶν.\FnParseFormGloss{h}{genitive neuter plural}{ἑρπετόν}{reptile} 
\MainTextVerseMark{1}{24}⸀Διὸ παρέδωκεν\FnParse{a}{aorist active indicative 3rd singular} αὐτοὺς ὁ θεὸς ἐν ταῖς ἐπιθυμίαις τῶν καρδιῶν αὐτῶν εἰς ἀκαθαρσίαν\FnParseFormGloss{b}{accusative feminine singular}{ἀκαθαρσία}{impurity} τοῦ ἀτιμάζεσθαι\FnParseFormGloss{c}{present middle infinitive}{ἀτιμάζω}{to dishonor} τὰ σώματα αὐτῶν ἐν ⸀αὐτοῖς, 
\MainTextVerseMark{1}{25}οἵτινες μετήλλαξαν\FnParseFormGloss{a}{aorist active indicative 3rd plural}{μεταλλάσσω}{to exchange} τὴν ἀλήθειαν τοῦ θεοῦ ἐν τῷ ψεύδει,\FnParseFormGloss{b}{dative neuter singular}{ψεῦδος}{lie} καὶ ἐσεβάσθησαν\FnParseFormGloss{c}{aorist passive indicative 3rd plural}{σεβάζομαι}{to worship} καὶ ἐλάτρευσαν\FnParseFormGloss{d}{aorist active indicative 3rd plural}{λατρεύω}{to serve} τῇ κτίσει\FnParseFormGloss{e}{dative feminine singular}{κτίσις}{creation} παρὰ τὸν κτίσαντα,\FnParseFormGloss{f}{aorist active participle accusative masculine singular}{κτίζω}{to create} ὅς ἐστιν\FnParse{g}{present active indicative 3rd singular} εὐλογητὸς\FnParseFormGloss{h}{nominative masculine singular}{εὐλογητός}{worthy of being praised} εἰς τοὺς αἰῶνας· ἀμήν. 
\MainTextVerseMark{1}{26}Διὰ τοῦτο παρέδωκεν\FnParse{a}{aorist active indicative 3rd singular} αὐτοὺς ὁ θεὸς εἰς πάθη\FnParseFormGloss{b}{accusative neuter plural}{πάθος}{lust} ἀτιμίας·\FnParseFormGloss{c}{genitive feminine singular}{ἀτιμία}{dishonor} αἵ τε γὰρ θήλειαι\FnParseFormGloss{d}{nominative feminine plural}{θῆλυς}{female} αὐτῶν μετήλλαξαν\FnParseFormGloss{e}{aorist active indicative 3rd plural}{μεταλλάσσω}{to exchange} τὴν φυσικὴν\FnParseFormGloss{f}{accusative feminine singular}{φυσικός}{natural} χρῆσιν\FnParseFormGloss{g}{accusative feminine singular}{χρῆσις}{relations} εἰς τὴν παρὰ φύσιν,\FnParseFormGloss{h}{accusative feminine singular}{φύσις}{nature} 
\MainTextVerseMark{1}{27}ὁμοίως τε καὶ οἱ ἄρσενες\FnParseFormGloss{a}{nominative masculine plural}{ἄρσην}{male} ἀφέντες\FnParse{b}{aorist active participle nominative masculine plural} τὴν φυσικὴν\FnParseFormGloss{c}{accusative feminine singular}{φυσικός}{natural} χρῆσιν\FnParseFormGloss{d}{accusative feminine singular}{χρῆσις}{relations} τῆς θηλείας\FnParseFormGloss{e}{genitive feminine singular}{θῆλυς}{female} ἐξεκαύθησαν\FnParseFormGloss{f}{aorist passive indicative 3rd plural}{ἐκκαίομαι}{to be inflamed} ἐν τῇ ὀρέξει\FnParseFormGloss{g}{dative feminine singular}{ὄρεξις}{lust} αὐτῶν εἰς ἀλλήλους, ἄρσενες\FnParseFormGloss{h}{nominative masculine plural}{ἄρσην}{male} ἐν ἄρσεσιν\FnParseFormGloss{i}{dative masculine plural}{ἄρσην}{male} τὴν ἀσχημοσύνην\FnParseFormGloss{j}{accusative feminine singular}{ἀσχημοσύνη}{indecent act} κατεργαζόμενοι\FnParseFormGloss{k}{present middle participle nominative masculine plural}{κατεργάζομαι}{to produce} καὶ τὴν ἀντιμισθίαν\FnParseFormGloss{l}{accusative feminine singular}{ἀντιμισθία}{an exchange} ἣν ἔδει\FnParse{m}{imperfect active indicative 3rd singular} τῆς πλάνης\FnParseFormGloss{n}{genitive feminine singular}{πλάνη}{error} αὐτῶν ἐν ⸀ἑαυτοῖς ἀπολαμβάνοντες.\FnParseFormGloss{o}{present active participle nominative masculine plural}{ἀπολαμβάνω}{to receive} 
\MainTextVerseMark{1}{28}Καὶ καθὼς οὐκ ἐδοκίμασαν\FnParseFormGloss{a}{aorist active indicative 3rd plural}{δοκιμάζω}{to test} τὸν θεὸν ἔχειν\FnParse{b}{present active infinitive} ἐν ἐπιγνώσει,\FnParseFormGloss{c}{dative feminine singular}{ἐπίγνωσις}{knowledge} παρέδωκεν\FnParse{d}{aorist active indicative 3rd singular} αὐτοὺς ὁ θεὸς εἰς ἀδόκιμον\FnParseFormGloss{e}{accusative masculine singular}{ἀδόκιμος}{failing the test} νοῦν,\FnParseFormGloss{f}{accusative masculine singular}{νοῦς}{mind} ποιεῖν\FnParse{g}{present active infinitive} τὰ μὴ καθήκοντα,\FnParseFormGloss{h}{present active participle accusative neuter plural}{καθήκω}{to be fitting} 
\MainTextVerseMark{1}{29}πεπληρωμένους\FnParse{a}{perfect passive participle accusative masculine plural} πάσῃ ⸀ἀδικίᾳ\FnParseFormGloss{b}{dative feminine singular}{ἀδικία}{wickedness} πονηρίᾳ\FnParseFormGloss{c}{dative feminine singular}{πονηρία}{evil} πλεονεξίᾳ\FnParseFormGloss{d}{dative feminine singular}{πλεονεξία}{greediness} κακίᾳ,\FnParseFormGloss{e}{dative feminine singular}{κακία}{evil} μεστοὺς\FnParseFormGloss{f}{accusative masculine plural}{μεστός}{full} φθόνου\FnParseFormGloss{g}{genitive masculine singular}{φθόνος}{envy} φόνου\FnParseFormGloss{h}{genitive masculine singular}{φόνος}{murder} ἔριδος\FnParseFormGloss{i}{genitive feminine singular}{ἔρις}{quarrel} δόλου\FnParseFormGloss{j}{genitive masculine singular}{δόλος}{deceit} κακοηθείας,\FnParseFormGloss{k}{genitive feminine singular}{κακοήθεια}{malice} ψιθυριστάς,\FnParseFormGloss{l}{accusative masculine plural}{ψιθυριστής}{gossip} 
\MainTextVerseMark{1}{30}καταλάλους,\FnParseFormGloss{a}{accusative masculine plural}{κατάλαλος}{slanderous} θεοστυγεῖς,\FnParseFormGloss{b}{accusative masculine plural}{θεοστυγής}{God-hating} ὑβριστάς,\FnParseFormGloss{c}{accusative masculine plural}{ὑβριστής}{insolent man} ὑπερηφάνους,\FnParseFormGloss{d}{accusative masculine plural}{ὑπερήφανος}{proud} ἀλαζόνας,\FnParseFormGloss{e}{accusative masculine plural}{ἀλαζών}{boaster} ἐφευρετὰς\FnParseFormGloss{f}{accusative masculine plural}{ἐφευρετής}{inventor} κακῶν, γονεῦσιν\FnParseFormGloss{g}{dative masculine plural}{γονεύς}{parents} ἀπειθεῖς,\FnParseFormGloss{h}{accusative masculine plural}{ἀπειθής}{disobedient} 
\MainTextVerseMark{1}{31}ἀσυνέτους,\FnParseFormGloss{a}{accusative masculine plural}{ἀσύνετος}{senseless} ἀσυνθέτους,\FnParseFormGloss{b}{accusative masculine plural}{ἀσύνθετος}{faithless} ⸀ἀστόργους,\FnParseFormGloss{c}{accusative masculine plural}{ἄστοργος}{without love} ἀνελεήμονας·\FnParseFormGloss{d}{accusative masculine plural}{ἀνελεήμων}{ruthless} 
\MainTextVerseMark{1}{32}οἵτινες τὸ δικαίωμα\FnParseFormGloss{a}{accusative neuter singular}{δικαίωμα}{regulation} τοῦ θεοῦ ἐπιγνόντες,\FnParse{b}{aorist active participle nominative masculine plural} ὅτι οἱ τὰ τοιαῦτα πράσσοντες\FnParse{c}{present active participle nominative masculine plural} ἄξιοι θανάτου εἰσίν,\FnParse{d}{present active indicative 3rd plural} οὐ μόνον αὐτὰ ποιοῦσιν\FnParse{e}{present active indicative 3rd plural} ἀλλὰ καὶ συνευδοκοῦσιν\FnParseFormGloss{f}{present active indicative 3rd plural}{συνευδοκέω}{to approve of} τοῖς πράσσουσιν.\FnParse{g}{present active participle dative masculine plural} 
\ChapterHeader{2}{Chapter 2}

\MainTextVerseMark{2}{1}Διὸ ἀναπολόγητος\FnParseFormGloss{a}{nominative masculine singular}{ἀναπολόγητος}{without excuse} εἶ,\FnParse{b}{present active indicative 2nd singular} ὦ\FnFormGloss{c}{ὦ}{O!} ἄνθρωπε πᾶς ὁ κρίνων·\FnParse{d}{present active participle nominative masculine singular} ἐν ᾧ γὰρ κρίνεις\FnParse{e}{present active indicative 2nd singular} τὸν ἕτερον, σεαυτὸν κατακρίνεις,\FnParseFormGloss{f}{present active indicative 2nd singular}{κατακρίνω}{to condemn} τὰ γὰρ αὐτὰ πράσσεις\FnParse{g}{present active indicative 2nd singular} ὁ κρίνων·\FnParse{h}{present active participle nominative masculine singular} 
\MainTextVerseMark{2}{2}οἴδαμεν\FnParse{a}{perfect active indicative 1st plural} δὲ ὅτι τὸ κρίμα\FnParseFormGloss{b}{nominative neuter singular}{κρίμα}{judgment} τοῦ θεοῦ ἐστιν\FnParse{c}{present active indicative 3rd singular} κατὰ ἀλήθειαν ἐπὶ τοὺς τὰ τοιαῦτα πράσσοντας.\FnParse{d}{present active participle accusative masculine plural} 
\MainTextVerseMark{2}{3}λογίζῃ\FnParse{a}{present middle indicative 2nd singular} δὲ τοῦτο, ὦ\FnFormGloss{b}{ὦ}{O!} ἄνθρωπε ὁ κρίνων\FnParse{c}{present active participle nominative masculine singular} τοὺς τὰ τοιαῦτα πράσσοντας\FnParse{d}{present active participle accusative masculine plural} καὶ ποιῶν\FnParse{e}{present active participle nominative masculine singular} αὐτά, ὅτι σὺ ἐκφεύξῃ\FnParseFormGloss{f}{future middle indicative 2nd singular}{ἐκφεύγω}{to escape} τὸ κρίμα\FnParseFormGloss{g}{accusative neuter singular}{κρίμα}{judgment} τοῦ θεοῦ; 
\MainTextVerseMark{2}{4}ἢ τοῦ πλούτου\FnParseFormGloss{a}{genitive masculine singular}{πλοῦτος}{riches} τῆς χρηστότητος\FnParseFormGloss{b}{genitive feminine singular}{χρηστότης}{kindness} αὐτοῦ καὶ τῆς ἀνοχῆς\FnParseFormGloss{c}{genitive feminine singular}{ἀνοχή}{tolerance} καὶ τῆς μακροθυμίας\FnParseFormGloss{d}{genitive feminine singular}{μακροθυμία}{patience} καταφρονεῖς,\FnParseFormGloss{e}{present active indicative 2nd singular}{καταφρονέω}{to despise} ἀγνοῶν\FnParseFormGloss{f}{present active participle nominative masculine singular}{ἀγνοέω}{to be ignorant} ὅτι τὸ χρηστὸν\FnParseFormGloss{g}{nominative neuter singular}{χρηστός}{easy} τοῦ θεοῦ εἰς μετάνοιάν\FnParseFormGloss{h}{accusative feminine singular}{μετάνοια}{change of mind} σε ἄγει;\FnParse{i}{present active indicative 3rd singular} 
\MainTextVerseMark{2}{5}κατὰ δὲ τὴν σκληρότητά\FnParseFormGloss{a}{accusative feminine singular}{σκληρότης}{hardness} σου καὶ ἀμετανόητον\FnParseFormGloss{b}{accusative feminine singular}{ἀμετανόητος}{unrepentant} καρδίαν θησαυρίζεις\FnParseFormGloss{c}{present active indicative 2nd singular}{θησαυρίζω}{to store up} σεαυτῷ ὀργὴν ἐν ἡμέρᾳ ὀργῆς καὶ ⸀ἀποκαλύψεως\FnParseFormGloss{d}{genitive feminine singular}{ἀποκάλυψις}{revelation} δικαιοκρισίας\FnParseFormGloss{e}{genitive feminine singular}{δικαιοκρισία}{righteous judgment} τοῦ θεοῦ, 
\MainTextVerseMark{2}{6}ὃς ἀποδώσει\FnParse{a}{future active indicative 3rd singular} ἑκάστῳ κατὰ τὰ ἔργα αὐτοῦ· 
\MainTextVerseMark{2}{7}τοῖς μὲν καθ’ ὑπομονὴν ἔργου ἀγαθοῦ δόξαν καὶ τιμὴν καὶ ἀφθαρσίαν\FnParseFormGloss{a}{accusative feminine singular}{ἀφθαρσία}{imperishableness} ζητοῦσιν\FnParse{b}{present active participle dative masculine plural} ζωὴν αἰώνιον· 
\MainTextVerseMark{2}{8}τοῖς δὲ ἐξ ἐριθείας\FnParseFormGloss{a}{genitive feminine singular}{ἐριθεία}{selfish ambition} καὶ ⸀ἀπειθοῦσι\FnParseFormGloss{b}{present active participle dative masculine plural}{ἀπειθέω}{to disobey} τῇ ἀληθείᾳ πειθομένοις\FnParse{c}{present passive participle dative masculine plural} δὲ τῇ ἀδικίᾳ\FnParseFormGloss{d}{dative feminine singular}{ἀδικία}{wickedness} ⸂ὀργὴ καὶ θυμός⸃,\FnParseFormGloss{e}{nominative masculine singular}{θυμός}{wrath} 
\MainTextVerseMark{2}{9}θλῖψις καὶ στενοχωρία,\FnParseFormGloss{a}{nominative feminine singular}{στενοχωρία}{distress} ἐπὶ πᾶσαν ψυχὴν ἀνθρώπου τοῦ κατεργαζομένου\FnParseFormGloss{b}{present middle participle genitive masculine singular}{κατεργάζομαι}{to produce} τὸ κακόν, Ἰουδαίου τε πρῶτον καὶ Ἕλληνος·\FnParseFormGloss{c}{genitive masculine singular}{Ἕλλην}{Greek} 
\MainTextVerseMark{2}{10}δόξα δὲ καὶ τιμὴ καὶ εἰρήνη παντὶ τῷ ἐργαζομένῳ\FnParse{a}{present middle participle dative masculine singular} τὸ ἀγαθόν, Ἰουδαίῳ τε πρῶτον καὶ Ἕλληνι·\FnParseFormGloss{b}{dative masculine singular}{Ἕλλην}{Greek} 
\MainTextVerseMark{2}{11}οὐ γάρ ἐστιν\FnParse{a}{present active indicative 3rd singular} προσωπολημψία\FnParseFormGloss{b}{nominative feminine singular}{προσωπολημψία}{favoritism} παρὰ τῷ θεῷ. 
\MainTextVerseMark{2}{12}Ὅσοι γὰρ ἀνόμως\FnFormGloss{a}{ἀνόμως}{without law} ἥμαρτον,\FnParse{b}{aorist active indicative 3rd plural} ἀνόμως\FnFormGloss{c}{ἀνόμως}{without law} καὶ ἀπολοῦνται·\FnParse{d}{future middle indicative 3rd plural} καὶ ὅσοι ἐν νόμῳ ἥμαρτον,\FnParse{e}{aorist active indicative 3rd plural} διὰ νόμου κριθήσονται·\FnParse{f}{future passive indicative 3rd plural} 
\MainTextVerseMark{2}{13}οὐ γὰρ οἱ ἀκροαταὶ\FnParseFormGloss{a}{nominative masculine plural}{ἀκροατής}{hearer} ⸀νόμου δίκαιοι παρὰ ⸀τῷ θεῷ, ἀλλ’ οἱ ποιηταὶ\FnParseFormGloss{b}{nominative masculine plural}{ποιητής}{doer} ⸁νόμου δικαιωθήσονται.\FnParse{c}{future passive indicative 3rd plural} 
\MainTextVerseMark{2}{14}ὅταν γὰρ ἔθνη τὰ μὴ νόμον ἔχοντα\FnParse{a}{present active participle nominative neuter plural} φύσει\FnParseFormGloss{b}{dative feminine singular}{φύσις}{nature} τὰ τοῦ νόμου ⸀ποιῶσιν,\FnParse{c}{present active subjunctive 3rd plural} οὗτοι νόμον μὴ ἔχοντες\FnParse{d}{present active participle nominative masculine plural} ἑαυτοῖς εἰσιν\FnParse{e}{present active indicative 3rd plural} νόμος· 
\MainTextVerseMark{2}{15}οἵτινες ἐνδείκνυνται\FnParseFormGloss{a}{present middle indicative 3rd plural}{ἐνδείκνυμι}{to show} τὸ ἔργον τοῦ νόμου γραπτὸν\FnParseFormGloss{b}{accusative neuter singular}{γραπτός}{written} ἐν ταῖς καρδίαις αὐτῶν, συμμαρτυρούσης\FnParseFormGloss{c}{present active participle genitive feminine singular}{συμμαρτυρέω}{to testify with} αὐτῶν τῆς συνειδήσεως καὶ μεταξὺ\FnFormGloss{d}{μεταξύ}{between} ἀλλήλων τῶν λογισμῶν\FnParseFormGloss{e}{genitive masculine plural}{λογισμός}{thought} κατηγορούντων\FnParseFormGloss{f}{present active participle genitive masculine plural}{κατηγορέω}{to accuse} ἢ καὶ ἀπολογουμένων,\FnParseFormGloss{g}{present middle participle genitive masculine plural}{ἀπολογέομαι}{to defend oneself} 
\MainTextVerseMark{2}{16}ἐν ⸂ἡμέρᾳ ὅτε⸃ ⸀κρίνει\FnParse{a}{present active indicative 3rd singular} ὁ θεὸς τὰ κρυπτὰ\FnParseFormGloss{b}{accusative neuter plural}{κρυπτός}{hidden} τῶν ἀνθρώπων κατὰ τὸ εὐαγγέλιόν μου διὰ ⸂Χριστοῦ Ἰησοῦ⸃. 
\MainTextVerseMark{2}{17}⸂Εἰ δὲ⸃ σὺ Ἰουδαῖος ἐπονομάζῃ\FnParseFormGloss{a}{present middle indicative 2nd singular}{ἐπονομάζομαι}{to be called} καὶ ἐπαναπαύῃ\FnParseFormGloss{b}{present middle indicative 2nd singular}{ἐπαναπαύομαι}{to rest on} ⸀νόμῳ καὶ καυχᾶσαι\FnParse{c}{present middle indicative 2nd singular} ἐν θεῷ 
\MainTextVerseMark{2}{18}καὶ γινώσκεις\FnParse{a}{present active indicative 2nd singular} τὸ θέλημα καὶ δοκιμάζεις\FnParseFormGloss{b}{present active indicative 2nd singular}{δοκιμάζω}{to test} τὰ διαφέροντα\FnParseFormGloss{c}{present active participle accusative neuter plural}{διαφέρω}{to carry} κατηχούμενος\FnParseFormGloss{d}{present passive participle nominative masculine singular}{κατηχέω}{to instruct} ἐκ τοῦ νόμου, 
\MainTextVerseMark{2}{19}πέποιθάς\FnParse{a}{perfect active indicative 2nd singular} τε σεαυτὸν ὁδηγὸν\FnParseFormGloss{b}{accusative masculine singular}{ὁδηγός}{guide} εἶναι\FnParse{c}{present active infinitive} τυφλῶν, φῶς τῶν ἐν σκότει, 
\MainTextVerseMark{2}{20}παιδευτὴν\FnParseFormGloss{a}{accusative masculine singular}{παιδευτής}{instructor} ἀφρόνων,\FnParseFormGloss{b}{genitive masculine plural}{ἄφρων}{foolish} διδάσκαλον νηπίων,\FnParseFormGloss{c}{genitive masculine plural}{νήπιος}{child} ἔχοντα\FnParse{d}{present active participle accusative masculine singular} τὴν μόρφωσιν\FnParseFormGloss{e}{accusative feminine singular}{μόρφωσις}{embodiment} τῆς γνώσεως\FnParseFormGloss{f}{genitive feminine singular}{γνῶσις}{knowledge} καὶ τῆς ἀληθείας ἐν τῷ νόμῳ— 
\MainTextVerseMark{2}{21}ὁ οὖν διδάσκων\FnParse{a}{present active participle nominative masculine singular} ἕτερον σεαυτὸν οὐ διδάσκεις;\FnParse{b}{present active indicative 2nd singular} ὁ κηρύσσων\FnParse{c}{present active participle nominative masculine singular} μὴ κλέπτειν\FnParseFormGloss{d}{present active infinitive}{κλέπτω}{steal} κλέπτεις;\FnParseFormGloss{e}{present active indicative 2nd singular}{κλέπτω}{steal} 
\MainTextVerseMark{2}{22}ὁ λέγων\FnParse{a}{present active participle nominative masculine singular} μὴ μοιχεύειν\FnParseFormGloss{b}{present active infinitive}{μοιχεύω}{to commit adultery} μοιχεύεις;\FnParseFormGloss{c}{present active indicative 2nd singular}{μοιχεύω}{to commit adultery} ὁ βδελυσσόμενος\FnParseFormGloss{d}{present middle participle nominative masculine singular}{βδελύσσομαι}{to abhor} τὰ εἴδωλα\FnParseFormGloss{e}{accusative neuter plural}{εἴδωλον}{idol} ἱεροσυλεῖς;\FnParseFormGloss{f}{present active indicative 2nd singular}{ἱεροσυλέω}{to rob temples} 
\MainTextVerseMark{2}{23}ὃς ἐν νόμῳ καυχᾶσαι,\FnParse{a}{present middle indicative 2nd singular} διὰ τῆς παραβάσεως\FnParseFormGloss{b}{genitive feminine singular}{παράβασις}{transgression} τοῦ νόμου τὸν θεὸν ἀτιμάζεις;\FnParseFormGloss{c}{present active indicative 2nd singular}{ἀτιμάζω}{to dishonor} 
\MainTextVerseMark{2}{24}τὸ γὰρ ὄνομα τοῦ θεοῦ δι’ ὑμᾶς βλασφημεῖται\FnParse{a}{present passive indicative 3rd singular} ἐν τοῖς ἔθνεσιν, καθὼς γέγραπται.\FnParse{b}{perfect passive indicative 3rd singular} 
\MainTextVerseMark{2}{25}Περιτομὴ μὲν γὰρ ὠφελεῖ\FnParseFormGloss{a}{present active indicative 3rd singular}{ὠφελέω}{to be of good use} ἐὰν νόμον πράσσῃς·\FnParse{b}{present active subjunctive 2nd singular} ἐὰν δὲ παραβάτης\FnParseFormGloss{c}{nominative masculine singular}{παραβάτης}{lawbreaker} νόμου ᾖς,\FnParse{d}{present active subjunctive 2nd singular} ἡ περιτομή σου ἀκροβυστία\FnParseFormGloss{e}{nominative feminine singular}{ἀκροβυστία}{uncircumcision} γέγονεν.\FnParse{f}{perfect active indicative 3rd singular} 
\MainTextVerseMark{2}{26}ἐὰν οὖν ἡ ἀκροβυστία\FnParseFormGloss{a}{nominative feminine singular}{ἀκροβυστία}{uncircumcision} τὰ δικαιώματα\FnParseFormGloss{b}{accusative neuter plural}{δικαίωμα}{regulation} τοῦ νόμου φυλάσσῃ,\FnParse{c}{present active subjunctive 3rd singular} ⸀οὐχ ἡ ἀκροβυστία\FnParseFormGloss{d}{nominative feminine singular}{ἀκροβυστία}{uncircumcision} αὐτοῦ εἰς περιτομὴν λογισθήσεται;\FnParse{e}{future passive indicative 3rd singular} 
\MainTextVerseMark{2}{27}καὶ κρινεῖ\FnParse{a}{future active indicative 3rd singular} ἡ ἐκ φύσεως\FnParseFormGloss{b}{genitive feminine singular}{φύσις}{nature} ἀκροβυστία\FnParseFormGloss{c}{nominative feminine singular}{ἀκροβυστία}{uncircumcision} τὸν νόμον τελοῦσα\FnParseFormGloss{d}{present active participle nominative feminine singular}{τελέω}{to finish} σὲ τὸν διὰ γράμματος\FnParseFormGloss{e}{genitive neuter singular}{γράμμα}{letter} καὶ περιτομῆς παραβάτην\FnParseFormGloss{f}{accusative masculine singular}{παραβάτης}{lawbreaker} νόμου. 
\MainTextVerseMark{2}{28}οὐ γὰρ ὁ ἐν τῷ φανερῷ\FnParseFormGloss{a}{dative neuter singular}{φανερός}{visible} Ἰουδαῖός ἐστιν,\FnParse{b}{present active indicative 3rd singular} οὐδὲ ἡ ἐν τῷ φανερῷ\FnParseFormGloss{c}{dative neuter singular}{φανερός}{visible} ἐν σαρκὶ περιτομή· 
\MainTextVerseMark{2}{29}ἀλλ’ ὁ ἐν τῷ κρυπτῷ\FnParseFormGloss{a}{dative neuter singular}{κρυπτός}{hidden} Ἰουδαῖος, καὶ περιτομὴ καρδίας ἐν πνεύματι οὐ γράμματι,\FnParseFormGloss{b}{dative neuter singular}{γράμμα}{letter} οὗ ὁ ἔπαινος\FnParseFormGloss{c}{nominative masculine singular}{ἔπαινος}{praise} οὐκ ἐξ ἀνθρώπων ἀλλ’ ἐκ τοῦ θεοῦ. 
\ChapterHeader{3}{Chapter 3}

\MainTextVerseMark{3}{1}Τί οὖν τὸ περισσὸν\FnParseFormGloss{a}{nominative neuter singular}{περισσός}{exceeding} τοῦ Ἰουδαίου, ἢ τίς ἡ ὠφέλεια\FnParseFormGloss{b}{nominative feminine singular}{ὠφέλεια}{value} τῆς περιτομῆς; 
\MainTextVerseMark{3}{2}πολὺ κατὰ πάντα τρόπον.\FnParseFormGloss{a}{accusative masculine singular}{τρόπος}{manner} πρῶτον μὲν ⸀γὰρ ὅτι ἐπιστεύθησαν\FnParse{b}{aorist passive indicative 3rd plural} τὰ λόγια\FnParseFormGloss{c}{accusative neuter plural}{λόγιον}{words} τοῦ θεοῦ. 
\MainTextVerseMark{3}{3}τί γάρ; εἰ ἠπίστησάν\FnParseFormGloss{a}{aorist active indicative 3rd plural}{ἀπιστέω}{to disbelieve} τινες, μὴ ἡ ἀπιστία\FnParseFormGloss{b}{nominative feminine singular}{ἀπιστία}{unbelief} αὐτῶν τὴν πίστιν τοῦ θεοῦ καταργήσει;\FnParseFormGloss{c}{future active indicative 3rd singular}{καταργέω}{to nullify} 
\MainTextVerseMark{3}{4}μὴ γένοιτο·\FnParse{a}{aorist middle optative 3rd singular} γινέσθω\FnParse{b}{present middle imperative 3rd singular} δὲ ὁ θεὸς ἀληθής,\FnParseFormGloss{c}{nominative masculine singular}{ἀληθής}{true} πᾶς δὲ ἄνθρωπος ψεύστης,\FnParseFormGloss{d}{nominative masculine singular}{ψεύστης}{liar} ⸀καθὼς γέγραπται·\FnParse{e}{perfect passive indicative 3rd singular} Ὅπως ἂν δικαιωθῇς\FnParse{f}{aorist passive subjunctive 2nd singular} ἐν τοῖς λόγοις σου καὶ ⸀νικήσεις\FnParseFormGloss{g}{future active indicative 2nd singular}{νικάω}{to overcome} ἐν τῷ κρίνεσθαί\FnParse{h}{present passive infinitive} σε. 
\MainTextVerseMark{3}{5}εἰ δὲ ἡ ἀδικία\FnParseFormGloss{a}{nominative feminine singular}{ἀδικία}{wickedness} ἡμῶν θεοῦ δικαιοσύνην συνίστησιν,\FnParseFormGloss{b}{present active indicative 3rd singular}{συνίστημι}{to commend} τί ἐροῦμεν;\FnParse{c}{future active indicative 1st plural} μὴ ἄδικος\FnParseFormGloss{d}{nominative masculine singular}{ἄδικος}{unjust} ὁ θεὸς ὁ ἐπιφέρων\FnParseFormGloss{e}{present active participle nominative masculine singular}{ἐπιφέρω}{to bring upon} τὴν ὀργήν; κατὰ ἄνθρωπον λέγω.\FnParse{f}{present active indicative 1st singular} 
\MainTextVerseMark{3}{6}μὴ γένοιτο·\FnParse{a}{aorist middle optative 3rd singular} ἐπεὶ\FnFormGloss{b}{ἐπεί}{since} πῶς κρινεῖ\FnParse{c}{future active indicative 3rd singular} ὁ θεὸς τὸν κόσμον; 
\MainTextVerseMark{3}{7}εἰ ⸀δὲ ἡ ἀλήθεια τοῦ θεοῦ ἐν τῷ ἐμῷ ψεύσματι\FnParseFormGloss{a}{dative neuter singular}{ψεῦσμα}{falsehood} ἐπερίσσευσεν\FnParse{b}{aorist active indicative 3rd singular} εἰς τὴν δόξαν αὐτοῦ, τί ἔτι κἀγὼ ὡς ἁμαρτωλὸς κρίνομαι,\FnParse{c}{present passive indicative 1st singular} 
\MainTextVerseMark{3}{8}καὶ μὴ καθὼς βλασφημούμεθα\FnParse{a}{present passive indicative 1st plural} καὶ καθώς φασίν\FnParse{b}{present active indicative 3rd plural} τινες ἡμᾶς λέγειν\FnParse{c}{present active infinitive} ὅτι Ποιήσωμεν\FnParse{d}{aorist active subjunctive 1st plural} τὰ κακὰ ἵνα ἔλθῃ\FnParse{e}{aorist active subjunctive 3rd singular} τὰ ἀγαθά; ὧν τὸ κρίμα\FnParseFormGloss{f}{nominative neuter singular}{κρίμα}{judgment} ἔνδικόν\FnParseFormGloss{g}{nominative neuter singular}{ἔνδικος}{just} ἐστιν.\FnParse{h}{present active indicative 3rd singular} 
\MainTextVerseMark{3}{9}Τί οὖν; προεχόμεθα;\FnParseFormGloss{a}{present middle indicative 1st plural}{προέχομαι}{to be better off} οὐ πάντως,\FnFormGloss{b}{πάντως}{surely} προῃτιασάμεθα\FnParseFormGloss{c}{aorist middle indicative 1st plural}{προαιτιάομαι}{to make a charge beforehand} γὰρ Ἰουδαίους τε καὶ Ἕλληνας\FnParseFormGloss{d}{accusative masculine plural}{Ἕλλην}{Greek} πάντας ὑφ’ ἁμαρτίαν εἶναι,\FnParse{e}{present active infinitive} 
\MainTextVerseMark{3}{10}καθὼς γέγραπται\FnParse{a}{perfect passive indicative 3rd singular} ὅτι Οὐκ ἔστιν\FnParse{b}{present active indicative 3rd singular} δίκαιος οὐδὲ εἷς, 
\MainTextVerseMark{3}{11}οὐκ ἔστιν\FnParse{a}{present active indicative 3rd singular} ⸀ὁ συνίων,\FnParseFormGloss{b}{present active participle nominative masculine singular}{συνίημι}{to understand} οὐκ ἔστιν\FnParse{c}{present active indicative 3rd singular} ⸁ὁ ἐκζητῶν\FnParseFormGloss{d}{present active participle nominative masculine singular}{ἐκζητέω}{to seek out} τὸν θεόν· 
\MainTextVerseMark{3}{12}πάντες ἐξέκλιναν,\FnParseFormGloss{a}{aorist active indicative 3rd plural}{ἐκκλίνω}{to turn away} ἅμα\FnFormGloss{b}{ἅμα}{together} ἠχρεώθησαν·\FnParseFormGloss{c}{aorist passive indicative 3rd plural}{ἀχρειόομαι}{to become worthless} οὐκ ⸀ἔστιν\FnParse{d}{present active indicative 3rd singular} ποιῶν\FnParse{e}{present active participle nominative masculine singular} χρηστότητα,\FnParseFormGloss{f}{accusative feminine singular}{χρηστότης}{kindness} οὐκ ἔστιν\FnParse{g}{present active indicative 3rd singular} ἕως ἑνός. 
\MainTextVerseMark{3}{13}τάφος\FnParseFormGloss{a}{nominative masculine singular}{τάφος}{tomb} ἀνεῳγμένος\FnParse{b}{perfect passive participle nominative masculine singular} ὁ λάρυγξ\FnParseFormGloss{c}{nominative masculine singular}{λάρυγξ}{throat} αὐτῶν, ταῖς γλώσσαις αὐτῶν ἐδολιοῦσαν,\FnParseFormGloss{d}{imperfect active indicative 3rd plural}{δολιόω}{to practice deceit} ἰὸς\FnParseFormGloss{e}{nominative masculine singular}{ἰός}{poison} ἀσπίδων\FnParseFormGloss{f}{genitive feminine plural}{ἀσπίς}{viper} ὑπὸ τὰ χείλη\FnParseFormGloss{g}{accusative neuter plural}{χεῖλος}{lip} αὐτῶν, 
\MainTextVerseMark{3}{14}ὧν τὸ στόμα ἀρᾶς\FnParseFormGloss{a}{genitive feminine singular}{ἀρά}{curse} καὶ πικρίας\FnParseFormGloss{b}{genitive feminine singular}{πικρία}{bitterness} γέμει·\FnParseFormGloss{c}{present active indicative 3rd singular}{γέμω}{to be full} 
\MainTextVerseMark{3}{15}ὀξεῖς\FnParseFormGloss{a}{nominative masculine plural}{ὀξύς}{sharp} οἱ πόδες αὐτῶν ἐκχέαι\FnParseFormGloss{b}{aorist active infinitive}{ἐκχέω}{to pour out} αἷμα, 
\MainTextVerseMark{3}{16}σύντριμμα\FnParseFormGloss{a}{nominative neuter singular}{σύντριμμα}{ruin} καὶ ταλαιπωρία\FnParseFormGloss{b}{nominative feminine singular}{ταλαιπωρία}{misery} ἐν ταῖς ὁδοῖς αὐτῶν, 
\MainTextVerseMark{3}{17}καὶ ὁδὸν εἰρήνης οὐκ ἔγνωσαν.\FnParse{a}{aorist active indicative 3rd plural} 
\MainTextVerseMark{3}{18}οὐκ ἔστιν\FnParse{a}{present active indicative 3rd singular} φόβος θεοῦ ἀπέναντι\FnFormGloss{b}{ἀπέναντι}{opposite} τῶν ὀφθαλμῶν αὐτῶν. 
\MainTextVerseMark{3}{19}Οἴδαμεν\FnParse{a}{perfect active indicative 1st plural} δὲ ὅτι ὅσα ὁ νόμος λέγει\FnParse{b}{present active indicative 3rd singular} τοῖς ἐν τῷ νόμῳ λαλεῖ,\FnParse{c}{present active indicative 3rd singular} ἵνα πᾶν στόμα φραγῇ\FnParseFormGloss{d}{aorist passive subjunctive 3rd singular}{φράσσω}{to shut} καὶ ὑπόδικος\FnParseFormGloss{e}{nominative masculine singular}{ὑπόδικος}{accountable} γένηται\FnParse{f}{aorist middle subjunctive 3rd singular} πᾶς ὁ κόσμος τῷ θεῷ· 
\MainTextVerseMark{3}{20}διότι\FnFormGloss{a}{διότι}{therefore} ἐξ ἔργων νόμου οὐ δικαιωθήσεται\FnParse{b}{future passive indicative 3rd singular} πᾶσα σὰρξ ἐνώπιον αὐτοῦ, διὰ γὰρ νόμου ἐπίγνωσις\FnParseFormGloss{c}{nominative feminine singular}{ἐπίγνωσις}{knowledge} ἁμαρτίας. 
\MainTextVerseMark{3}{21}Νυνὶ\FnFormGloss{a}{νυνί}{now} δὲ χωρὶς νόμου δικαιοσύνη θεοῦ πεφανέρωται,\FnParse{b}{perfect passive indicative 3rd singular} μαρτυρουμένη\FnParse{c}{present passive participle nominative feminine singular} ὑπὸ τοῦ νόμου καὶ τῶν προφητῶν, 
\MainTextVerseMark{3}{22}δικαιοσύνη δὲ θεοῦ διὰ πίστεως Ἰησοῦ Χριστοῦ, εἰς ⸀πάντας τοὺς πιστεύοντας,\FnParse{a}{present active participle accusative masculine plural} οὐ γάρ ἐστιν\FnParse{b}{present active indicative 3rd singular} διαστολή.\FnParseFormGloss{c}{nominative feminine singular}{διαστολή}{difference} 
\MainTextVerseMark{3}{23}πάντες γὰρ ἥμαρτον\FnParse{a}{aorist active indicative 3rd plural} καὶ ὑστεροῦνται\FnParseFormGloss{b}{present passive indicative 3rd plural}{ὑστερέω}{to lack} τῆς δόξης τοῦ θεοῦ, 
\MainTextVerseMark{3}{24}δικαιούμενοι\FnParse{a}{present passive participle nominative masculine plural} δωρεὰν\FnFormGloss{b}{δωρεάν}{freely} τῇ αὐτοῦ χάριτι διὰ τῆς ἀπολυτρώσεως\FnParseFormGloss{c}{genitive feminine singular}{ἀπολύτρωσις}{redemption} τῆς ἐν Χριστῷ Ἰησοῦ· 
\MainTextVerseMark{3}{25}ὃν προέθετο\FnParseFormGloss{a}{aorist middle indicative 3rd singular}{προτίθεμαι}{to plan} ὁ θεὸς ἱλαστήριον\FnParseFormGloss{b}{accusative neuter singular}{ἱλαστήριον}{atoning sacrifice} ⸀διὰ πίστεως ἐν τῷ αὐτοῦ αἵματι εἰς ἔνδειξιν\FnParseFormGloss{c}{accusative feminine singular}{ἔνδειξις}{demonstration} τῆς δικαιοσύνης αὐτοῦ διὰ τὴν πάρεσιν\FnParseFormGloss{d}{accusative feminine singular}{πάρεσις}{leaving unpunished} τῶν προγεγονότων\FnParseFormGloss{e}{perfect active participle genitive neuter plural}{προγίνομαι}{to commit beforehand} ἁμαρτημάτων\FnParseFormGloss{f}{genitive neuter plural}{ἁμάρτημα}{sin} 
\MainTextVerseMark{3}{26}ἐν τῇ ἀνοχῇ\FnParseFormGloss{a}{dative feminine singular}{ἀνοχή}{tolerance} τοῦ θεοῦ, πρὸς ⸀τὴν ἔνδειξιν\FnParseFormGloss{b}{accusative feminine singular}{ἔνδειξις}{demonstration} τῆς δικαιοσύνης αὐτοῦ ἐν τῷ νῦν καιρῷ, εἰς τὸ εἶναι\FnParse{c}{present active infinitive} αὐτὸν δίκαιον καὶ δικαιοῦντα\FnParse{d}{present active participle accusative masculine singular} τὸν ἐκ πίστεως Ἰησοῦ. 
\MainTextVerseMark{3}{27}Ποῦ οὖν ἡ καύχησις;\FnParseFormGloss{a}{nominative feminine singular}{καύχησις}{boasting} ἐξεκλείσθη.\FnParseFormGloss{b}{aorist passive indicative 3rd singular}{ἐκκλείω}{to alienate} διὰ ποίου νόμου; τῶν ἔργων; οὐχί, ἀλλὰ διὰ νόμου πίστεως. 
\MainTextVerseMark{3}{28}λογιζόμεθα\FnParse{a}{present middle indicative 1st plural} ⸀γὰρ ⸂δικαιοῦσθαι\FnParse{b}{present passive infinitive} πίστει⸃ ἄνθρωπον χωρὶς ἔργων νόμου. 
\MainTextVerseMark{3}{29}ἢ Ἰουδαίων ὁ θεὸς μόνον; ⸀οὐχὶ καὶ ἐθνῶν; ναὶ καὶ ἐθνῶν, 
\MainTextVerseMark{3}{30}⸀εἴπερ\FnFormGloss{a}{εἴπερ}{if indeed} εἷς ὁ θεός, ὃς δικαιώσει\FnParse{b}{future active indicative 3rd singular} περιτομὴν ἐκ πίστεως καὶ ἀκροβυστίαν\FnParseFormGloss{c}{accusative feminine singular}{ἀκροβυστία}{uncircumcision} διὰ τῆς πίστεως. 
\MainTextVerseMark{3}{31}νόμον οὖν καταργοῦμεν\FnParseFormGloss{a}{present active indicative 1st plural}{καταργέω}{to nullify} διὰ τῆς πίστεως; μὴ γένοιτο,\FnParse{b}{aorist middle optative 3rd singular} ἀλλὰ νόμον ⸀ἱστάνομεν.\FnParse{c}{present active indicative 1st plural} 
\ChapterHeader{4}{Chapter 4}

\MainTextVerseMark{4}{1}Τί οὖν ἐροῦμεν\FnParse{a}{future active indicative 1st plural} ⸀εὑρηκέναι\FnParse{b}{perfect active infinitive} Ἀβραὰμ τὸν ⸀προπάτορα\FnParseFormGloss{c}{accusative masculine singular}{προπάτωρ}{forefather} ⸀ἡμῶν κατὰ σάρκα; 
\MainTextVerseMark{4}{2}εἰ γὰρ Ἀβραὰμ ἐξ ἔργων ἐδικαιώθη,\FnParse{a}{aorist passive indicative 3rd singular} ἔχει\FnParse{b}{present active indicative 3rd singular} καύχημα·\FnParseFormGloss{c}{accusative neuter singular}{καύχημα}{something to boast about} ἀλλ’ οὐ ⸀πρὸς θεόν, 
\MainTextVerseMark{4}{3}τί γὰρ ἡ γραφὴ λέγει;\FnParse{a}{present active indicative 3rd singular} Ἐπίστευσεν\FnParse{b}{aorist active indicative 3rd singular} δὲ Ἀβραὰμ τῷ θεῷ καὶ ἐλογίσθη\FnParse{c}{aorist passive indicative 3rd singular} αὐτῷ εἰς δικαιοσύνην. 
\MainTextVerseMark{4}{4}τῷ δὲ ἐργαζομένῳ\FnParse{a}{present middle participle dative masculine singular} ὁ μισθὸς\FnParseFormGloss{b}{nominative masculine singular}{μισθός}{wage} οὐ λογίζεται\FnParse{c}{present passive indicative 3rd singular} κατὰ χάριν ἀλλὰ κατὰ ὀφείλημα·\FnParseFormGloss{d}{accusative neuter singular}{ὀφείλημα}{debt} 
\MainTextVerseMark{4}{5}τῷ δὲ μὴ ἐργαζομένῳ,\FnParse{a}{present middle participle dative masculine singular} πιστεύοντι\FnParse{b}{present active participle dative masculine singular} δὲ ἐπὶ τὸν δικαιοῦντα\FnParse{c}{present active participle accusative masculine singular} τὸν ἀσεβῆ,\FnParseFormGloss{d}{accusative masculine singular}{ἀσεβής}{ungodly} λογίζεται\FnParse{e}{present passive indicative 3rd singular} ἡ πίστις αὐτοῦ εἰς δικαιοσύνην, 
\MainTextVerseMark{4}{6}καθάπερ\FnFormGloss{a}{καθάπερ}{as} καὶ Δαυὶδ λέγει\FnParse{b}{present active indicative 3rd singular} τὸν μακαρισμὸν\FnParseFormGloss{c}{accusative masculine singular}{μακαρισμός}{blessedness} τοῦ ἀνθρώπου ᾧ ὁ θεὸς λογίζεται\FnParse{d}{present middle indicative 3rd singular} δικαιοσύνην χωρὶς ἔργων· 
\MainTextVerseMark{4}{7}Μακάριοι ὧν ἀφέθησαν\FnParse{a}{aorist passive indicative 3rd plural} αἱ ἀνομίαι\FnParseFormGloss{b}{nominative feminine plural}{ἀνομία}{wickedness} καὶ ὧν ἐπεκαλύφθησαν\FnParseFormGloss{c}{aorist passive indicative 3rd plural}{ἐπικαλύπτω}{to be covered} αἱ ἁμαρτίαι, 
\MainTextVerseMark{4}{8}μακάριος ἀνὴρ ⸀οὗ οὐ μὴ λογίσηται\FnParse{a}{aorist middle subjunctive 3rd singular} κύριος ἁμαρτίαν. 
\MainTextVerseMark{4}{9}Ὁ μακαρισμὸς\FnParseFormGloss{a}{nominative masculine singular}{μακαρισμός}{blessedness} οὖν οὗτος ἐπὶ τὴν περιτομὴν ἢ καὶ ἐπὶ τὴν ἀκροβυστίαν;\FnParseFormGloss{b}{accusative feminine singular}{ἀκροβυστία}{uncircumcision} λέγομεν\FnParse{c}{present active indicative 1st plural} ⸀γάρ· Ἐλογίσθη\FnParse{d}{aorist passive indicative 3rd singular} τῷ Ἀβραὰμ ἡ πίστις εἰς δικαιοσύνην. 
\MainTextVerseMark{4}{10}πῶς οὖν ἐλογίσθη;\FnParse{a}{aorist passive indicative 3rd singular} ἐν περιτομῇ ὄντι\FnParse{b}{present active participle dative masculine singular} ἢ ἐν ἀκροβυστίᾳ;\FnParseFormGloss{c}{dative feminine singular}{ἀκροβυστία}{uncircumcision} οὐκ ἐν περιτομῇ ἀλλ’ ἐν ἀκροβυστίᾳ·\FnParseFormGloss{d}{dative feminine singular}{ἀκροβυστία}{uncircumcision} 
\MainTextVerseMark{4}{11}καὶ σημεῖον ἔλαβεν\FnParse{a}{aorist active indicative 3rd singular} περιτομῆς, σφραγῖδα\FnParseFormGloss{b}{accusative feminine singular}{σφραγίς}{seal} τῆς δικαιοσύνης τῆς πίστεως τῆς ἐν τῇ ἀκροβυστίᾳ,\FnParseFormGloss{c}{dative feminine singular}{ἀκροβυστία}{uncircumcision} εἰς τὸ εἶναι\FnParse{d}{present active infinitive} αὐτὸν πατέρα πάντων τῶν πιστευόντων\FnParse{e}{present active participle genitive masculine plural} δι’ ἀκροβυστίας,\FnParseFormGloss{f}{genitive feminine singular}{ἀκροβυστία}{uncircumcision} εἰς τὸ ⸀λογισθῆναι\FnParse{g}{aorist passive infinitive} αὐτοῖς τὴν δικαιοσύνην, 
\MainTextVerseMark{4}{12}καὶ πατέρα περιτομῆς τοῖς οὐκ ἐκ περιτομῆς μόνον ἀλλὰ καὶ τοῖς στοιχοῦσιν\FnParseFormGloss{a}{present active participle dative masculine plural}{στοιχέω}{to follow} τοῖς ἴχνεσιν\FnParseFormGloss{b}{dative neuter plural}{ἴχνος}{step} τῆς ⸂ἐν ἀκροβυστίᾳ\FnParseFormGloss{c}{dative feminine singular}{ἀκροβυστία}{uncircumcision} πίστεως⸃ τοῦ πατρὸς ἡμῶν Ἀβραάμ. 
\MainTextVerseMark{4}{13}Οὐ γὰρ διὰ νόμου ἡ ἐπαγγελία τῷ Ἀβραὰμ ἢ τῷ σπέρματι αὐτοῦ, τὸ κληρονόμον\FnParseFormGloss{a}{accusative masculine singular}{κληρονόμος}{heir} αὐτὸν ⸀εἶναι\FnParse{b}{present active infinitive} κόσμου, ἀλλὰ διὰ δικαιοσύνης πίστεως· 
\MainTextVerseMark{4}{14}εἰ γὰρ οἱ ἐκ νόμου κληρονόμοι,\FnParseFormGloss{a}{nominative masculine plural}{κληρονόμος}{heir} κεκένωται\FnParseFormGloss{b}{perfect passive indicative 3rd singular}{κενόω}{to empty} ἡ πίστις καὶ κατήργηται\FnParseFormGloss{c}{perfect passive indicative 3rd singular}{καταργέω}{to nullify} ἡ ἐπαγγελία· 
\MainTextVerseMark{4}{15}ὁ γὰρ νόμος ὀργὴν κατεργάζεται,\FnParseFormGloss{a}{present middle indicative 3rd singular}{κατεργάζομαι}{to produce} οὗ\FnFormGloss{b}{οὗ}{where} ⸀δὲ οὐκ ἔστιν\FnParse{c}{present active indicative 3rd singular} νόμος, οὐδὲ παράβασις.\FnParseFormGloss{d}{nominative feminine singular}{παράβασις}{transgression} 
\MainTextVerseMark{4}{16}Διὰ τοῦτο ἐκ πίστεως, ἵνα κατὰ χάριν, εἰς τὸ εἶναι\FnParse{a}{present active infinitive} βεβαίαν\FnParseFormGloss{b}{accusative feminine singular}{βέβαιος}{firm} τὴν ἐπαγγελίαν παντὶ τῷ σπέρματι, οὐ τῷ ἐκ τοῦ νόμου μόνον ἀλλὰ καὶ τῷ ἐκ πίστεως Ἀβραάμ (ὅς ἐστιν\FnParse{c}{present active indicative 3rd singular} πατὴρ πάντων ἡμῶν, 
\MainTextVerseMark{4}{17}καθὼς γέγραπται\FnParse{a}{perfect passive indicative 3rd singular} ὅτι Πατέρα πολλῶν ἐθνῶν τέθεικά\FnParse{b}{perfect active indicative 1st singular} σε), κατέναντι\FnFormGloss{c}{κατέναντι}{ahead} οὗ ἐπίστευσεν\FnParse{d}{aorist active indicative 3rd singular} θεοῦ τοῦ ζῳοποιοῦντος\FnParseFormGloss{e}{present active participle genitive masculine singular}{ζῳοποιέω}{to make alive} τοὺς νεκροὺς καὶ καλοῦντος\FnParse{f}{present active participle genitive masculine singular} τὰ μὴ ὄντα\FnParse{g}{present active participle accusative neuter plural} ὡς ὄντα·\FnParse{h}{present active participle accusative neuter plural} 
\MainTextVerseMark{4}{18}ὃς παρ’ ἐλπίδα ἐπ’ ἐλπίδι ἐπίστευσεν\FnParse{a}{aorist active indicative 3rd singular} εἰς τὸ γενέσθαι\FnParse{b}{aorist middle infinitive} αὐτὸν πατέρα πολλῶν ἐθνῶν κατὰ τὸ εἰρημένον·\FnParse{c}{perfect passive participle accusative neuter singular} Οὕτως ἔσται\FnParse{d}{future middle indicative 3rd singular} τὸ σπέρμα σου· 
\MainTextVerseMark{4}{19}καὶ μὴ ἀσθενήσας\FnParse{a}{aorist active participle nominative masculine singular} τῇ ⸀πίστει κατενόησεν\FnParseFormGloss{b}{aorist active indicative 3rd singular}{κατανοέω}{to pay attention} τὸ ἑαυτοῦ ⸀σῶμα νενεκρωμένον,\FnParseFormGloss{c}{perfect passive participle accusative neuter singular}{νεκρόω}{to put to death} ἑκατονταετής\FnParseFormGloss{d}{nominative masculine singular}{ἑκατονταετής}{a hundred years old} που\FnFormGloss{e}{πού}{somewhere} ὑπάρχων,\FnParse{f}{present active participle nominative masculine singular} καὶ τὴν νέκρωσιν\FnParseFormGloss{g}{accusative feminine singular}{νέκρωσις}{death} τῆς μήτρας\FnParseFormGloss{h}{genitive feminine singular}{μήτρα}{womb} Σάρρας,\FnParseFormGloss{i}{genitive feminine singular}{Σάρρα}{Sarah} 
\MainTextVerseMark{4}{20}εἰς δὲ τὴν ἐπαγγελίαν τοῦ θεοῦ οὐ διεκρίθη\FnParseFormGloss{a}{aorist passive indicative 3rd singular}{διακρίνω}{to make a distinction} τῇ ἀπιστίᾳ\FnParseFormGloss{b}{dative feminine singular}{ἀπιστία}{unbelief} ἀλλὰ ἐνεδυναμώθη\FnParseFormGloss{c}{aorist passive indicative 3rd singular}{ἐνδυναμόω}{to give strength} τῇ πίστει, δοὺς\FnParse{d}{aorist active participle nominative masculine singular} δόξαν τῷ θεῷ 
\MainTextVerseMark{4}{21}καὶ πληροφορηθεὶς\FnParseFormGloss{a}{aorist passive participle nominative masculine singular}{πληροφορέω}{to fulfill} ὅτι ὃ ἐπήγγελται\FnParseFormGloss{b}{perfect middle indicative 3rd singular}{ἐπαγγέλλομαι}{to promise} δυνατός ἐστιν\FnParse{c}{present active indicative 3rd singular} καὶ ποιῆσαι.\FnParse{d}{aorist active infinitive} 
\MainTextVerseMark{4}{22}⸀διὸ ἐλογίσθη\FnParse{a}{aorist passive indicative 3rd singular} αὐτῷ εἰς δικαιοσύνην. 
\MainTextVerseMark{4}{23}Οὐκ ἐγράφη\FnParse{a}{aorist passive indicative 3rd singular} δὲ δι’ αὐτὸν μόνον ὅτι ἐλογίσθη\FnParse{b}{aorist passive indicative 3rd singular} αὐτῷ, 
\MainTextVerseMark{4}{24}ἀλλὰ καὶ δι’ ἡμᾶς οἷς μέλλει\FnParse{a}{present active indicative 3rd singular} λογίζεσθαι,\FnParse{b}{present passive infinitive} τοῖς πιστεύουσιν\FnParse{c}{present active participle dative masculine plural} ἐπὶ τὸν ἐγείραντα\FnParse{d}{aorist active participle accusative masculine singular} Ἰησοῦν τὸν κύριον ἡμῶν ἐκ νεκρῶν, 
\MainTextVerseMark{4}{25}ὃς παρεδόθη\FnParse{a}{aorist passive indicative 3rd singular} διὰ τὰ παραπτώματα\FnParseFormGloss{b}{accusative neuter plural}{παράπτωμα}{trespass} ἡμῶν καὶ ἠγέρθη\FnParse{c}{aorist passive indicative 3rd singular} διὰ τὴν δικαίωσιν\FnParseFormGloss{d}{accusative feminine singular}{δικαίωσις}{justification} ἡμῶν. 
\ChapterHeader{5}{Chapter 5}

\MainTextVerseMark{5}{1}Δικαιωθέντες\FnParse{a}{aorist passive participle nominative masculine plural} οὖν ἐκ πίστεως εἰρήνην ⸀ἔχομεν\FnParse{b}{present active indicative 1st plural} πρὸς τὸν θεὸν διὰ τοῦ κυρίου ἡμῶν Ἰησοῦ Χριστοῦ, 
\MainTextVerseMark{5}{2}δι’ οὗ καὶ τὴν προσαγωγὴν\FnParseFormGloss{a}{accusative feminine singular}{προσαγωγή}{access} ἐσχήκαμεν\FnParse{b}{perfect active indicative 1st plural} τῇ πίστει εἰς τὴν χάριν ταύτην ἐν ᾗ ἑστήκαμεν,\FnParse{c}{perfect active indicative 1st plural} καὶ καυχώμεθα\FnParse{d}{present middle indicative 1st plural} ἐπ’ ἐλπίδι τῆς δόξης τοῦ θεοῦ· 
\MainTextVerseMark{5}{3}οὐ μόνον δέ, ἀλλὰ καὶ ⸀καυχώμεθα\FnParse{a}{present middle indicative 1st plural} ἐν ταῖς θλίψεσιν, εἰδότες\FnParse{b}{perfect active participle nominative masculine plural} ὅτι ἡ θλῖψις ὑπομονὴν κατεργάζεται,\FnParseFormGloss{c}{present middle indicative 3rd singular}{κατεργάζομαι}{to produce} 
\MainTextVerseMark{5}{4}ἡ δὲ ὑπομονὴ δοκιμήν,\FnParseFormGloss{a}{accusative feminine singular}{δοκιμή}{character} ἡ δὲ δοκιμὴ\FnParseFormGloss{b}{nominative feminine singular}{δοκιμή}{character} ἐλπίδα. 
\MainTextVerseMark{5}{5}ἡ δὲ ἐλπὶς οὐ καταισχύνει·\FnParseFormGloss{a}{present active indicative 3rd singular}{καταισχύνω}{to dishonor} ὅτι ἡ ἀγάπη τοῦ θεοῦ ἐκκέχυται\FnParseFormGloss{b}{perfect passive indicative 3rd singular}{ἐκχέω}{to pour out} ἐν ταῖς καρδίαις ἡμῶν διὰ πνεύματος ἁγίου τοῦ δοθέντος\FnParse{c}{aorist passive participle genitive neuter singular} ἡμῖν. 
\MainTextVerseMark{5}{6}⸂Ἔτι γὰρ⸃ Χριστὸς ὄντων\FnParse{a}{present active participle genitive masculine plural} ἡμῶν ἀσθενῶν\FnParseFormGloss{b}{genitive masculine plural}{ἀσθενής}{weak} ⸀ἔτι κατὰ καιρὸν ὑπὲρ ἀσεβῶν\FnParseFormGloss{c}{genitive masculine plural}{ἀσεβής}{ungodly} ἀπέθανεν.\FnParse{d}{aorist active indicative 3rd singular} 
\MainTextVerseMark{5}{7}μόλις\FnFormGloss{a}{μόλις}{with difficulty} γὰρ ὑπὲρ δικαίου τις ἀποθανεῖται·\FnParse{b}{future middle indicative 3rd singular} ὑπὲρ γὰρ τοῦ ἀγαθοῦ τάχα\FnFormGloss{c}{τάχα}{perhaps} τις καὶ τολμᾷ\FnParseFormGloss{d}{present active indicative 3rd singular}{τολμάω}{to dare} ἀποθανεῖν·\FnParse{e}{aorist active infinitive} 
\MainTextVerseMark{5}{8}συνίστησιν\FnParseFormGloss{a}{present active indicative 3rd singular}{συνίστημι}{to commend} δὲ τὴν ἑαυτοῦ ἀγάπην εἰς ἡμᾶς ὁ θεὸς ὅτι ἔτι ἁμαρτωλῶν ὄντων\FnParse{b}{present active participle genitive masculine plural} ἡμῶν Χριστὸς ὑπὲρ ἡμῶν ἀπέθανεν.\FnParse{c}{aorist active indicative 3rd singular} 
\MainTextVerseMark{5}{9}πολλῷ οὖν μᾶλλον δικαιωθέντες\FnParse{a}{aorist passive participle nominative masculine plural} νῦν ἐν τῷ αἵματι αὐτοῦ σωθησόμεθα\FnParse{b}{future passive indicative 1st plural} δι’ αὐτοῦ ἀπὸ τῆς ὀργῆς. 
\MainTextVerseMark{5}{10}εἰ γὰρ ἐχθροὶ ὄντες\FnParse{a}{present active participle nominative masculine plural} κατηλλάγημεν\FnParseFormGloss{b}{aorist passive indicative 1st plural}{καταλλάσσω}{to reconcile} τῷ θεῷ διὰ τοῦ θανάτου τοῦ υἱοῦ αὐτοῦ, πολλῷ μᾶλλον καταλλαγέντες\FnParseFormGloss{c}{aorist passive participle nominative masculine plural}{καταλλάσσω}{to reconcile} σωθησόμεθα\FnParse{d}{future passive indicative 1st plural} ἐν τῇ ζωῇ αὐτοῦ· 
\MainTextVerseMark{5}{11}οὐ μόνον δέ, ἀλλὰ καὶ καυχώμενοι\FnParse{a}{present middle participle nominative masculine plural} ἐν τῷ θεῷ διὰ τοῦ κυρίου ἡμῶν Ἰησοῦ Χριστοῦ, δι’ οὗ νῦν τὴν καταλλαγὴν\FnParseFormGloss{b}{accusative feminine singular}{καταλλαγή}{reconciliation} ἐλάβομεν.\FnParse{c}{aorist active indicative 1st plural} 
\MainTextVerseMark{5}{12}Διὰ τοῦτο ὥσπερ δι’ ἑνὸς ἀνθρώπου ἡ ἁμαρτία εἰς τὸν κόσμον εἰσῆλθεν\FnParse{a}{aorist active indicative 3rd singular} καὶ διὰ τῆς ἁμαρτίας ὁ θάνατος, καὶ οὕτως εἰς πάντας ἀνθρώπους ὁ θάνατος διῆλθεν\FnParse{b}{aorist active indicative 3rd singular} ἐφ’ ᾧ πάντες ἥμαρτον—\FnParse{c}{aorist active indicative 3rd plural} 
\MainTextVerseMark{5}{13}ἄχρι γὰρ νόμου ἁμαρτία ἦν\FnParse{a}{imperfect active indicative 3rd singular} ἐν κόσμῳ, ἁμαρτία δὲ οὐκ ἐλλογεῖται\FnParseFormGloss{b}{present passive indicative 3rd singular}{ἐλλογέω}{to be charged to one's account} μὴ ὄντος\FnParse{c}{present active participle genitive masculine singular} νόμου, 
\MainTextVerseMark{5}{14}ἀλλὰ ἐβασίλευσεν\FnParseFormGloss{a}{aorist active indicative 3rd singular}{βασιλεύω}{to reign as a king} ὁ θάνατος ἀπὸ Ἀδὰμ\FnParseFormGloss{b}{genitive masculine singular}{Ἀδάμ}{Adam} μέχρι\FnFormGloss{c}{μέχρι(ς)}{until} Μωϋσέως καὶ ἐπὶ τοὺς μὴ ἁμαρτήσαντας\FnParse{d}{aorist active participle accusative masculine plural} ἐπὶ τῷ ὁμοιώματι\FnParseFormGloss{e}{dative neuter singular}{ὁμοίωμα}{likeness} τῆς παραβάσεως\FnParseFormGloss{f}{genitive feminine singular}{παράβασις}{transgression} Ἀδάμ,\FnParseFormGloss{g}{genitive masculine singular}{Ἀδάμ}{Adam} ὅς ἐστιν\FnParse{h}{present active indicative 3rd singular} τύπος\FnParseFormGloss{i}{nominative masculine singular}{τύπος}{pattern} τοῦ μέλλοντος.\FnParse{j}{present active participle genitive masculine singular} 
\MainTextVerseMark{5}{15}Ἀλλ’ οὐχ ὡς τὸ παράπτωμα,\FnParseFormGloss{a}{nominative neuter singular}{παράπτωμα}{trespass} οὕτως καὶ τὸ χάρισμα·\FnParseFormGloss{b}{nominative neuter singular}{χάρισμα}{gracious gift} εἰ γὰρ τῷ τοῦ ἑνὸς παραπτώματι\FnParseFormGloss{c}{dative neuter singular}{παράπτωμα}{trespass} οἱ πολλοὶ ἀπέθανον,\FnParse{d}{aorist active indicative 3rd plural} πολλῷ μᾶλλον ἡ χάρις τοῦ θεοῦ καὶ ἡ δωρεὰ\FnParseFormGloss{e}{nominative feminine singular}{δωρεά}{gift} ἐν χάριτι τῇ τοῦ ἑνὸς ἀνθρώπου Ἰησοῦ Χριστοῦ εἰς τοὺς πολλοὺς ἐπερίσσευσεν.\FnParse{f}{aorist active indicative 3rd singular} 
\MainTextVerseMark{5}{16}καὶ οὐχ ὡς δι’ ἑνὸς ἁμαρτήσαντος\FnParse{a}{aorist active participle genitive masculine singular} τὸ δώρημα·\FnParseFormGloss{b}{nominative neuter singular}{δώρημα}{gift} τὸ μὲν γὰρ κρίμα\FnParseFormGloss{c}{nominative neuter singular}{κρίμα}{judgment} ἐξ ἑνὸς εἰς κατάκριμα,\FnParseFormGloss{d}{accusative neuter singular}{κατάκριμα}{condemnation} τὸ δὲ χάρισμα\FnParseFormGloss{e}{nominative neuter singular}{χάρισμα}{gracious gift} ἐκ πολλῶν παραπτωμάτων\FnParseFormGloss{f}{genitive neuter plural}{παράπτωμα}{trespass} εἰς δικαίωμα.\FnParseFormGloss{g}{accusative neuter singular}{δικαίωμα}{regulation} 
\MainTextVerseMark{5}{17}εἰ γὰρ τῷ τοῦ ἑνὸς παραπτώματι\FnParseFormGloss{a}{dative neuter singular}{παράπτωμα}{trespass} ὁ θάνατος ἐβασίλευσεν\FnParseFormGloss{b}{aorist active indicative 3rd singular}{βασιλεύω}{to reign as a king} διὰ τοῦ ἑνός, πολλῷ μᾶλλον οἱ τὴν περισσείαν\FnParseFormGloss{c}{accusative feminine singular}{περισσεία}{abundance} τῆς χάριτος καὶ τῆς δωρεᾶς\FnParseFormGloss{d}{genitive feminine singular}{δωρεά}{gift} τῆς δικαιοσύνης λαμβάνοντες\FnParse{e}{present active participle nominative masculine plural} ἐν ζωῇ βασιλεύσουσιν\FnParseFormGloss{f}{future active indicative 3rd plural}{βασιλεύω}{to reign as a king} διὰ τοῦ ἑνὸς Ἰησοῦ Χριστοῦ. 
\MainTextVerseMark{5}{18}Ἄρα οὖν ὡς δι’ ἑνὸς παραπτώματος\FnParseFormGloss{a}{genitive neuter singular}{παράπτωμα}{trespass} εἰς πάντας ἀνθρώπους εἰς κατάκριμα,\FnParseFormGloss{b}{accusative neuter singular}{κατάκριμα}{condemnation} οὕτως καὶ δι’ ἑνὸς δικαιώματος\FnParseFormGloss{c}{genitive neuter singular}{δικαίωμα}{regulation} εἰς πάντας ἀνθρώπους εἰς δικαίωσιν\FnParseFormGloss{d}{accusative feminine singular}{δικαίωσις}{justification} ζωῆς· 
\MainTextVerseMark{5}{19}ὥσπερ γὰρ διὰ τῆς παρακοῆς\FnParseFormGloss{a}{genitive feminine singular}{παρακοή}{disobedience} τοῦ ἑνὸς ἀνθρώπου ἁμαρτωλοὶ κατεστάθησαν\FnParseFormGloss{b}{aorist passive indicative 3rd plural}{καθίστημι}{to put in charge} οἱ πολλοί, οὕτως καὶ διὰ τῆς ὑπακοῆς\FnParseFormGloss{c}{genitive feminine singular}{ὑπακοή}{obedience} τοῦ ἑνὸς δίκαιοι κατασταθήσονται\FnParseFormGloss{d}{future passive indicative 3rd plural}{καθίστημι}{to put in charge} οἱ πολλοί. 
\MainTextVerseMark{5}{20}νόμος δὲ παρεισῆλθεν\FnParseFormGloss{a}{aorist active indicative 3rd singular}{παρεισέρχομαι}{to come in} ἵνα πλεονάσῃ\FnParseFormGloss{b}{aorist active subjunctive 3rd singular}{πλεονάζω}{to make increase} τὸ παράπτωμα·\FnParseFormGloss{c}{nominative neuter singular}{παράπτωμα}{trespass} οὗ\FnFormGloss{d}{οὗ}{where} δὲ ἐπλεόνασεν\FnParseFormGloss{e}{aorist active indicative 3rd singular}{πλεονάζω}{to make increase} ἡ ἁμαρτία, ὑπερεπερίσσευσεν\FnParseFormGloss{f}{aorist active indicative 3rd singular}{ὑπερπερισσεύω}{to increase all the more} ἡ χάρις, 
\MainTextVerseMark{5}{21}ἵνα ὥσπερ ἐβασίλευσεν\FnParseFormGloss{a}{aorist active indicative 3rd singular}{βασιλεύω}{to reign as a king} ἡ ἁμαρτία ἐν τῷ θανάτῳ, οὕτως καὶ ἡ χάρις βασιλεύσῃ\FnParseFormGloss{b}{aorist active subjunctive 3rd singular}{βασιλεύω}{to reign as a king} διὰ δικαιοσύνης εἰς ζωὴν αἰώνιον διὰ Ἰησοῦ Χριστοῦ τοῦ κυρίου ἡμῶν. 
\ChapterHeader{6}{Chapter 6}

\MainTextVerseMark{6}{1}Τί οὖν ἐροῦμεν;\FnParse{a}{future active indicative 1st plural} ⸀ἐπιμένωμεν\FnParseFormGloss{b}{present active subjunctive 1st plural}{ἐπιμένω}{to stay} τῇ ἁμαρτίᾳ, ἵνα ἡ χάρις πλεονάσῃ;\FnParseFormGloss{c}{aorist active subjunctive 3rd singular}{πλεονάζω}{to make increase} 
\MainTextVerseMark{6}{2}μὴ γένοιτο·\FnParse{a}{aorist middle optative 3rd singular} οἵτινες ἀπεθάνομεν\FnParse{b}{aorist active indicative 1st plural} τῇ ἁμαρτίᾳ, πῶς ἔτι ζήσομεν\FnParse{c}{future active indicative 1st plural} ἐν αὐτῇ; 
\MainTextVerseMark{6}{3}ἢ ἀγνοεῖτε\FnParseFormGloss{a}{present active indicative 2nd plural}{ἀγνοέω}{to be ignorant} ὅτι ὅσοι ἐβαπτίσθημεν\FnParse{b}{aorist passive indicative 1st plural} εἰς Χριστὸν Ἰησοῦν εἰς τὸν θάνατον αὐτοῦ ἐβαπτίσθημεν;\FnParse{c}{aorist passive indicative 1st plural} 
\MainTextVerseMark{6}{4}συνετάφημεν\FnParseFormGloss{a}{aorist passive indicative 1st plural}{συνθάπτομαι}{to be buried with} οὖν αὐτῷ διὰ τοῦ βαπτίσματος\FnParseFormGloss{b}{genitive neuter singular}{βάπτισμα}{baptism} εἰς τὸν θάνατον, ἵνα ὥσπερ ἠγέρθη\FnParse{c}{aorist passive indicative 3rd singular} Χριστὸς ἐκ νεκρῶν διὰ τῆς δόξης τοῦ πατρός, οὕτως καὶ ἡμεῖς ἐν καινότητι\FnParseFormGloss{d}{dative feminine singular}{καινότης}{newness} ζωῆς περιπατήσωμεν.\FnParse{e}{aorist active subjunctive 1st plural} 
\MainTextVerseMark{6}{5}Εἰ γὰρ σύμφυτοι\FnParseFormGloss{a}{nominative masculine plural}{σύμφυτος}{united} γεγόναμεν\FnParse{b}{perfect active indicative 1st plural} τῷ ὁμοιώματι\FnParseFormGloss{c}{dative neuter singular}{ὁμοίωμα}{likeness} τοῦ θανάτου αὐτοῦ, ἀλλὰ καὶ τῆς ἀναστάσεως ἐσόμεθα·\FnParse{d}{future middle indicative 1st plural} 
\MainTextVerseMark{6}{6}τοῦτο γινώσκοντες\FnParse{a}{present active participle nominative masculine plural} ὅτι ὁ παλαιὸς\FnParseFormGloss{b}{nominative masculine singular}{παλαιός}{old} ἡμῶν ἄνθρωπος συνεσταυρώθη,\FnParseFormGloss{c}{aorist passive indicative 3rd singular}{συσταυρόομαι}{to be crucified with} ἵνα καταργηθῇ\FnParseFormGloss{d}{aorist passive subjunctive 3rd singular}{καταργέω}{to nullify} τὸ σῶμα τῆς ἁμαρτίας, τοῦ μηκέτι\FnFormGloss{e}{μηκέτι}{no longer} δουλεύειν\FnParseFormGloss{f}{present active infinitive}{δουλεύω}{to serve} ἡμᾶς τῇ ἁμαρτίᾳ, 
\MainTextVerseMark{6}{7}ὁ γὰρ ἀποθανὼν\FnParse{a}{aorist active participle nominative masculine singular} δεδικαίωται\FnParse{b}{perfect passive indicative 3rd singular} ἀπὸ τῆς ἁμαρτίας. 
\MainTextVerseMark{6}{8}εἰ δὲ ἀπεθάνομεν\FnParse{a}{aorist active indicative 1st plural} σὺν Χριστῷ, πιστεύομεν\FnParse{b}{present active indicative 1st plural} ὅτι καὶ συζήσομεν\FnParseFormGloss{c}{future active indicative 1st plural}{συζάω}{to live with} αὐτῷ· 
\MainTextVerseMark{6}{9}εἰδότες\FnParse{a}{perfect active participle nominative masculine plural} ὅτι Χριστὸς ἐγερθεὶς\FnParse{b}{aorist passive participle nominative masculine singular} ἐκ νεκρῶν οὐκέτι ἀποθνῄσκει,\FnParse{c}{present active indicative 3rd singular} θάνατος αὐτοῦ οὐκέτι κυριεύει·\FnParseFormGloss{d}{present active indicative 3rd singular}{κυριεύω}{to lord over} 
\MainTextVerseMark{6}{10}ὃ γὰρ ἀπέθανεν,\FnParse{a}{aorist active indicative 3rd singular} τῇ ἁμαρτίᾳ ἀπέθανεν\FnParse{b}{aorist active indicative 3rd singular} ἐφάπαξ·\FnFormGloss{c}{ἐφάπαξ}{once for all} ὃ δὲ ζῇ,\FnParse{d}{present active indicative 3rd singular} ζῇ\FnParse{e}{present active indicative 3rd singular} τῷ θεῷ. 
\MainTextVerseMark{6}{11}οὕτως καὶ ὑμεῖς λογίζεσθε\FnParse{a}{present middle imperative 2nd plural} ἑαυτοὺς ⸂εἶναι\FnParse{b}{present active infinitive} νεκροὺς μὲν⸃ τῇ ἁμαρτίᾳ ζῶντας\FnParse{c}{present active participle accusative masculine plural} δὲ τῷ θεῷ ἐν Χριστῷ ⸀Ἰησοῦ. 
\MainTextVerseMark{6}{12}Μὴ οὖν βασιλευέτω\FnParseFormGloss{a}{present active imperative 3rd singular}{βασιλεύω}{to reign as a king} ἡ ἁμαρτία ἐν τῷ θνητῷ\FnParseFormGloss{b}{dative neuter singular}{θνητός}{mortal} ὑμῶν σώματι εἰς τὸ ⸀ὑπακούειν\FnParseFormGloss{c}{present active infinitive}{ὑπακούω}{to obey} ταῖς ἐπιθυμίαις αὐτοῦ, 
\MainTextVerseMark{6}{13}μηδὲ παριστάνετε\FnParse{a}{present active imperative 2nd plural} τὰ μέλη ὑμῶν ὅπλα\FnParseFormGloss{b}{accusative neuter plural}{ὅπλον}{instrument} ἀδικίας\FnParseFormGloss{c}{genitive feminine singular}{ἀδικία}{wickedness} τῇ ἁμαρτίᾳ, ἀλλὰ παραστήσατε\FnParse{d}{aorist active imperative 2nd plural} ἑαυτοὺς τῷ θεῷ ⸀ὡσεὶ\FnFormGloss{e}{ὡσεί}{like} ἐκ νεκρῶν ζῶντας\FnParse{f}{present active participle accusative masculine plural} καὶ τὰ μέλη ὑμῶν ὅπλα\FnParseFormGloss{g}{accusative neuter plural}{ὅπλον}{instrument} δικαιοσύνης τῷ θεῷ. 
\MainTextVerseMark{6}{14}ἁμαρτία γὰρ ὑμῶν οὐ κυριεύσει,\FnParseFormGloss{a}{future active indicative 3rd singular}{κυριεύω}{to lord over} οὐ γάρ ἐστε\FnParse{b}{present active indicative 2nd plural} ὑπὸ νόμον ἀλλὰ ὑπὸ χάριν. 
\MainTextVerseMark{6}{15}Τί οὖν; ⸀ἁμαρτήσωμεν\FnParse{a}{aorist active subjunctive 1st plural} ὅτι οὐκ ἐσμὲν\FnParse{b}{present active indicative 1st plural} ὑπὸ νόμον ἀλλὰ ὑπὸ χάριν; μὴ γένοιτο·\FnParse{c}{aorist middle optative 3rd singular} 
\MainTextVerseMark{6}{16}οὐκ οἴδατε\FnParse{a}{perfect active indicative 2nd plural} ὅτι ᾧ παριστάνετε\FnParse{b}{present active indicative 2nd plural} ἑαυτοὺς δούλους εἰς ὑπακοήν,\FnParseFormGloss{c}{accusative feminine singular}{ὑπακοή}{obedience} δοῦλοί ἐστε\FnParse{d}{present active indicative 2nd plural} ᾧ ὑπακούετε,\FnParseFormGloss{e}{present active indicative 2nd plural}{ὑπακούω}{to obey} ἤτοι\FnFormGloss{f}{ἤτοι}{whether...or} ἁμαρτίας εἰς θάνατον ἢ ὑπακοῆς\FnParseFormGloss{g}{genitive feminine singular}{ὑπακοή}{obedience} εἰς δικαιοσύνην; 
\MainTextVerseMark{6}{17}χάρις δὲ τῷ θεῷ ὅτι ἦτε\FnParse{a}{imperfect active indicative 2nd plural} δοῦλοι τῆς ἁμαρτίας ὑπηκούσατε\FnParseFormGloss{b}{aorist active indicative 2nd plural}{ὑπακούω}{to obey} δὲ ἐκ καρδίας εἰς ὃν παρεδόθητε\FnParse{c}{aorist passive indicative 2nd plural} τύπον\FnParseFormGloss{d}{accusative masculine singular}{τύπος}{pattern} διδαχῆς, 
\MainTextVerseMark{6}{18}ἐλευθερωθέντες\FnParseFormGloss{a}{aorist passive participle nominative masculine plural}{ἐλευθερόω}{to set free} δὲ ἀπὸ τῆς ἁμαρτίας ἐδουλώθητε\FnParseFormGloss{b}{aorist passive indicative 2nd plural}{δουλόω}{to enslave} τῇ δικαιοσύνῃ· 
\MainTextVerseMark{6}{19}ἀνθρώπινον\FnParseFormGloss{a}{accusative neuter singular}{ἀνθρώπινος}{human} λέγω\FnParse{b}{present active indicative 1st singular} διὰ τὴν ἀσθένειαν\FnParseFormGloss{c}{accusative feminine singular}{ἀσθένεια}{weakness} τῆς σαρκὸς ὑμῶν· ὥσπερ γὰρ παρεστήσατε\FnParse{d}{aorist active indicative 2nd plural} τὰ μέλη ὑμῶν δοῦλα τῇ ἀκαθαρσίᾳ\FnParseFormGloss{e}{dative feminine singular}{ἀκαθαρσία}{impurity} καὶ τῇ ἀνομίᾳ\FnParseFormGloss{f}{dative feminine singular}{ἀνομία}{wickedness} εἰς τὴν ἀνομίαν,\FnParseFormGloss{g}{accusative feminine singular}{ἀνομία}{wickedness} οὕτως νῦν παραστήσατε\FnParse{h}{aorist active imperative 2nd plural} τὰ μέλη ὑμῶν δοῦλα τῇ δικαιοσύνῃ εἰς ἁγιασμόν.\FnParseFormGloss{i}{accusative masculine singular}{ἁγιασμός}{holiness} 
\MainTextVerseMark{6}{20}Ὅτε γὰρ δοῦλοι ἦτε\FnParse{a}{imperfect active indicative 2nd plural} τῆς ἁμαρτίας, ἐλεύθεροι\FnParseFormGloss{b}{nominative masculine plural}{ἐλεύθερος}{free} ἦτε\FnParse{c}{imperfect active indicative 2nd plural} τῇ δικαιοσύνῃ. 
\MainTextVerseMark{6}{21}τίνα οὖν καρπὸν εἴχετε\FnParse{a}{imperfect active indicative 2nd plural} τότε ἐφ’ οἷς νῦν ἐπαισχύνεσθε;\FnParseFormGloss{b}{present middle indicative 2nd plural}{ἐπαισχύνομαι}{to be ashamed of} τὸ γὰρ τέλος ἐκείνων θάνατος· 
\MainTextVerseMark{6}{22}νυνὶ\FnFormGloss{a}{νυνί}{now} δέ, ἐλευθερωθέντες\FnParseFormGloss{b}{aorist passive participle nominative masculine plural}{ἐλευθερόω}{to set free} ἀπὸ τῆς ἁμαρτίας δουλωθέντες\FnParseFormGloss{c}{aorist passive participle nominative masculine plural}{δουλόω}{to enslave} δὲ τῷ θεῷ, ἔχετε\FnParse{d}{present active indicative 2nd plural} τὸν καρπὸν ὑμῶν εἰς ἁγιασμόν,\FnParseFormGloss{e}{accusative masculine singular}{ἁγιασμός}{holiness} τὸ δὲ τέλος ζωὴν αἰώνιον. 
\MainTextVerseMark{6}{23}τὰ γὰρ ὀψώνια\FnParseFormGloss{a}{nominative neuter plural}{ὀψώνιον}{pay} τῆς ἁμαρτίας θάνατος, τὸ δὲ χάρισμα\FnParseFormGloss{b}{nominative neuter singular}{χάρισμα}{gracious gift} τοῦ θεοῦ ζωὴ αἰώνιος ἐν Χριστῷ Ἰησοῦ τῷ κυρίῳ ἡμῶν. 
\ChapterHeader{7}{Chapter 7}

\MainTextVerseMark{7}{1}Ἢ ἀγνοεῖτε,\FnParseFormGloss{a}{present active indicative 2nd plural}{ἀγνοέω}{to be ignorant} ἀδελφοί, γινώσκουσιν\FnParse{b}{present active participle dative masculine plural} γὰρ νόμον λαλῶ,\FnParse{c}{present active indicative 1st singular} ὅτι ὁ νόμος κυριεύει\FnParseFormGloss{d}{present active indicative 3rd singular}{κυριεύω}{to lord over} τοῦ ἀνθρώπου ἐφ’ ὅσον χρόνον ζῇ;\FnParse{e}{present active indicative 3rd singular} 
\MainTextVerseMark{7}{2}ἡ γὰρ ὕπανδρος\FnParseFormGloss{a}{nominative feminine singular}{ὕπανδρος}{married} γυνὴ τῷ ζῶντι\FnParse{b}{present active participle dative masculine singular} ἀνδρὶ δέδεται\FnParse{c}{perfect passive indicative 3rd singular} νόμῳ· ἐὰν δὲ ἀποθάνῃ\FnParse{d}{aorist active subjunctive 3rd singular} ὁ ἀνήρ, κατήργηται\FnParseFormGloss{e}{perfect passive indicative 3rd singular}{καταργέω}{to nullify} ἀπὸ τοῦ νόμου τοῦ ἀνδρός. 
\MainTextVerseMark{7}{3}ἄρα οὖν ζῶντος\FnParse{a}{present active participle genitive masculine singular} τοῦ ἀνδρὸς μοιχαλὶς\FnParseFormGloss{b}{nominative feminine singular}{μοιχαλίς}{adulteress} χρηματίσει\FnParseFormGloss{c}{future active indicative 3rd singular}{χρηματίζω}{to warn} ἐὰν γένηται\FnParse{d}{aorist middle subjunctive 3rd singular} ἀνδρὶ ἑτέρῳ· ἐὰν δὲ ἀποθάνῃ\FnParse{e}{aorist active subjunctive 3rd singular} ὁ ἀνήρ, ἐλευθέρα\FnParseFormGloss{f}{nominative feminine singular}{ἐλεύθερος}{free} ἐστὶν\FnParse{g}{present active indicative 3rd singular} ἀπὸ τοῦ νόμου, τοῦ μὴ εἶναι\FnParse{h}{present active infinitive} αὐτὴν μοιχαλίδα\FnParseFormGloss{i}{accusative feminine singular}{μοιχαλίς}{adulteress} γενομένην\FnParse{j}{aorist middle participle accusative feminine singular} ἀνδρὶ ἑτέρῳ. 
\MainTextVerseMark{7}{4}Ὥστε, ἀδελφοί μου, καὶ ὑμεῖς ἐθανατώθητε\FnParseFormGloss{a}{aorist passive indicative 2nd plural}{θανατόω}{to put to death} τῷ νόμῳ διὰ τοῦ σώματος τοῦ Χριστοῦ, εἰς τὸ γενέσθαι\FnParse{b}{aorist middle infinitive} ὑμᾶς ἑτέρῳ, τῷ ἐκ νεκρῶν ἐγερθέντι\FnParse{c}{aorist passive participle dative masculine singular} ἵνα καρποφορήσωμεν\FnParseFormGloss{d}{aorist active subjunctive 1st plural}{καρποφορέω}{to produce a crop} τῷ θεῷ. 
\MainTextVerseMark{7}{5}ὅτε γὰρ ἦμεν\FnParse{a}{imperfect active indicative 1st plural} ἐν τῇ σαρκί, τὰ παθήματα\FnParseFormGloss{b}{nominative neuter plural}{πάθημα}{suffering} τῶν ἁμαρτιῶν τὰ διὰ τοῦ νόμου ἐνηργεῖτο\FnParseFormGloss{c}{imperfect middle indicative 3rd singular}{ἐνεργέω}{to be at work in} ἐν τοῖς μέλεσιν ἡμῶν εἰς τὸ καρποφορῆσαι\FnParseFormGloss{d}{aorist active infinitive}{καρποφορέω}{to produce a crop} τῷ θανάτῳ· 
\MainTextVerseMark{7}{6}νυνὶ\FnFormGloss{a}{νυνί}{now} δὲ κατηργήθημεν\FnParseFormGloss{b}{aorist passive indicative 1st plural}{καταργέω}{to nullify} ἀπὸ τοῦ νόμου, ἀποθανόντες\FnParse{c}{aorist active participle nominative masculine plural} ἐν ᾧ κατειχόμεθα,\FnParseFormGloss{d}{imperfect passive indicative 1st plural}{κατέχω}{to hold back} ὥστε δουλεύειν\FnParseFormGloss{e}{present active infinitive}{δουλεύω}{to serve} ἡμᾶς ἐν καινότητι\FnParseFormGloss{f}{dative feminine singular}{καινότης}{newness} πνεύματος καὶ οὐ παλαιότητι\FnParseFormGloss{g}{dative feminine singular}{παλαιότης}{the old way} γράμματος.\FnParseFormGloss{h}{genitive neuter singular}{γράμμα}{letter} 
\MainTextVerseMark{7}{7}Τί οὖν ἐροῦμεν;\FnParse{a}{future active indicative 1st plural} ὁ νόμος ἁμαρτία; μὴ γένοιτο·\FnParse{b}{aorist middle optative 3rd singular} ἀλλὰ τὴν ἁμαρτίαν οὐκ ἔγνων\FnParse{c}{aorist active indicative 1st singular} εἰ μὴ διὰ νόμου, τήν τε γὰρ ἐπιθυμίαν οὐκ ᾔδειν\FnParse{d}{pluperfect active indicative 1st singular} εἰ μὴ ὁ νόμος ἔλεγεν·\FnParse{e}{imperfect active indicative 3rd singular} Οὐκ ἐπιθυμήσεις·\FnParseFormGloss{f}{future active indicative 2nd singular}{ἐπιθυμέω}{to long for} 
\MainTextVerseMark{7}{8}ἀφορμὴν\FnParseFormGloss{a}{accusative feminine singular}{ἀφορμή}{opportunity} δὲ λαβοῦσα\FnParse{b}{aorist active participle nominative feminine singular} ἡ ἁμαρτία διὰ τῆς ἐντολῆς κατειργάσατο\FnParseFormGloss{c}{aorist middle indicative 3rd singular}{κατεργάζομαι}{to produce} ἐν ἐμοὶ πᾶσαν ἐπιθυμίαν, χωρὶς γὰρ νόμου ἁμαρτία νεκρά. 
\MainTextVerseMark{7}{9}ἐγὼ δὲ ἔζων\FnParse{a}{imperfect active indicative 1st singular} χωρὶς νόμου ποτέ·\FnFormGloss{b}{ποτέ}{once} ἐλθούσης\FnParse{c}{aorist active participle genitive feminine singular} δὲ τῆς ἐντολῆς ἡ ἁμαρτία ἀνέζησεν,\FnParseFormGloss{d}{aorist active indicative 3rd singular}{ἀναζάω}{to become alive again} 
\MainTextVerseMark{7}{10}ἐγὼ δὲ ἀπέθανον,\FnParse{a}{aorist active indicative 1st singular} καὶ εὑρέθη\FnParse{b}{aorist passive indicative 3rd singular} μοι ἡ ἐντολὴ ἡ εἰς ζωὴν αὕτη εἰς θάνατον· 
\MainTextVerseMark{7}{11}ἡ γὰρ ἁμαρτία ἀφορμὴν\FnParseFormGloss{a}{accusative feminine singular}{ἀφορμή}{opportunity} λαβοῦσα\FnParse{b}{aorist active participle nominative feminine singular} διὰ τῆς ἐντολῆς ἐξηπάτησέν\FnParseFormGloss{c}{aorist active indicative 3rd singular}{ἐξαπατάω}{to deceive} με καὶ δι’ αὐτῆς ἀπέκτεινεν.\FnParse{d}{aorist active indicative 3rd singular} 
\MainTextVerseMark{7}{12}ὥστε ὁ μὲν νόμος ἅγιος, καὶ ἡ ἐντολὴ ἁγία καὶ δικαία καὶ ἀγαθή. 
\MainTextVerseMark{7}{13}Τὸ οὖν ἀγαθὸν ἐμοὶ ⸀ἐγένετο\FnParse{a}{aorist middle indicative 3rd singular} θάνατος; μὴ γένοιτο·\FnParse{b}{aorist middle optative 3rd singular} ἀλλὰ ἡ ἁμαρτία, ἵνα φανῇ\FnParse{c}{aorist passive subjunctive 3rd singular} ἁμαρτία διὰ τοῦ ἀγαθοῦ μοι κατεργαζομένη\FnParseFormGloss{d}{present middle participle nominative feminine singular}{κατεργάζομαι}{to produce} θάνατον· ἵνα γένηται\FnParse{e}{aorist middle subjunctive 3rd singular} καθ’ ὑπερβολὴν\FnParseFormGloss{f}{accusative feminine singular}{ὑπερβολή}{all-surpassing} ἁμαρτωλὸς ἡ ἁμαρτία διὰ τῆς ἐντολῆς. 
\MainTextVerseMark{7}{14}Οἴδαμεν\FnParse{a}{perfect active indicative 1st plural} γὰρ ὅτι ὁ νόμος πνευματικός\FnParseFormGloss{b}{nominative masculine singular}{πνευματικός}{spiritual} ἐστιν·\FnParse{c}{present active indicative 3rd singular} ἐγὼ δὲ ⸀σάρκινός\FnParseFormGloss{d}{nominative masculine singular}{σάρκινος}{fleshly} εἰμι,\FnParse{e}{present active indicative 1st singular} πεπραμένος\FnParseFormGloss{f}{perfect passive participle nominative masculine singular}{πιπράσκω}{to sell} ὑπὸ τὴν ἁμαρτίαν. 
\MainTextVerseMark{7}{15}ὃ γὰρ κατεργάζομαι\FnParseFormGloss{a}{present middle indicative 1st singular}{κατεργάζομαι}{to produce} οὐ γινώσκω·\FnParse{b}{present active indicative 1st singular} οὐ γὰρ ὃ θέλω\FnParse{c}{present active indicative 1st singular} τοῦτο πράσσω,\FnParse{d}{present active indicative 1st singular} ἀλλ’ ὃ μισῶ\FnParse{e}{present active indicative 1st singular} τοῦτο ποιῶ.\FnParse{f}{present active indicative 1st singular} 
\MainTextVerseMark{7}{16}εἰ δὲ ὃ οὐ θέλω\FnParse{a}{present active indicative 1st singular} τοῦτο ποιῶ,\FnParse{b}{present active indicative 1st singular} σύμφημι\FnParseFormGloss{c}{present active indicative 1st singular}{σύμφημι}{to agree with} τῷ νόμῳ ὅτι καλός. 
\MainTextVerseMark{7}{17}νυνὶ\FnFormGloss{a}{νυνί}{now} δὲ οὐκέτι ἐγὼ κατεργάζομαι\FnParseFormGloss{b}{present middle indicative 1st singular}{κατεργάζομαι}{to produce} αὐτὸ ἀλλὰ ἡ ⸀οἰκοῦσα\FnParseFormGloss{c}{present active participle nominative feminine singular}{οἰκέω}{to live} ἐν ἐμοὶ ἁμαρτία. 
\MainTextVerseMark{7}{18}οἶδα\FnParse{a}{perfect active indicative 1st singular} γὰρ ὅτι οὐκ οἰκεῖ\FnParseFormGloss{b}{present active indicative 3rd singular}{οἰκέω}{to live} ἐν ἐμοί, τοῦτ’ ἔστιν\FnParse{c}{present active indicative 3rd singular} ἐν τῇ σαρκί μου, ἀγαθόν· τὸ γὰρ θέλειν\FnParse{d}{present active infinitive} παράκειταί\FnParseFormGloss{e}{present middle indicative 3rd singular}{παράκειμαι}{to be present} μοι, τὸ δὲ κατεργάζεσθαι\FnParseFormGloss{f}{present middle infinitive}{κατεργάζομαι}{to produce} τὸ καλὸν ⸀οὔ·\FnFormGloss{g}{οὔ}{no} 
\MainTextVerseMark{7}{19}οὐ γὰρ ὃ θέλω\FnParse{a}{present active indicative 1st singular} ποιῶ\FnParse{b}{present active indicative 1st singular} ἀγαθόν, ἀλλὰ ὃ οὐ θέλω\FnParse{c}{present active indicative 1st singular} κακὸν τοῦτο πράσσω.\FnParse{d}{present active indicative 1st singular} 
\MainTextVerseMark{7}{20}εἰ δὲ ὃ οὐ ⸀θέλω\FnParse{a}{present active indicative 1st singular} τοῦτο ποιῶ,\FnParse{b}{present active indicative 1st singular} οὐκέτι ἐγὼ κατεργάζομαι\FnParseFormGloss{c}{present middle indicative 1st singular}{κατεργάζομαι}{to produce} αὐτὸ ἀλλὰ ἡ οἰκοῦσα\FnParseFormGloss{d}{present active participle nominative feminine singular}{οἰκέω}{to live} ἐν ἐμοὶ ἁμαρτία. 
\MainTextVerseMark{7}{21}Εὑρίσκω\FnParse{a}{present active indicative 1st singular} ἄρα τὸν νόμον τῷ θέλοντι\FnParse{b}{present active participle dative masculine singular} ἐμοὶ ποιεῖν\FnParse{c}{present active infinitive} τὸ καλὸν ὅτι ἐμοὶ τὸ κακὸν παράκειται·\FnParseFormGloss{d}{present middle indicative 3rd singular}{παράκειμαι}{to be present} 
\MainTextVerseMark{7}{22}συνήδομαι\FnParseFormGloss{a}{present middle indicative 1st singular}{συνήδομαι}{to delight in agreement} γὰρ τῷ νόμῳ τοῦ θεοῦ κατὰ τὸν ἔσω\FnFormGloss{b}{ἔσω}{in} ἄνθρωπον, 
\MainTextVerseMark{7}{23}βλέπω\FnParse{a}{present active indicative 1st singular} δὲ ἕτερον νόμον ἐν τοῖς μέλεσίν μου ἀντιστρατευόμενον\FnParseFormGloss{b}{present middle participle accusative masculine singular}{ἀντιστρατεύομαι}{to wage war against} τῷ νόμῳ τοῦ νοός\FnParseFormGloss{c}{genitive masculine singular}{νοῦς}{mind} μου καὶ αἰχμαλωτίζοντά\FnParseFormGloss{d}{present active participle accusative masculine singular}{αἰχμαλωτίζω}{to take captive} με ⸀ἐν τῷ νόμῳ τῆς ἁμαρτίας τῷ ὄντι\FnParse{e}{present active participle dative masculine singular} ἐν τοῖς μέλεσίν μου. 
\MainTextVerseMark{7}{24}ταλαίπωρος\FnParseFormGloss{a}{nominative masculine singular}{ταλαίπωρος}{wretched} ἐγὼ ἄνθρωπος· τίς με ῥύσεται\FnParseFormGloss{b}{future middle indicative 3rd singular}{ῥύομαι}{to rescue} ἐκ τοῦ σώματος τοῦ θανάτου τούτου; 
\MainTextVerseMark{7}{25}⸀χάρις τῷ θεῷ διὰ Ἰησοῦ Χριστοῦ τοῦ κυρίου ἡμῶν. Ἄρα οὖν αὐτὸς ἐγὼ τῷ μὲν νοῒ\FnParseFormGloss{a}{dative masculine singular}{νοῦς}{mind} δουλεύω\FnParseFormGloss{b}{present active indicative 1st singular}{δουλεύω}{to serve} νόμῳ θεοῦ, τῇ δὲ σαρκὶ νόμῳ ἁμαρτίας. 
\ChapterHeader{8}{Chapter 8}

\MainTextVerseMark{8}{1}Οὐδὲν ἄρα νῦν κατάκριμα\FnParseFormGloss{a}{nominative neuter singular}{κατάκριμα}{condemnation} τοῖς ἐν Χριστῷ ⸀Ἰησοῦ· 
\MainTextVerseMark{8}{2}ὁ γὰρ νόμος τοῦ πνεύματος τῆς ζωῆς ἐν Χριστῷ Ἰησοῦ ἠλευθέρωσέν\FnParseFormGloss{a}{aorist active indicative 3rd singular}{ἐλευθερόω}{to set free} ⸀σε ἀπὸ τοῦ νόμου τῆς ἁμαρτίας καὶ τοῦ θανάτου. 
\MainTextVerseMark{8}{3}τὸ γὰρ ἀδύνατον\FnParseFormGloss{a}{accusative neuter singular}{ἀδύνατος}{impossible} τοῦ νόμου, ἐν ᾧ ἠσθένει\FnParse{b}{imperfect active indicative 3rd singular} διὰ τῆς σαρκός, ὁ θεὸς τὸν ἑαυτοῦ υἱὸν πέμψας\FnParse{c}{aorist active participle nominative masculine singular} ἐν ὁμοιώματι\FnParseFormGloss{d}{dative neuter singular}{ὁμοίωμα}{likeness} σαρκὸς ἁμαρτίας καὶ περὶ ἁμαρτίας κατέκρινε\FnParseFormGloss{e}{aorist active indicative 3rd singular}{κατακρίνω}{to condemn} τὴν ἁμαρτίαν ἐν τῇ σαρκί, 
\MainTextVerseMark{8}{4}ἵνα τὸ δικαίωμα\FnParseFormGloss{a}{nominative neuter singular}{δικαίωμα}{regulation} τοῦ νόμου πληρωθῇ\FnParse{b}{aorist passive subjunctive 3rd singular} ἐν ἡμῖν τοῖς μὴ κατὰ σάρκα περιπατοῦσιν\FnParse{c}{present active participle dative masculine plural} ἀλλὰ κατὰ πνεῦμα· 
\MainTextVerseMark{8}{5}οἱ γὰρ κατὰ σάρκα ὄντες\FnParse{a}{present active participle nominative masculine plural} τὰ τῆς σαρκὸς φρονοῦσιν,\FnParseFormGloss{b}{present active indicative 3rd plural}{φρονέω}{to think} οἱ δὲ κατὰ πνεῦμα τὰ τοῦ πνεύματος. 
\MainTextVerseMark{8}{6}τὸ γὰρ φρόνημα\FnParseFormGloss{a}{nominative neuter singular}{φρόνημα}{mind} τῆς σαρκὸς θάνατος, τὸ δὲ φρόνημα\FnParseFormGloss{b}{nominative neuter singular}{φρόνημα}{mind} τοῦ πνεύματος ζωὴ καὶ εἰρήνη· 
\MainTextVerseMark{8}{7}διότι\FnFormGloss{a}{διότι}{therefore} τὸ φρόνημα\FnParseFormGloss{b}{nominative neuter singular}{φρόνημα}{mind} τῆς σαρκὸς ἔχθρα\FnParseFormGloss{c}{nominative feminine singular}{ἔχθρα}{hostility} εἰς θεόν, τῷ γὰρ νόμῳ τοῦ θεοῦ οὐχ ὑποτάσσεται,\FnParse{d}{present passive indicative 3rd singular} οὐδὲ γὰρ δύναται·\FnParse{e}{present middle indicative 3rd singular} 
\MainTextVerseMark{8}{8}οἱ δὲ ἐν σαρκὶ ὄντες\FnParse{a}{present active participle nominative masculine plural} θεῷ ἀρέσαι\FnParseFormGloss{b}{aorist active infinitive}{ἀρέσκω}{to please} οὐ δύνανται.\FnParse{c}{present middle indicative 3rd plural} 
\MainTextVerseMark{8}{9}Ὑμεῖς δὲ οὐκ ἐστὲ\FnParse{a}{present active indicative 2nd plural} ἐν σαρκὶ ἀλλὰ ἐν πνεύματι, εἴπερ\FnFormGloss{b}{εἴπερ}{if indeed} πνεῦμα θεοῦ οἰκεῖ\FnParseFormGloss{c}{present active indicative 3rd singular}{οἰκέω}{to live} ἐν ὑμῖν. εἰ δέ τις πνεῦμα Χριστοῦ οὐκ ἔχει,\FnParse{d}{present active indicative 3rd singular} οὗτος οὐκ ἔστιν\FnParse{e}{present active indicative 3rd singular} αὐτοῦ. 
\MainTextVerseMark{8}{10}εἰ δὲ Χριστὸς ἐν ὑμῖν, τὸ μὲν σῶμα νεκρὸν διὰ ἁμαρτίαν, τὸ δὲ πνεῦμα ζωὴ διὰ δικαιοσύνην. 
\MainTextVerseMark{8}{11}εἰ δὲ τὸ πνεῦμα τοῦ ἐγείραντος\FnParse{a}{aorist active participle genitive masculine singular} ⸀τὸν Ἰησοῦν ἐκ νεκρῶν οἰκεῖ\FnParseFormGloss{b}{present active indicative 3rd singular}{οἰκέω}{to live} ἐν ὑμῖν, ὁ ἐγείρας\FnParse{c}{aorist active participle nominative masculine singular} ⸂ἐκ νεκρῶν Χριστὸν Ἰησοῦν⸃ ζῳοποιήσει\FnParseFormGloss{d}{future active indicative 3rd singular}{ζῳοποιέω}{to make alive} καὶ τὰ θνητὰ\FnParseFormGloss{e}{accusative neuter plural}{θνητός}{mortal} σώματα ὑμῶν διὰ ⸂τὸ ἐνοικοῦν\FnParseFormGloss{f}{present active participle accusative neuter singular}{ἐνοικέω}{to live in} αὐτοῦ πνεῦμα⸃ ἐν ὑμῖν. 
\MainTextVerseMark{8}{12}Ἄρα οὖν, ἀδελφοί, ὀφειλέται\FnParseFormGloss{a}{nominative masculine plural}{ὀφειλέτης}{debtor} ἐσμέν,\FnParse{b}{present active indicative 1st plural} οὐ τῇ σαρκὶ τοῦ κατὰ σάρκα ζῆν,\FnParse{c}{present active infinitive} 
\MainTextVerseMark{8}{13}εἰ γὰρ κατὰ σάρκα ζῆτε\FnParse{a}{present active indicative 2nd plural} μέλλετε\FnParse{b}{present active indicative 2nd plural} ἀποθνῄσκειν,\FnParse{c}{present active infinitive} εἰ δὲ πνεύματι τὰς πράξεις\FnParseFormGloss{d}{accusative feminine plural}{πρᾶξις}{deed} τοῦ σώματος θανατοῦτε,\FnParseFormGloss{e}{present active indicative 2nd plural}{θανατόω}{to put to death} ζήσεσθε.\FnParse{f}{future middle indicative 2nd plural} 
\MainTextVerseMark{8}{14}ὅσοι γὰρ πνεύματι θεοῦ ἄγονται,\FnParse{a}{present passive indicative 3rd plural} οὗτοι ⸂υἱοί εἰσιν\FnParse{b}{present active indicative 3rd plural} θεοῦ⸃. 
\MainTextVerseMark{8}{15}οὐ γὰρ ἐλάβετε\FnParse{a}{aorist active indicative 2nd plural} πνεῦμα δουλείας\FnParseFormGloss{b}{genitive feminine singular}{δουλεία}{slavery} πάλιν εἰς φόβον, ἀλλὰ ἐλάβετε\FnParse{c}{aorist active indicative 2nd plural} πνεῦμα υἱοθεσίας\FnParseFormGloss{d}{genitive feminine singular}{υἱοθεσία}{adoption as sons} ἐν ᾧ κράζομεν·\FnParse{e}{present active indicative 1st plural} Αββα\FnParseFormGloss{f}{masculine singular}{αββα}{father} ὁ πατήρ· 
\MainTextVerseMark{8}{16}αὐτὸ τὸ πνεῦμα συμμαρτυρεῖ\FnParseFormGloss{a}{present active indicative 3rd singular}{συμμαρτυρέω}{to testify with} τῷ πνεύματι ἡμῶν ὅτι ἐσμὲν\FnParse{b}{present active indicative 1st plural} τέκνα θεοῦ. 
\MainTextVerseMark{8}{17}εἰ δὲ τέκνα, καὶ κληρονόμοι·\FnParseFormGloss{a}{nominative masculine plural}{κληρονόμος}{heir} κληρονόμοι\FnParseFormGloss{b}{nominative masculine plural}{κληρονόμος}{heir} μὲν θεοῦ, συγκληρονόμοι\FnParseFormGloss{c}{nominative masculine plural}{συγκληρονόμος}{inheriting together} δὲ Χριστοῦ, εἴπερ\FnFormGloss{d}{εἴπερ}{if indeed} συμπάσχομεν\FnParseFormGloss{e}{present active indicative 1st plural}{συμπάσχω}{to suffer with} ἵνα καὶ συνδοξασθῶμεν.\FnParseFormGloss{f}{aorist passive subjunctive 1st plural}{συνδοξάζομαι}{to be glorified with} 
\MainTextVerseMark{8}{18}Λογίζομαι\FnParse{a}{present middle indicative 1st singular} γὰρ ὅτι οὐκ ἄξια τὰ παθήματα\FnParseFormGloss{b}{nominative neuter plural}{πάθημα}{suffering} τοῦ νῦν καιροῦ πρὸς τὴν μέλλουσαν\FnParse{c}{present active participle accusative feminine singular} δόξαν ἀποκαλυφθῆναι\FnParseFormGloss{d}{aorist passive infinitive}{ἀποκαλύπτω}{to reveal} εἰς ἡμᾶς. 
\MainTextVerseMark{8}{19}ἡ γὰρ ἀποκαραδοκία\FnParseFormGloss{a}{nominative feminine singular}{ἀποκαραδοκία}{eager expectation} τῆς κτίσεως\FnParseFormGloss{b}{genitive feminine singular}{κτίσις}{creation} τὴν ἀποκάλυψιν\FnParseFormGloss{c}{accusative feminine singular}{ἀποκάλυψις}{revelation} τῶν υἱῶν τοῦ θεοῦ ἀπεκδέχεται·\FnParseFormGloss{d}{present middle indicative 3rd singular}{ἀπεκδέχομαι}{wait eagerly for} 
\MainTextVerseMark{8}{20}τῇ γὰρ ματαιότητι\FnParseFormGloss{a}{dative feminine singular}{ματαιότης}{emptiness} ἡ κτίσις\FnParseFormGloss{b}{nominative feminine singular}{κτίσις}{creation} ὑπετάγη,\FnParse{c}{aorist passive indicative 3rd singular} οὐχ ἑκοῦσα\FnParseFormGloss{d}{nominative feminine singular}{ἑκών}{voluntarily} ἀλλὰ διὰ τὸν ὑποτάξαντα,\FnParse{e}{aorist active participle accusative masculine singular} ⸀ἐφ’ ἑλπίδι 
\MainTextVerseMark{8}{21}ὅτι καὶ αὐτὴ ἡ κτίσις\FnParseFormGloss{a}{nominative feminine singular}{κτίσις}{creation} ἐλευθερωθήσεται\FnParseFormGloss{b}{future passive indicative 3rd singular}{ἐλευθερόω}{to set free} ἀπὸ τῆς δουλείας\FnParseFormGloss{c}{genitive feminine singular}{δουλεία}{slavery} τῆς φθορᾶς\FnParseFormGloss{d}{genitive feminine singular}{φθορά}{perishableness} εἰς τὴν ἐλευθερίαν\FnParseFormGloss{e}{accusative feminine singular}{ἐλευθερία}{freedom} τῆς δόξης τῶν τέκνων τοῦ θεοῦ. 
\MainTextVerseMark{8}{22}οἴδαμεν\FnParse{a}{perfect active indicative 1st plural} γὰρ ὅτι πᾶσα ἡ κτίσις\FnParseFormGloss{b}{nominative feminine singular}{κτίσις}{creation} συστενάζει\FnParseFormGloss{c}{present active indicative 3rd singular}{συστενάζω}{to join in groaning} καὶ συνωδίνει\FnParseFormGloss{d}{present active indicative 3rd singular}{συνωδίνω}{suffer together} ἄχρι τοῦ νῦν· 
\MainTextVerseMark{8}{23}οὐ μόνον δέ, ἀλλὰ ⸂καὶ αὐτοὶ⸃ τὴν ἀπαρχὴν\FnParseFormGloss{a}{accusative feminine singular}{ἀπαρχή}{firstfruits} τοῦ πνεύματος ἔχοντες\FnParse{b}{present active participle nominative masculine plural} ⸂ἡμεῖς καὶ⸃ αὐτοὶ ἐν ἑαυτοῖς στενάζομεν,\FnParseFormGloss{c}{present active indicative 1st plural}{στενάζω}{to groan} υἱοθεσίαν\FnParseFormGloss{d}{accusative feminine singular}{υἱοθεσία}{adoption as sons} ἀπεκδεχόμενοι\FnParseFormGloss{e}{present middle participle nominative masculine plural}{ἀπεκδέχομαι}{wait eagerly for} τὴν ἀπολύτρωσιν\FnParseFormGloss{f}{accusative feminine singular}{ἀπολύτρωσις}{redemption} τοῦ σώματος ἡμῶν. 
\MainTextVerseMark{8}{24}τῇ γὰρ ἐλπίδι ἐσώθημεν·\FnParse{a}{aorist passive indicative 1st plural} ἐλπὶς δὲ βλεπομένη\FnParse{b}{present passive participle nominative feminine singular} οὐκ ἔστιν\FnParse{c}{present active indicative 3rd singular} ἐλπίς, ὃ γὰρ βλέπει\FnParse{d}{present active indicative 3rd singular} ⸀τίς ἐλπίζει;\FnParse{e}{present active indicative 3rd singular} 
\MainTextVerseMark{8}{25}εἰ δὲ ὃ οὐ βλέπομεν\FnParse{a}{present active indicative 1st plural} ἐλπίζομεν,\FnParse{b}{present active indicative 1st plural} δι’ ὑπομονῆς ἀπεκδεχόμεθα.\FnParseFormGloss{c}{present middle indicative 1st plural}{ἀπεκδέχομαι}{wait eagerly for} 
\MainTextVerseMark{8}{26}Ὡσαύτως\FnFormGloss{a}{ὡσαύτως}{in the same way} δὲ καὶ τὸ πνεῦμα συναντιλαμβάνεται\FnParseFormGloss{b}{present middle indicative 3rd singular}{συναντιλαμβάνομαι}{to help} ⸂τῇ ἀσθενείᾳ⸃\FnParseFormGloss{c}{dative feminine singular}{ἀσθένεια}{weakness} ἡμῶν· τὸ γὰρ τί ⸀προσευξώμεθα\FnParse{d}{aorist middle subjunctive 1st plural} καθὸ\FnFormGloss{e}{καθό}{insofar as} δεῖ\FnParse{f}{present active indicative 3rd singular} οὐκ οἴδαμεν,\FnParse{g}{perfect active indicative 1st plural} ἀλλὰ αὐτὸ τὸ πνεῦμα ⸀ὑπερεντυγχάνει\FnParseFormGloss{h}{present active indicative 3rd singular}{ὑπερεντυγχάνω}{to intercede} στεναγμοῖς\FnParseFormGloss{i}{dative masculine plural}{στεναγμός}{groan} ἀλαλήτοις,\FnParseFormGloss{j}{dative masculine plural}{ἀλάλητος}{inexpressible} 
\MainTextVerseMark{8}{27}ὁ δὲ ἐραυνῶν\FnParseFormGloss{a}{present active participle nominative masculine singular}{ἐραυνάω}{to search} τὰς καρδίας οἶδεν\FnParse{b}{perfect active indicative 3rd singular} τί τὸ φρόνημα\FnParseFormGloss{c}{nominative neuter singular}{φρόνημα}{mind} τοῦ πνεύματος, ὅτι κατὰ θεὸν ἐντυγχάνει\FnParseFormGloss{d}{present active indicative 3rd singular}{ἐντυγχάνω}{to intercede} ὑπὲρ ἁγίων. 
\MainTextVerseMark{8}{28}Οἴδαμεν\FnParse{a}{perfect active indicative 1st plural} δὲ ὅτι τοῖς ἀγαπῶσι\FnParse{b}{present active participle dative masculine plural} τὸν θεὸν πάντα ⸀συνεργεῖ\FnParseFormGloss{c}{present active indicative 3rd singular}{συνεργέω}{to work together} εἰς ἀγαθόν, τοῖς κατὰ πρόθεσιν\FnParseFormGloss{d}{accusative feminine singular}{πρόθεσις}{plan} κλητοῖς\FnParseFormGloss{e}{dative masculine plural}{κλητός}{called} οὖσιν.\FnParse{f}{present active participle dative masculine plural} 
\MainTextVerseMark{8}{29}ὅτι οὓς προέγνω,\FnParseFormGloss{a}{aorist active indicative 3rd singular}{προγινώσκω}{to know beforehand} καὶ προώρισεν\FnParseFormGloss{b}{aorist active indicative 3rd singular}{προορίζω}{to predestine} συμμόρφους\FnParseFormGloss{c}{accusative masculine plural}{σύμμορφος}{conformed} τῆς εἰκόνος\FnParseFormGloss{d}{genitive feminine singular}{εἰκών}{image} τοῦ υἱοῦ αὐτοῦ, εἰς τὸ εἶναι\FnParse{e}{present active infinitive} αὐτὸν πρωτότοκον\FnParseFormGloss{f}{accusative masculine singular}{πρωτότοκος}{firstborn} ἐν πολλοῖς ἀδελφοῖς· 
\MainTextVerseMark{8}{30}οὓς δὲ προώρισεν,\FnParseFormGloss{a}{aorist active indicative 3rd singular}{προορίζω}{to predestine} τούτους καὶ ἐκάλεσεν·\FnParse{b}{aorist active indicative 3rd singular} καὶ οὓς ἐκάλεσεν,\FnParse{c}{aorist active indicative 3rd singular} τούτους καὶ ἐδικαίωσεν·\FnParse{d}{aorist active indicative 3rd singular} οὓς δὲ ἐδικαίωσεν,\FnParse{e}{aorist active indicative 3rd singular} τούτους καὶ ἐδόξασεν.\FnParse{f}{aorist active indicative 3rd singular} 
\MainTextVerseMark{8}{31}Τί οὖν ἐροῦμεν\FnParse{a}{future active indicative 1st plural} πρὸς ταῦτα; εἰ ὁ θεὸς ὑπὲρ ἡμῶν, τίς καθ’ ἡμῶν; 
\MainTextVerseMark{8}{32}ὅς γε\FnFormGloss{a}{γέ}{indeed} τοῦ ἰδίου υἱοῦ οὐκ ἐφείσατο,\FnParseFormGloss{b}{aorist middle indicative 3rd singular}{φείδομαι}{to spare} ἀλλὰ ὑπὲρ ἡμῶν πάντων παρέδωκεν\FnParse{c}{aorist active indicative 3rd singular} αὐτόν, πῶς οὐχὶ καὶ σὺν αὐτῷ τὰ πάντα ἡμῖν χαρίσεται;\FnParseFormGloss{d}{future middle indicative 3rd singular}{χαρίζομαι}{to give grace} 
\MainTextVerseMark{8}{33}τίς ἐγκαλέσει\FnParseFormGloss{a}{future active indicative 3rd singular}{ἐγκαλέω}{to bring charges} κατὰ ἐκλεκτῶν\FnParseFormGloss{b}{genitive masculine plural}{ἐκλεκτός}{elect} θεοῦ; θεὸς ὁ δικαιῶν·\FnParse{c}{present active participle nominative masculine singular} 
\MainTextVerseMark{8}{34}τίς ὁ ⸀κατακρινῶν;\FnParseFormGloss{a}{present active participle nominative masculine singular}{κατακρίνω}{to condemn} ⸀Χριστὸς ὁ ἀποθανών,\FnParse{b}{aorist active participle nominative masculine singular} μᾶλλον ⸀δὲ ⸀ἐγερθείς,\FnParse{c}{aorist passive participle nominative masculine singular} ὅς ⸀καί ἐστιν\FnParse{d}{present active indicative 3rd singular} ἐν δεξιᾷ τοῦ θεοῦ, ὃς καὶ ἐντυγχάνει\FnParseFormGloss{e}{present active indicative 3rd singular}{ἐντυγχάνω}{to intercede} ὑπὲρ ἡμῶν· 
\MainTextVerseMark{8}{35}τίς ἡμᾶς χωρίσει\FnParseFormGloss{a}{future active indicative 3rd singular}{χωρίζω}{to divide} ἀπὸ τῆς ἀγάπης τοῦ Χριστοῦ; θλῖψις ἢ στενοχωρία\FnParseFormGloss{b}{nominative feminine singular}{στενοχωρία}{distress} ἢ διωγμὸς\FnParseFormGloss{c}{nominative masculine singular}{διωγμός}{persecution} ἢ λιμὸς\FnParseFormGloss{d}{nominative masculine singular}{λιμός}{hunger} ἢ γυμνότης\FnParseFormGloss{e}{nominative feminine singular}{γυμνότης}{nakedness} ἢ κίνδυνος\FnParseFormGloss{f}{nominative masculine singular}{κίνδυνος}{danger} ἢ μάχαιρα;\FnParseFormGloss{g}{nominative feminine singular}{μάχαιρα}{sword} 
\MainTextVerseMark{8}{36}καθὼς γέγραπται\FnParse{a}{perfect passive indicative 3rd singular} ὅτι Ἕνεκεν\FnFormGloss{b}{ἕνεκεν}{for} σοῦ θανατούμεθα\FnParseFormGloss{c}{present passive indicative 1st plural}{θανατόω}{to put to death} ὅλην τὴν ἡμέραν, ἐλογίσθημεν\FnParse{d}{aorist passive indicative 1st plural} ὡς πρόβατα σφαγῆς.\FnParseFormGloss{e}{genitive feminine singular}{σφαγή}{slaughter} 
\MainTextVerseMark{8}{37}ἀλλ’ ἐν τούτοις πᾶσιν ὑπερνικῶμεν\FnParseFormGloss{a}{present active indicative 1st plural}{ὑπερνικάω}{to thoroughly conquer} διὰ τοῦ ἀγαπήσαντος\FnParse{b}{aorist active participle genitive masculine singular} ἡμᾶς. 
\MainTextVerseMark{8}{38}πέπεισμαι\FnParse{a}{perfect passive indicative 1st singular} γὰρ ὅτι οὔτε θάνατος οὔτε ζωὴ οὔτε ἄγγελοι οὔτε ἀρχαὶ οὔτε ⸂ἐνεστῶτα\FnParseFormGloss{b}{perfect active participle nominative neuter plural}{ἐνίστημι}{to be present} οὔτε μέλλοντα\FnParse{c}{present active participle nominative neuter plural} οὔτε δυνάμεις⸃ 
\MainTextVerseMark{8}{39}οὔτε ὕψωμα\FnParseFormGloss{a}{nominative neuter singular}{ὕψωμα}{height} οὔτε βάθος\FnParseFormGloss{b}{nominative neuter singular}{βάθος}{depth} οὔτε τις κτίσις\FnParseFormGloss{c}{nominative feminine singular}{κτίσις}{creation} ἑτέρα δυνήσεται\FnParse{d}{future middle indicative 3rd singular} ἡμᾶς χωρίσαι\FnParseFormGloss{e}{aorist active infinitive}{χωρίζω}{to divide} ἀπὸ τῆς ἀγάπης τοῦ θεοῦ τῆς ἐν Χριστῷ Ἰησοῦ τῷ κυρίῳ ἡμῶν. 
\ChapterHeader{9}{Chapter 9}

\MainTextVerseMark{9}{1}Ἀλήθειαν λέγω\FnParse{a}{present active indicative 1st singular} ἐν Χριστῷ, οὐ ψεύδομαι,\FnParseFormGloss{b}{present middle indicative 1st singular}{ψεύδομαι}{to lie} συμμαρτυρούσης\FnParseFormGloss{c}{present active participle genitive feminine singular}{συμμαρτυρέω}{to testify with} μοι τῆς συνειδήσεώς μου ἐν πνεύματι ἁγίῳ, 
\MainTextVerseMark{9}{2}ὅτι λύπη\FnParseFormGloss{a}{nominative feminine singular}{λύπη}{sorrow} μοί ἐστιν\FnParse{b}{present active indicative 3rd singular} μεγάλη καὶ ἀδιάλειπτος\FnParseFormGloss{c}{nominative feminine singular}{ἀδιάλειπτος}{constant} ὀδύνη\FnParseFormGloss{d}{nominative feminine singular}{ὀδύνη}{anguish} τῇ καρδίᾳ μου· 
\MainTextVerseMark{9}{3}ηὐχόμην\FnParseFormGloss{a}{imperfect middle indicative 1st singular}{εὔχομαι}{to pray for} γὰρ ⸂ἀνάθεμα\FnParseFormGloss{b}{nominative neuter singular}{ἀνάθεμα}{curse} εἶναι\FnParse{c}{present active infinitive} αὐτὸς ἐγὼ⸃ ἀπὸ τοῦ Χριστοῦ ὑπὲρ τῶν ἀδελφῶν μου τῶν συγγενῶν\FnParseFormGloss{d}{genitive masculine plural}{συγγενής}{family} μου κατὰ σάρκα, 
\MainTextVerseMark{9}{4}οἵτινές εἰσιν\FnParse{a}{present active indicative 3rd plural} Ἰσραηλῖται,\FnParseFormGloss{b}{nominative masculine plural}{Ἰσραηλίτης}{Israelite} ὧν ἡ υἱοθεσία\FnParseFormGloss{c}{nominative feminine singular}{υἱοθεσία}{adoption as sons} καὶ ἡ δόξα καὶ αἱ διαθῆκαι καὶ ἡ νομοθεσία\FnParseFormGloss{d}{nominative feminine singular}{νομοθεσία}{law} καὶ ἡ λατρεία\FnParseFormGloss{e}{nominative feminine singular}{λατρεία}{worship} καὶ αἱ ἐπαγγελίαι, 
\MainTextVerseMark{9}{5}ὧν οἱ πατέρες, καὶ ἐξ ὧν ὁ χριστὸς τὸ κατὰ σάρκα, ὁ ὢν\FnParse{a}{present active participle nominative masculine singular} ἐπὶ πάντων, θεὸς εὐλογητὸς\FnParseFormGloss{b}{nominative masculine singular}{εὐλογητός}{worthy of being praised} εἰς τοὺς αἰῶνας· ἀμήν. 
\MainTextVerseMark{9}{6}Οὐχ οἷον\FnParseFormGloss{a}{nominative neuter singular}{οἷος}{what sort of} δὲ ὅτι ἐκπέπτωκεν\FnParseFormGloss{b}{perfect active indicative 3rd singular}{ἐκπίπτω}{to fall off} ὁ λόγος τοῦ θεοῦ. οὐ γὰρ πάντες οἱ ἐξ Ἰσραήλ, οὗτοι Ἰσραήλ· 
\MainTextVerseMark{9}{7}οὐδ’ ὅτι εἰσὶν\FnParse{a}{present active indicative 3rd plural} σπέρμα Ἀβραάμ, πάντες τέκνα, ἀλλ’· Ἐν Ἰσαὰκ\FnParseFormGloss{b}{dative masculine singular}{Ἰσαάκ}{Isaac} κληθήσεταί\FnParse{c}{future passive indicative 3rd singular} σοι σπέρμα. 
\MainTextVerseMark{9}{8}τοῦτ’ ἔστιν,\FnParse{a}{present active indicative 3rd singular} οὐ τὰ τέκνα τῆς σαρκὸς ταῦτα τέκνα τοῦ θεοῦ, ἀλλὰ τὰ τέκνα τῆς ἐπαγγελίας λογίζεται\FnParse{b}{present passive indicative 3rd singular} εἰς σπέρμα· 
\MainTextVerseMark{9}{9}ἐπαγγελίας γὰρ ὁ λόγος οὗτος· Κατὰ τὸν καιρὸν τοῦτον ἐλεύσομαι\FnParse{a}{future middle indicative 1st singular} καὶ ἔσται\FnParse{b}{future middle indicative 3rd singular} τῇ Σάρρᾳ\FnParseFormGloss{c}{dative feminine singular}{Σάρρα}{Sarah} υἱός. 
\MainTextVerseMark{9}{10}οὐ μόνον δέ, ἀλλὰ καὶ Ῥεβέκκα\FnParseFormGloss{a}{nominative feminine singular}{Ῥεβέκκα}{Rebekah} ἐξ ἑνὸς κοίτην\FnParseFormGloss{b}{accusative feminine singular}{κοίτη}{bed} ἔχουσα,\FnParse{c}{present active participle nominative feminine singular} Ἰσαὰκ\FnParseFormGloss{d}{genitive masculine singular}{Ἰσαάκ}{Isaac} τοῦ πατρὸς ἡμῶν· 
\MainTextVerseMark{9}{11}μήπω\FnFormGloss{a}{μήπω}{not yet} γὰρ γεννηθέντων\FnParse{b}{aorist passive participle genitive masculine plural} μηδὲ πραξάντων\FnParse{c}{aorist active participle genitive masculine plural} τι ἀγαθὸν ἢ ⸀φαῦλον,\FnParseFormGloss{d}{accusative neuter singular}{φαῦλος}{evil} ἵνα ἡ κατ’ ἐκλογὴν\FnParseFormGloss{e}{accusative feminine singular}{ἐκλογή}{election} πρόθεσις\FnParseFormGloss{f}{nominative feminine singular}{πρόθεσις}{plan} τοῦ θεοῦ μένῃ,\FnParse{g}{present active subjunctive 3rd singular} 
\MainTextVerseMark{9}{12}οὐκ ἐξ ἔργων ἀλλ’ ἐκ τοῦ καλοῦντος,\FnParse{a}{present active participle genitive masculine singular} ἐρρέθη\FnParse{b}{aorist passive indicative 3rd singular} αὐτῇ ὅτι Ὁ μείζων δουλεύσει\FnParseFormGloss{c}{future active indicative 3rd singular}{δουλεύω}{to serve} τῷ ἐλάσσονι·\FnParseFormGloss{d}{dative masculine singular}{ἐλάσσων}{lesser} 
\MainTextVerseMark{9}{13}⸀καθὼς γέγραπται·\FnParse{a}{perfect passive indicative 3rd singular} Τὸν Ἰακὼβ\FnParseFormGloss{b}{accusative masculine singular}{Ἰακώβ}{Jacob} ἠγάπησα,\FnParse{c}{aorist active indicative 1st singular} τὸν δὲ Ἠσαῦ\FnParseFormGloss{d}{accusative masculine singular}{Ἠσαῦ}{Esau} ἐμίσησα.\FnParse{e}{aorist active indicative 1st singular} 
\MainTextVerseMark{9}{14}Τί οὖν ἐροῦμεν;\FnParse{a}{future active indicative 1st plural} μὴ ἀδικία\FnParseFormGloss{b}{nominative feminine singular}{ἀδικία}{wickedness} παρὰ τῷ θεῷ; μὴ γένοιτο·\FnParse{c}{aorist middle optative 3rd singular} 
\MainTextVerseMark{9}{15}τῷ ⸂Μωϋσεῖ γὰρ⸃ λέγει·\FnParse{a}{present active indicative 3rd singular} Ἐλεήσω\FnParse{b}{future active indicative 1st singular} ὃν ἂν ἐλεῶ,\FnParse{c}{present active subjunctive 1st singular} καὶ οἰκτιρήσω\FnParseFormGloss{d}{future active indicative 1st singular}{οἰκτίρω}{to have compassion on} ὃν ἂν οἰκτίρω.\FnParseFormGloss{e}{present active subjunctive 1st singular}{οἰκτίρω}{to have compassion on} 
\MainTextVerseMark{9}{16}ἄρα οὖν οὐ τοῦ θέλοντος\FnParse{a}{present active participle genitive masculine singular} οὐδὲ τοῦ τρέχοντος\FnParseFormGloss{b}{present active participle genitive masculine singular}{τρέχω}{to run} ἀλλὰ τοῦ ἐλεῶντος\FnParse{c}{present active participle genitive masculine singular} θεοῦ. 
\MainTextVerseMark{9}{17}λέγει\FnParse{a}{present active indicative 3rd singular} γὰρ ἡ γραφὴ τῷ Φαραὼ\FnParseFormGloss{b}{dative masculine singular}{Φαραώ}{Pharaoh} ὅτι Εἰς αὐτὸ τοῦτο ἐξήγειρά\FnParseFormGloss{c}{aorist active indicative 1st singular}{ἐξεγείρω}{to raise} σε ὅπως ἐνδείξωμαι\FnParseFormGloss{d}{aorist middle subjunctive 1st singular}{ἐνδείκνυμι}{to show} ἐν σοὶ τὴν δύναμίν μου, καὶ ὅπως διαγγελῇ\FnParseFormGloss{e}{aorist passive subjunctive 3rd singular}{διαγγέλλω}{to proclaim} τὸ ὄνομά μου ἐν πάσῃ τῇ γῇ. 
\MainTextVerseMark{9}{18}ἄρα οὖν ὃν θέλει\FnParse{a}{present active indicative 3rd singular} ἐλεεῖ,\FnParse{b}{present active indicative 3rd singular} ὃν δὲ θέλει\FnParse{c}{present active indicative 3rd singular} σκληρύνει.\FnParseFormGloss{d}{present active indicative 3rd singular}{σκληρύνω}{to harden} 
\MainTextVerseMark{9}{19}Ἐρεῖς\FnParse{a}{future active indicative 2nd singular} ⸂μοι οὖν⸃· Τί ⸀οὖν ἔτι μέμφεται;\FnParseFormGloss{b}{present middle indicative 3rd singular}{μέμφομαι}{to find fault with} τῷ γὰρ βουλήματι\FnParseFormGloss{c}{dative neuter singular}{βούλημα}{plan} αὐτοῦ τίς ἀνθέστηκεν;\FnParseFormGloss{d}{perfect active indicative 3rd singular}{ἀνθίστημι}{to resist} 
\MainTextVerseMark{9}{20}⸂ὦ\FnFormGloss{a}{ὦ}{O!} ἄνθρωπε, μενοῦνγε⸃\FnFormGloss{b}{μενοῦνγε}{rather} σὺ τίς εἶ\FnParse{c}{present active indicative 2nd singular} ὁ ἀνταποκρινόμενος\FnParseFormGloss{d}{present middle participle nominative masculine singular}{ἀνταποκρίνομαι}{to talk back} τῷ θεῷ; μὴ ἐρεῖ\FnParse{e}{future active indicative 3rd singular} τὸ πλάσμα\FnParseFormGloss{f}{nominative neuter singular}{πλάσμα}{what is formed} τῷ πλάσαντι\FnParseFormGloss{g}{aorist active participle dative masculine singular}{πλάσσω}{to form} Τί με ἐποίησας\FnParse{h}{aorist active indicative 2nd singular} οὕτως; 
\MainTextVerseMark{9}{21}ἢ οὐκ ἔχει\FnParse{a}{present active indicative 3rd singular} ἐξουσίαν ὁ κεραμεὺς\FnParseFormGloss{b}{nominative masculine singular}{κεραμεύς}{potter} τοῦ πηλοῦ\FnParseFormGloss{c}{genitive masculine singular}{πηλός}{mud} ἐκ τοῦ αὐτοῦ φυράματος\FnParseFormGloss{d}{genitive neuter singular}{φύραμα}{lump} ποιῆσαι\FnParse{e}{aorist active infinitive} ὃ μὲν εἰς τιμὴν σκεῦος\FnParseFormGloss{f}{accusative neuter singular}{σκεῦος}{possession} ὃ δὲ εἰς ἀτιμίαν;\FnParseFormGloss{g}{accusative feminine singular}{ἀτιμία}{dishonor} 
\MainTextVerseMark{9}{22}εἰ δὲ θέλων\FnParse{a}{present active participle nominative masculine singular} ὁ θεὸς ἐνδείξασθαι\FnParseFormGloss{b}{aorist middle infinitive}{ἐνδείκνυμι}{to show} τὴν ὀργὴν καὶ γνωρίσαι\FnParseFormGloss{c}{aorist active infinitive}{γνωρίζω}{to make known} τὸ δυνατὸν αὐτοῦ ἤνεγκεν\FnParse{d}{aorist active indicative 3rd singular} ἐν πολλῇ μακροθυμίᾳ\FnParseFormGloss{e}{dative feminine singular}{μακροθυμία}{patience} σκεύη\FnParseFormGloss{f}{accusative neuter plural}{σκεῦος}{possession} ὀργῆς κατηρτισμένα\FnParseFormGloss{g}{perfect passive participle accusative neuter plural}{καταρτίζω}{to restore} εἰς ἀπώλειαν,\FnParseFormGloss{h}{accusative feminine singular}{ἀπώλεια}{destruction} 
\MainTextVerseMark{9}{23}⸀καὶ ἵνα γνωρίσῃ\FnParseFormGloss{a}{aorist active subjunctive 3rd singular}{γνωρίζω}{to make known} τὸν πλοῦτον\FnParseFormGloss{b}{accusative masculine singular}{πλοῦτος}{riches} τῆς δόξης αὐτοῦ ἐπὶ σκεύη\FnParseFormGloss{c}{accusative neuter plural}{σκεῦος}{possession} ἐλέους,\FnParseFormGloss{d}{genitive neuter singular}{ἔλεος}{mercy} ἃ προητοίμασεν\FnParseFormGloss{e}{aorist active indicative 3rd singular}{προετοιμάζω}{to prepare in advance} εἰς δόξαν, 
\MainTextVerseMark{9}{24}οὓς καὶ ἐκάλεσεν\FnParse{a}{aorist active indicative 3rd singular} ἡμᾶς οὐ μόνον ἐξ Ἰουδαίων ἀλλὰ καὶ ἐξ ἐθνῶν;— 
\MainTextVerseMark{9}{25}ὡς καὶ ἐν τῷ Ὡσηὲ\FnParseFormGloss{a}{dative masculine singular}{Ὡσηέ}{Hosea} λέγει·\FnParse{b}{present active indicative 3rd singular} Καλέσω\FnParse{c}{future active indicative 1st singular} τὸν οὐ λαόν μου λαόν μου καὶ τὴν οὐκ ἠγαπημένην\FnParse{d}{perfect passive participle accusative feminine singular} ἠγαπημένην·\FnParse{e}{perfect passive participle accusative feminine singular} 
\MainTextVerseMark{9}{26}καὶ ἔσται\FnParse{a}{future middle indicative 3rd singular} ἐν τῷ τόπῳ οὗ\FnFormGloss{b}{οὗ}{where} ἐρρέθη\FnParse{c}{aorist passive indicative 3rd singular} ⸀αὐτοῖς· Οὐ λαός μου ὑμεῖς, ἐκεῖ κληθήσονται\FnParse{d}{future passive indicative 3rd plural} υἱοὶ θεοῦ ζῶντος.\FnParse{e}{present active participle genitive masculine singular} 
\MainTextVerseMark{9}{27}Ἠσαΐας\FnParseFormGloss{a}{nominative masculine singular}{Ἠσαΐας}{Isaiah} δὲ κράζει\FnParse{b}{present active indicative 3rd singular} ὑπὲρ τοῦ Ἰσραήλ· Ἐὰν ᾖ\FnParse{c}{present active subjunctive 3rd singular} ὁ ἀριθμὸς\FnParseFormGloss{d}{nominative masculine singular}{ἀριθμός}{number} τῶν υἱῶν Ἰσραὴλ ὡς ἡ ἄμμος\FnParseFormGloss{e}{nominative feminine singular}{ἄμμος}{sand} τῆς θαλάσσης, τὸ ⸀ὑπόλειμμα\FnParseFormGloss{f}{nominative neuter singular}{ὑπόλειμμα}{remnant} σωθήσεται·\FnParse{g}{future passive indicative 3rd singular} 
\MainTextVerseMark{9}{28}λόγον γὰρ συντελῶν\FnParseFormGloss{a}{present active participle nominative masculine singular}{συντελέω}{to finish} καὶ ⸀συντέμνων\FnParseFormGloss{b}{present active participle nominative masculine singular}{συντέμνω}{to cut short} ποιήσει\FnParse{c}{future active indicative 3rd singular} κύριος ἐπὶ τῆς γῆς. 
\MainTextVerseMark{9}{29}καὶ καθὼς προείρηκεν\FnParseFormGloss{a}{perfect active indicative 3rd singular}{προλέγω}{to tell beforehand} Ἠσαΐας·\FnParseFormGloss{b}{nominative masculine singular}{Ἠσαΐας}{Isaiah} Εἰ μὴ κύριος Σαβαὼθ\FnParseFormGloss{c}{genitive masculine plural}{Σαβαώθ}{Almighty} ἐγκατέλιπεν\FnParseFormGloss{d}{aorist active indicative 3rd singular}{ἐγκαταλείπω}{to forsake} ἡμῖν σπέρμα, ὡς Σόδομα\FnParseFormGloss{e}{nominative neuter plural}{Σόδομα}{Sodom} ἂν ἐγενήθημεν\FnParse{f}{aorist passive indicative 1st plural} καὶ ὡς Γόμορρα\FnParseFormGloss{g}{nominative feminine singular}{Γόμορρα}{Gomorrah} ἂν ὡμοιώθημεν.\FnParseFormGloss{h}{aorist passive indicative 1st plural}{ὁμοιόω}{to make like} 
\MainTextVerseMark{9}{30}Τί οὖν ἐροῦμεν;\FnParse{a}{future active indicative 1st plural} ὅτι ἔθνη τὰ μὴ διώκοντα\FnParse{b}{present active participle nominative neuter plural} δικαιοσύνην κατέλαβεν\FnParseFormGloss{c}{aorist active indicative 3rd singular}{καταλαμβάνω}{to obtain} δικαιοσύνην, δικαιοσύνην δὲ τὴν ἐκ πίστεως· 
\MainTextVerseMark{9}{31}Ἰσραὴλ δὲ διώκων\FnParse{a}{present active participle nominative masculine singular} νόμον δικαιοσύνης εἰς ⸀νόμον οὐκ ἔφθασεν.\FnParseFormGloss{b}{aorist active indicative 3rd singular}{φθάνω}{to precede} 
\MainTextVerseMark{9}{32}διὰ τί; ὅτι οὐκ ἐκ πίστεως ἀλλ’ ὡς ἐξ ⸀ἔργων· ⸀προσέκοψαν\FnParseFormGloss{a}{aorist active indicative 3rd plural}{προσκόπτω}{to strike} τῷ λίθῳ τοῦ προσκόμματος,\FnParseFormGloss{b}{genitive neuter singular}{πρόσκομμα}{stumbling block} 
\MainTextVerseMark{9}{33}καθὼς γέγραπται·\FnParse{a}{perfect passive indicative 3rd singular} Ἰδοὺ τίθημι\FnParse{b}{present active indicative 1st singular} ἐν Σιὼν\FnParseFormGloss{c}{dative feminine singular}{Σιών}{Zion} λίθον προσκόμματος\FnParseFormGloss{d}{genitive neuter singular}{πρόσκομμα}{stumbling block} καὶ πέτραν\FnParseFormGloss{e}{accusative feminine singular}{πέτρα}{rock} σκανδάλου,\FnParseFormGloss{f}{genitive neuter singular}{σκάνδαλον}{stumbling block} ⸀καὶ ὁ πιστεύων\FnParse{g}{present active participle nominative masculine singular} ἐπ’ αὐτῷ οὐ καταισχυνθήσεται.\FnParseFormGloss{h}{future passive indicative 3rd singular}{καταισχύνω}{to dishonor} 
\ChapterHeader{10}{Chapter 10}

\MainTextVerseMark{10}{1}Ἀδελφοί, ἡ μὲν εὐδοκία\FnParseFormGloss{a}{nominative feminine singular}{εὐδοκία}{goodwill} τῆς ἐμῆς καρδίας καὶ ἡ ⸀δέησις\FnParseFormGloss{b}{nominative feminine singular}{δέησις}{prayer} πρὸς τὸν θεὸν ὑπὲρ ⸀αὐτῶν εἰς σωτηρίαν. 
\MainTextVerseMark{10}{2}μαρτυρῶ\FnParse{a}{present active indicative 1st singular} γὰρ αὐτοῖς ὅτι ζῆλον\FnParseFormGloss{b}{accusative masculine singular}{ζῆλος}{zeal} θεοῦ ἔχουσιν·\FnParse{c}{present active indicative 3rd plural} ἀλλ’ οὐ κατ’ ἐπίγνωσιν,\FnParseFormGloss{d}{accusative feminine singular}{ἐπίγνωσις}{knowledge} 
\MainTextVerseMark{10}{3}ἀγνοοῦντες\FnParseFormGloss{a}{present active participle nominative masculine plural}{ἀγνοέω}{to be ignorant} γὰρ τὴν τοῦ θεοῦ δικαιοσύνην, καὶ τὴν ⸀ἰδίαν ζητοῦντες\FnParse{b}{present active participle nominative masculine plural} στῆσαι,\FnParse{c}{aorist active infinitive} τῇ δικαιοσύνῃ τοῦ θεοῦ οὐχ ὑπετάγησαν·\FnParse{d}{aorist passive indicative 3rd plural} 
\MainTextVerseMark{10}{4}τέλος γὰρ νόμου Χριστὸς εἰς δικαιοσύνην παντὶ τῷ πιστεύοντι.\FnParse{a}{present active participle dative masculine singular} 
\MainTextVerseMark{10}{5}Μωϋσῆς γὰρ γράφει\FnParse{a}{present active indicative 3rd singular} ⸂ὅτι τὴν δικαιοσύνην τὴν ἐκ ⸀τοῦ νόμου⸃ ὁ ⸀ποιήσας\FnParse{b}{aorist active participle nominative masculine singular} ἄνθρωπος ζήσεται\FnParse{c}{future middle indicative 3rd singular} ἐν ⸀αὐτῇ. 
\MainTextVerseMark{10}{6}ἡ δὲ ἐκ πίστεως δικαιοσύνη οὕτως λέγει·\FnParse{a}{present active indicative 3rd singular} Μὴ εἴπῃς\FnParse{b}{aorist active subjunctive 2nd singular} ἐν τῇ καρδίᾳ σου· Τίς ἀναβήσεται\FnParse{c}{future middle indicative 3rd singular} εἰς τὸν οὐρανόν; τοῦτ’ ἔστιν\FnParse{d}{present active indicative 3rd singular} Χριστὸν καταγαγεῖν·\FnParseFormGloss{e}{aorist active infinitive}{κατάγω}{to bring} 
\MainTextVerseMark{10}{7}ἤ· Τίς καταβήσεται\FnParse{a}{future middle indicative 3rd singular} εἰς τὴν ἄβυσσον;\FnParseFormGloss{b}{accusative feminine singular}{ἄβυσσος}{Abyss} τοῦτ’ ἔστιν\FnParse{c}{present active indicative 3rd singular} Χριστὸν ἐκ νεκρῶν ἀναγαγεῖν.\FnParseFormGloss{d}{aorist active infinitive}{ἀνάγω}{to lead up} 
\MainTextVerseMark{10}{8}ἀλλὰ τί λέγει;\FnParse{a}{present active indicative 3rd singular} Ἐγγύς σου τὸ ῥῆμά ἐστιν,\FnParse{b}{present active indicative 3rd singular} ἐν τῷ στόματί σου καὶ ἐν τῇ καρδίᾳ σου, τοῦτ’ ἔστιν\FnParse{c}{present active indicative 3rd singular} τὸ ῥῆμα τῆς πίστεως ὃ κηρύσσομεν.\FnParse{d}{present active indicative 1st plural} 
\MainTextVerseMark{10}{9}ὅτι ἐὰν ⸀ὁμολογήσῃς\FnParseFormGloss{a}{aorist active subjunctive 2nd singular}{ὁμολογέω}{to confess} ἐν τῷ στόματί ⸀σου κύριον Ἰησοῦν, καὶ πιστεύσῃς\FnParse{b}{aorist active subjunctive 2nd singular} ἐν τῇ καρδίᾳ σου ὅτι ὁ θεὸς αὐτὸν ἤγειρεν\FnParse{c}{aorist active indicative 3rd singular} ἐκ νεκρῶν, σωθήσῃ·\FnParse{d}{future passive indicative 2nd singular} 
\MainTextVerseMark{10}{10}καρδίᾳ γὰρ πιστεύεται\FnParse{a}{present passive indicative 3rd singular} εἰς δικαιοσύνην, στόματι δὲ ὁμολογεῖται\FnParseFormGloss{b}{present passive indicative 3rd singular}{ὁμολογέω}{to confess} εἰς σωτηρίαν· 
\MainTextVerseMark{10}{11}λέγει\FnParse{a}{present active indicative 3rd singular} γὰρ ἡ γραφή· Πᾶς ὁ πιστεύων\FnParse{b}{present active participle nominative masculine singular} ἐπ’ αὐτῷ οὐ καταισχυνθήσεται.\FnParseFormGloss{c}{future passive indicative 3rd singular}{καταισχύνω}{to dishonor} 
\MainTextVerseMark{10}{12}οὐ γάρ ἐστιν\FnParse{a}{present active indicative 3rd singular} διαστολὴ\FnParseFormGloss{b}{nominative feminine singular}{διαστολή}{difference} Ἰουδαίου τε καὶ Ἕλληνος,\FnParseFormGloss{c}{genitive masculine singular}{Ἕλλην}{Greek} ὁ γὰρ αὐτὸς κύριος πάντων, πλουτῶν\FnParseFormGloss{d}{present active participle nominative masculine singular}{πλουτέω}{to be rich} εἰς πάντας τοὺς ἐπικαλουμένους\FnParse{e}{present middle participle accusative masculine plural} αὐτόν· 
\MainTextVerseMark{10}{13}Πᾶς γὰρ ὃς ἂν ἐπικαλέσηται\FnParse{a}{aorist middle subjunctive 3rd singular} τὸ ὄνομα κυρίου σωθήσεται.\FnParse{b}{future passive indicative 3rd singular} 
\MainTextVerseMark{10}{14}Πῶς οὖν ⸀ἐπικαλέσωνται\FnParse{a}{aorist middle subjunctive 3rd plural} εἰς ὃν οὐκ ἐπίστευσαν;\FnParse{b}{aorist active indicative 3rd plural} πῶς δὲ ⸀πιστεύσωσιν\FnParse{c}{aorist active subjunctive 3rd plural} οὗ οὐκ ἤκουσαν;\FnParse{d}{aorist active indicative 3rd plural} πῶς δὲ ⸀ἀκούσωσιν\FnParse{e}{aorist active subjunctive 3rd plural} χωρὶς κηρύσσοντος;\FnParse{f}{present active participle genitive masculine singular} 
\MainTextVerseMark{10}{15}πῶς δὲ ⸀κηρύξωσιν\FnParse{a}{aorist active subjunctive 3rd plural} ἐὰν μὴ ἀποσταλῶσιν;\FnParse{b}{aorist passive subjunctive 3rd plural} ⸀καθὼς γέγραπται·\FnParse{c}{perfect passive indicative 3rd singular} Ὡς ὡραῖοι\FnParseFormGloss{d}{nominative masculine plural}{ὡραῖος}{beautiful} οἱ πόδες τῶν ⸀εὐαγγελιζομένων\FnParse{e}{present middle participle genitive masculine plural} ⸀τὰ ἀγαθά. 
\MainTextVerseMark{10}{16}ἀλλ’ οὐ πάντες ὑπήκουσαν\FnParseFormGloss{a}{aorist active indicative 3rd plural}{ὑπακούω}{to obey} τῷ εὐαγγελίῳ· Ἠσαΐας\FnParseFormGloss{b}{nominative masculine singular}{Ἠσαΐας}{Isaiah} γὰρ λέγει·\FnParse{c}{present active indicative 3rd singular} Κύριε, τίς ἐπίστευσεν\FnParse{d}{aorist active indicative 3rd singular} τῇ ἀκοῇ\FnParseFormGloss{e}{dative feminine singular}{ἀκοή}{hearing} ἡμῶν; 
\MainTextVerseMark{10}{17}ἄρα ἡ πίστις ἐξ ἀκοῆς,\FnParseFormGloss{a}{genitive feminine singular}{ἀκοή}{hearing} ἡ δὲ ἀκοὴ\FnParseFormGloss{b}{nominative feminine singular}{ἀκοή}{hearing} διὰ ῥήματος ⸀Χριστοῦ. 
\MainTextVerseMark{10}{18}Ἀλλὰ λέγω,\FnParse{a}{present active indicative 1st singular} μὴ οὐκ ἤκουσαν;\FnParse{b}{aorist active indicative 3rd plural} μενοῦνγε·\FnFormGloss{c}{μενοῦνγε}{rather} Εἰς πᾶσαν τὴν γῆν ἐξῆλθεν\FnParse{d}{aorist active indicative 3rd singular} ὁ φθόγγος\FnParseFormGloss{e}{nominative masculine singular}{φθόγγος}{voice} αὐτῶν, καὶ εἰς τὰ πέρατα\FnParseFormGloss{f}{accusative neuter plural}{πέρας}{end} τῆς οἰκουμένης\FnParseFormGloss{g}{genitive feminine singular}{οἰκουμένη}{the world} τὰ ῥήματα αὐτῶν. 
\MainTextVerseMark{10}{19}ἀλλὰ λέγω,\FnParse{a}{present active indicative 1st singular} μὴ ⸂Ἰσραὴλ οὐκ ἔγνω⸃;\FnParse{b}{aorist active indicative 3rd singular} πρῶτος Μωϋσῆς λέγει·\FnParse{c}{present active indicative 3rd singular} Ἐγὼ παραζηλώσω\FnParseFormGloss{d}{future active indicative 1st singular}{παραζηλόω}{to make envious} ὑμᾶς ἐπ’ οὐκ ἔθνει, ἐπ’ ἔθνει ἀσυνέτῳ\FnParseFormGloss{e}{dative neuter singular}{ἀσύνετος}{senseless} παροργιῶ\FnParseFormGloss{f}{future active indicative 1st singular}{παροργίζω}{to anger} ὑμᾶς. 
\MainTextVerseMark{10}{20}Ἠσαΐας\FnParseFormGloss{a}{nominative masculine singular}{Ἠσαΐας}{Isaiah} δὲ ἀποτολμᾷ\FnParseFormGloss{b}{present active indicative 3rd singular}{ἀποτολμάω}{to bring forth boldly} καὶ λέγει·\FnParse{c}{present active indicative 3rd singular} Εὑρέθην\FnParse{d}{aorist passive indicative 1st singular} ⸀ἐν τοῖς ἐμὲ μὴ ζητοῦσιν,\FnParse{e}{present active participle dative masculine plural} ἐμφανὴς\FnParseFormGloss{f}{nominative masculine singular}{ἐμφανής}{seen} ⸀ἐγενόμην\FnParse{g}{aorist middle indicative 1st singular} τοῖς ἐμὲ μὴ ἐπερωτῶσιν.\FnParse{h}{present active participle dative masculine plural} 
\MainTextVerseMark{10}{21}πρὸς δὲ τὸν Ἰσραὴλ λέγει·\FnParse{a}{present active indicative 3rd singular} Ὅλην τὴν ἡμέραν ἐξεπέτασα\FnParseFormGloss{b}{aorist active indicative 1st singular}{ἐκπετάννυμι}{to hold out} τὰς χεῖράς μου πρὸς λαὸν ἀπειθοῦντα\FnParseFormGloss{c}{present active participle accusative masculine singular}{ἀπειθέω}{to disobey} καὶ ἀντιλέγοντα.\FnParseFormGloss{d}{present active participle accusative masculine singular}{ἀντιλέγω}{to speak against} 
\ChapterHeader{11}{Chapter 11}

\MainTextVerseMark{11}{1}Λέγω\FnParse{a}{present active indicative 1st singular} οὖν, μὴ ἀπώσατο\FnParseFormGloss{b}{aorist middle indicative 3rd singular}{ἀπωθέομαι}{to reject} ὁ θεὸς τὸν λαὸν αὐτοῦ; μὴ γένοιτο·\FnParse{c}{aorist middle optative 3rd singular} καὶ γὰρ ἐγὼ Ἰσραηλίτης\FnParseFormGloss{d}{nominative masculine singular}{Ἰσραηλίτης}{Israelite} εἰμί,\FnParse{e}{present active indicative 1st singular} ἐκ σπέρματος Ἀβραάμ, φυλῆς Βενιαμίν.\FnParseFormGloss{f}{genitive masculine singular}{Βενιαμίν}{Benjamin} 
\MainTextVerseMark{11}{2}οὐκ ἀπώσατο\FnParseFormGloss{a}{aorist middle indicative 3rd singular}{ἀπωθέομαι}{to reject} ὁ θεὸς τὸν λαὸν αὐτοῦ ὃν προέγνω.\FnParseFormGloss{b}{aorist active indicative 3rd singular}{προγινώσκω}{to know beforehand} ἢ οὐκ οἴδατε\FnParse{c}{perfect active indicative 2nd plural} ἐν Ἠλίᾳ\FnParseFormGloss{d}{dative masculine singular}{Ἠλίας}{Elijah} τί λέγει\FnParse{e}{present active indicative 3rd singular} ἡ γραφή, ὡς ἐντυγχάνει\FnParseFormGloss{f}{present active indicative 3rd singular}{ἐντυγχάνω}{to intercede} τῷ θεῷ κατὰ τοῦ ⸀Ἰσραήλ; 
\MainTextVerseMark{11}{3}Κύριε, τοὺς προφήτας σου ἀπέκτειναν,\FnParse{a}{aorist active indicative 3rd plural} ⸀τὰ θυσιαστήριά\FnParseFormGloss{b}{accusative neuter plural}{θυσιαστήριον}{altar} σου κατέσκαψαν,\FnParseFormGloss{c}{aorist active indicative 3rd plural}{κατασκάπτω}{to tear down} κἀγὼ ὑπελείφθην\FnParseFormGloss{d}{aorist passive indicative 1st singular}{ὑπολείπω}{to be left} μόνος, καὶ ζητοῦσιν\FnParse{e}{present active indicative 3rd plural} τὴν ψυχήν μου. 
\MainTextVerseMark{11}{4}ἀλλὰ τί λέγει\FnParse{a}{present active indicative 3rd singular} αὐτῷ ὁ χρηματισμός;\FnParseFormGloss{b}{nominative masculine singular}{χρηματισμός}{a response from God} Κατέλιπον\FnParseFormGloss{c}{aorist active indicative 1st singular}{καταλείπω}{to leave} ἐμαυτῷ ἑπτακισχιλίους\FnParseFormGloss{d}{accusative masculine plural}{ἑπτακισχίλιοι}{seven thousand} ἄνδρας, οἵτινες οὐκ ἔκαμψαν\FnParseFormGloss{e}{aorist active indicative 3rd plural}{κάμπτω}{to bend} γόνυ\FnParseFormGloss{f}{accusative neuter singular}{γόνυ}{knee} τῇ Βάαλ.\FnParseFormGloss{g}{genitive masculine singular}{Βάαλ}{Baal} 
\MainTextVerseMark{11}{5}οὕτως οὖν καὶ ἐν τῷ νῦν καιρῷ λεῖμμα\FnParseFormGloss{a}{nominative neuter singular}{λεῖμμα}{remnant} κατ’ ἐκλογὴν\FnParseFormGloss{b}{accusative feminine singular}{ἐκλογή}{election} χάριτος γέγονεν·\FnParse{c}{perfect active indicative 3rd singular} 
\MainTextVerseMark{11}{6}εἰ δὲ χάριτι, οὐκέτι ἐξ ἔργων, ἐπεὶ\FnFormGloss{a}{ἐπεί}{since} ἡ χάρις οὐκέτι γίνεται\FnParse{b}{present middle indicative 3rd singular} ⸀χάρις. 
\MainTextVerseMark{11}{7}τί οὖν; ὃ ἐπιζητεῖ\FnParseFormGloss{a}{present active indicative 3rd singular}{ἐπιζητέω}{to look for} Ἰσραήλ, τοῦτο οὐκ ἐπέτυχεν,\FnParseFormGloss{b}{aorist active indicative 3rd singular}{ἐπιτυγχάνω}{to obtain} ἡ δὲ ἐκλογὴ\FnParseFormGloss{c}{nominative feminine singular}{ἐκλογή}{election} ἐπέτυχεν·\FnParseFormGloss{d}{aorist active indicative 3rd singular}{ἐπιτυγχάνω}{to obtain} οἱ δὲ λοιποὶ ἐπωρώθησαν,\FnParseFormGloss{e}{aorist passive indicative 3rd plural}{πωρόω}{to harden} 
\MainTextVerseMark{11}{8}⸀καθὼς γέγραπται·\FnParse{a}{perfect passive indicative 3rd singular} Ἔδωκεν\FnParse{b}{aorist active indicative 3rd singular} αὐτοῖς ὁ θεὸς πνεῦμα κατανύξεως,\FnParseFormGloss{c}{genitive feminine singular}{κατάνυξις}{stupor} ὀφθαλμοὺς τοῦ μὴ βλέπειν\FnParse{d}{present active infinitive} καὶ ὦτα τοῦ μὴ ἀκούειν,\FnParse{e}{present active infinitive} ἕως τῆς σήμερον ἡμέρας. 
\MainTextVerseMark{11}{9}καὶ Δαυὶδ λέγει·\FnParse{a}{present active indicative 3rd singular} Γενηθήτω\FnParse{b}{aorist passive imperative 3rd singular} ἡ τράπεζα\FnParseFormGloss{c}{nominative feminine singular}{τράπεζα}{table} αὐτῶν εἰς παγίδα\FnParseFormGloss{d}{accusative feminine singular}{παγίς}{trap} καὶ εἰς θήραν\FnParseFormGloss{e}{accusative feminine singular}{θήρα}{trap} καὶ εἰς σκάνδαλον\FnParseFormGloss{f}{accusative neuter singular}{σκάνδαλον}{stumbling block} καὶ εἰς ἀνταπόδομα\FnParseFormGloss{g}{accusative neuter singular}{ἀνταπόδομα}{repayment} αὐτοῖς, 
\MainTextVerseMark{11}{10}σκοτισθήτωσαν\FnParseFormGloss{a}{aorist passive imperative 3rd plural}{σκοτίζομαι}{to be dark} οἱ ὀφθαλμοὶ αὐτῶν τοῦ μὴ βλέπειν,\FnParse{b}{present active infinitive} καὶ τὸν νῶτον\FnParseFormGloss{c}{accusative masculine singular}{νῶτος}{back} αὐτῶν διὰ παντὸς σύγκαμψον.\FnParseFormGloss{d}{aorist active imperative 2nd singular}{συγκάμπτω}{to be bent over} 
\MainTextVerseMark{11}{11}Λέγω\FnParse{a}{present active indicative 1st singular} οὖν, μὴ ἔπταισαν\FnParseFormGloss{b}{aorist active indicative 3rd plural}{πταίω}{to stumble} ἵνα πέσωσιν;\FnParse{c}{aorist active subjunctive 3rd plural} μὴ γένοιτο·\FnParse{d}{aorist middle optative 3rd singular} ἀλλὰ τῷ αὐτῶν παραπτώματι\FnParseFormGloss{e}{dative neuter singular}{παράπτωμα}{trespass} ἡ σωτηρία τοῖς ἔθνεσιν, εἰς τὸ παραζηλῶσαι\FnParseFormGloss{f}{aorist active infinitive}{παραζηλόω}{to make envious} αὐτούς. 
\MainTextVerseMark{11}{12}εἰ δὲ τὸ παράπτωμα\FnParseFormGloss{a}{nominative neuter singular}{παράπτωμα}{trespass} αὐτῶν πλοῦτος\FnParseFormGloss{b}{nominative masculine singular}{πλοῦτος}{riches} κόσμου καὶ τὸ ἥττημα\FnParseFormGloss{c}{nominative neuter singular}{ἥττημα}{loss} αὐτῶν πλοῦτος\FnParseFormGloss{d}{nominative masculine singular}{πλοῦτος}{riches} ἐθνῶν, πόσῳ\FnParseFormGloss{e}{dative neuter singular}{πόσος}{how great?} μᾶλλον τὸ πλήρωμα\FnParseFormGloss{f}{nominative neuter singular}{πλήρωμα}{fullness} αὐτῶν. 
\MainTextVerseMark{11}{13}Ὑμῖν ⸀δὲ λέγω\FnParse{a}{present active indicative 1st singular} τοῖς ἔθνεσιν. ἐφ’ ὅσον μὲν ⸀οὖν εἰμι\FnParse{b}{present active indicative 1st singular} ἐγὼ ἐθνῶν ἀπόστολος, τὴν διακονίαν μου δοξάζω,\FnParse{c}{present active indicative 1st singular} 
\MainTextVerseMark{11}{14}εἴ πως\FnFormGloss{a}{πώς}{in any way} παραζηλώσω\FnParseFormGloss{b}{future active indicative 1st singular}{παραζηλόω}{to make envious} μου τὴν σάρκα καὶ σώσω\FnParse{c}{future active indicative 1st singular} τινὰς ἐξ αὐτῶν. 
\MainTextVerseMark{11}{15}εἰ γὰρ ἡ ἀποβολὴ\FnParseFormGloss{a}{nominative feminine singular}{ἀποβολή}{rejection} αὐτῶν καταλλαγὴ\FnParseFormGloss{b}{nominative feminine singular}{καταλλαγή}{reconciliation} κόσμου, τίς ἡ πρόσλημψις\FnParseFormGloss{c}{nominative feminine singular}{πρόσλημψις}{acceptance} εἰ μὴ ζωὴ ἐκ νεκρῶν; 
\MainTextVerseMark{11}{16}εἰ δὲ ἡ ἀπαρχὴ\FnParseFormGloss{a}{nominative feminine singular}{ἀπαρχή}{firstfruits} ἁγία, καὶ τὸ φύραμα·\FnParseFormGloss{b}{nominative neuter singular}{φύραμα}{lump} καὶ εἰ ἡ ῥίζα\FnParseFormGloss{c}{nominative feminine singular}{ῥίζα}{root} ἁγία, καὶ οἱ κλάδοι.\FnParseFormGloss{d}{nominative masculine plural}{κλάδος}{branch} 
\MainTextVerseMark{11}{17}Εἰ δέ τινες τῶν κλάδων\FnParseFormGloss{a}{genitive masculine plural}{κλάδος}{branch} ἐξεκλάσθησαν,\FnParseFormGloss{b}{aorist passive indicative 3rd plural}{ἐκκλάω}{to be broken off} σὺ δὲ ἀγριέλαιος\FnParseFormGloss{c}{nominative feminine singular}{ἀγριέλαιος}{wild olive tree} ὢν\FnParse{d}{present active participle nominative masculine singular} ἐνεκεντρίσθης\FnParseFormGloss{e}{aorist passive indicative 2nd singular}{ἐγκεντρίζω}{to graft into} ἐν αὐτοῖς καὶ συγκοινωνὸς\FnParseFormGloss{f}{nominative masculine singular}{συγκοινωνός}{sharer} τῆς ⸀ῥίζης\FnParseFormGloss{g}{genitive feminine singular}{ῥίζα}{root} τῆς πιότητος\FnParseFormGloss{h}{genitive feminine singular}{πιότης}{richness} τῆς ἐλαίας\FnParseFormGloss{i}{genitive feminine singular}{ἐλαία}{olive} ἐγένου,\FnParse{j}{aorist middle indicative 2nd singular} 
\MainTextVerseMark{11}{18}μὴ κατακαυχῶ\FnParseFormGloss{a}{present middle imperative 2nd singular}{κατακαυχάομαι}{to boast about} τῶν κλάδων·\FnParseFormGloss{b}{genitive masculine plural}{κλάδος}{branch} εἰ δὲ κατακαυχᾶσαι,\FnParseFormGloss{c}{present middle indicative 2nd singular}{κατακαυχάομαι}{to boast about} οὐ σὺ τὴν ῥίζαν\FnParseFormGloss{d}{accusative feminine singular}{ῥίζα}{root} βαστάζεις\FnParseFormGloss{e}{present active indicative 2nd singular}{βαστάζω}{to carry} ἀλλὰ ἡ ῥίζα\FnParseFormGloss{f}{nominative feminine singular}{ῥίζα}{root} σέ. 
\MainTextVerseMark{11}{19}ἐρεῖς\FnParse{a}{future active indicative 2nd singular} οὖν· Ἐξεκλάσθησαν\FnParseFormGloss{b}{aorist passive indicative 3rd plural}{ἐκκλάω}{to be broken off} κλάδοι\FnParseFormGloss{c}{nominative masculine plural}{κλάδος}{branch} ἵνα ἐγὼ ἐγκεντρισθῶ.\FnParseFormGloss{d}{aorist passive subjunctive 1st singular}{ἐγκεντρίζω}{to graft into} 
\MainTextVerseMark{11}{20}καλῶς· τῇ ἀπιστίᾳ\FnParseFormGloss{a}{dative feminine singular}{ἀπιστία}{unbelief} ἐξεκλάσθησαν,\FnParseFormGloss{b}{aorist passive indicative 3rd plural}{ἐκκλάω}{to be broken off} σὺ δὲ τῇ πίστει ἕστηκας.\FnParse{c}{perfect active indicative 2nd singular} μὴ ⸂ὑψηλὰ\FnParseFormGloss{d}{accusative neuter plural}{ὑψηλός}{high} φρόνει⸃,\FnParseFormGloss{e}{present active imperative 2nd singular}{φρονέω}{to think} ἀλλὰ φοβοῦ·\FnParse{f}{present middle imperative 2nd singular} 
\MainTextVerseMark{11}{21}εἰ γὰρ ὁ θεὸς τῶν κατὰ φύσιν\FnParseFormGloss{a}{accusative feminine singular}{φύσις}{nature} κλάδων\FnParseFormGloss{b}{genitive masculine plural}{κλάδος}{branch} οὐκ ἐφείσατο,\FnParseFormGloss{c}{aorist middle indicative 3rd singular}{φείδομαι}{to spare} ⸀οὐδὲ σοῦ φείσεται.\FnParseFormGloss{d}{future middle indicative 3rd singular}{φείδομαι}{to spare} 
\MainTextVerseMark{11}{22}ἴδε οὖν χρηστότητα\FnParseFormGloss{a}{accusative feminine singular}{χρηστότης}{kindness} καὶ ἀποτομίαν\FnParseFormGloss{b}{accusative feminine singular}{ἀποτομία}{sternness} θεοῦ· ἐπὶ μὲν τοὺς πεσόντας\FnParse{c}{aorist active participle accusative masculine plural} ⸀ἀποτομία,\FnParseFormGloss{d}{nominative feminine singular}{ἀποτομία}{sternness} ἐπὶ δὲ σὲ ⸂χρηστότης\FnParseFormGloss{e}{nominative feminine singular}{χρηστότης}{kindness} θεοῦ⸃, ἐὰν ⸀ἐπιμένῃς\FnParseFormGloss{f}{present active subjunctive 2nd singular}{ἐπιμένω}{to stay} τῇ χρηστότητι,\FnParseFormGloss{g}{dative feminine singular}{χρηστότης}{kindness} ἐπεὶ\FnFormGloss{h}{ἐπεί}{since} καὶ σὺ ἐκκοπήσῃ.\FnParseFormGloss{i}{future passive indicative 2nd singular}{ἐκκόπτω}{to cut off} 
\MainTextVerseMark{11}{23}κἀκεῖνοι\FnParseFormGloss{a}{nominative masculine plural}{κἀκεῖνος}{and that one} δέ, ἐὰν μὴ ⸀ἐπιμένωσι\FnParseFormGloss{b}{present active subjunctive 3rd plural}{ἐπιμένω}{to stay} τῇ ἀπιστίᾳ,\FnParseFormGloss{c}{dative feminine singular}{ἀπιστία}{unbelief} ἐγκεντρισθήσονται·\FnParseFormGloss{d}{future passive indicative 3rd plural}{ἐγκεντρίζω}{to graft into} δυνατὸς γάρ ⸂ἐστιν\FnParse{e}{present active indicative 3rd singular} ὁ θεὸς⸃ πάλιν ἐγκεντρίσαι\FnParseFormGloss{f}{aorist active infinitive}{ἐγκεντρίζω}{to graft into} αὐτούς. 
\MainTextVerseMark{11}{24}εἰ γὰρ σὺ ἐκ τῆς κατὰ φύσιν\FnParseFormGloss{a}{accusative feminine singular}{φύσις}{nature} ἐξεκόπης\FnParseFormGloss{b}{aorist passive indicative 2nd singular}{ἐκκόπτω}{to cut off} ἀγριελαίου\FnParseFormGloss{c}{genitive feminine singular}{ἀγριέλαιος}{wild olive tree} καὶ παρὰ φύσιν\FnParseFormGloss{d}{accusative feminine singular}{φύσις}{nature} ἐνεκεντρίσθης\FnParseFormGloss{e}{aorist passive indicative 2nd singular}{ἐγκεντρίζω}{to graft into} εἰς καλλιέλαιον,\FnParseFormGloss{f}{accusative feminine singular}{καλλιέλαιος}{cultivated olive tree} πόσῳ\FnParseFormGloss{g}{dative neuter singular}{πόσος}{how great?} μᾶλλον οὗτοι οἱ κατὰ φύσιν\FnParseFormGloss{h}{accusative feminine singular}{φύσις}{nature} ἐγκεντρισθήσονται\FnParseFormGloss{i}{future passive indicative 3rd plural}{ἐγκεντρίζω}{to graft into} τῇ ἰδίᾳ ἐλαίᾳ.\FnParseFormGloss{j}{dative feminine singular}{ἐλαία}{olive} 
\MainTextVerseMark{11}{25}Οὐ γὰρ θέλω\FnParse{a}{present active indicative 1st singular} ὑμᾶς ἀγνοεῖν,\FnParseFormGloss{b}{present active infinitive}{ἀγνοέω}{to be ignorant} ἀδελφοί, τὸ μυστήριον\FnParseFormGloss{c}{accusative neuter singular}{μυστήριον}{mystery} τοῦτο, ἵνα μὴ ⸀ἦτε\FnParse{d}{present active subjunctive 2nd plural} ἑαυτοῖς φρόνιμοι,\FnParseFormGloss{e}{nominative masculine plural}{φρόνιμος}{wise} ὅτι πώρωσις\FnParseFormGloss{f}{nominative feminine singular}{πώρωσις}{hardening} ἀπὸ μέρους τῷ Ἰσραὴλ γέγονεν\FnParse{g}{perfect active indicative 3rd singular} ἄχρι οὗ τὸ πλήρωμα\FnParseFormGloss{h}{nominative neuter singular}{πλήρωμα}{fullness} τῶν ἐθνῶν εἰσέλθῃ,\FnParse{i}{aorist active subjunctive 3rd singular} 
\MainTextVerseMark{11}{26}καὶ οὕτως πᾶς Ἰσραὴλ σωθήσεται·\FnParse{a}{future passive indicative 3rd singular} καθὼς γέγραπται·\FnParse{b}{perfect passive indicative 3rd singular} Ἥξει\FnParseFormGloss{c}{future active indicative 3rd singular}{ἥκω}{to come} ἐκ Σιὼν\FnParseFormGloss{d}{genitive feminine singular}{Σιών}{Zion} ὁ ⸀ῥυόμενος,\FnParseFormGloss{e}{present middle participle nominative masculine singular}{ῥύομαι}{to rescue} ἀποστρέψει\FnParseFormGloss{f}{future active indicative 3rd singular}{ἀποστρέφω}{to turn away from} ἀσεβείας\FnParseFormGloss{g}{accusative feminine plural}{ἀσέβεια}{ungodliness} ἀπὸ Ἰακώβ.\FnParseFormGloss{h}{genitive masculine singular}{Ἰακώβ}{Jacob} 
\MainTextVerseMark{11}{27}καὶ αὕτη αὐτοῖς ἡ παρ’ ἐμοῦ διαθήκη, ὅταν ἀφέλωμαι\FnParseFormGloss{a}{aorist middle subjunctive 1st singular}{ἀφαιρέω}{to take away from} τὰς ἁμαρτίας αὐτῶν. 
\MainTextVerseMark{11}{28}κατὰ μὲν τὸ εὐαγγέλιον ἐχθροὶ δι’ ὑμᾶς, κατὰ δὲ τὴν ἐκλογὴν\FnParseFormGloss{a}{accusative feminine singular}{ἐκλογή}{election} ἀγαπητοὶ διὰ τοὺς πατέρας· 
\MainTextVerseMark{11}{29}ἀμεταμέλητα\FnParseFormGloss{a}{nominative neuter plural}{ἀμεταμέλητος}{without regret} γὰρ τὰ χαρίσματα\FnParseFormGloss{b}{nominative neuter plural}{χάρισμα}{gracious gift} καὶ ἡ κλῆσις\FnParseFormGloss{c}{nominative feminine singular}{κλῆσις}{call} τοῦ θεοῦ. 
\MainTextVerseMark{11}{30}ὥσπερ ⸀γὰρ ὑμεῖς ποτε\FnFormGloss{a}{ποτέ}{once} ἠπειθήσατε\FnParseFormGloss{b}{aorist active indicative 2nd plural}{ἀπειθέω}{to disobey} τῷ θεῷ, νῦν δὲ ἠλεήθητε\FnParse{c}{aorist passive indicative 2nd plural} τῇ τούτων ἀπειθείᾳ,\FnParseFormGloss{d}{dative feminine singular}{ἀπείθεια}{disobedience} 
\MainTextVerseMark{11}{31}οὕτως καὶ οὗτοι νῦν ἠπείθησαν\FnParseFormGloss{a}{aorist active indicative 3rd plural}{ἀπειθέω}{to disobey} τῷ ὑμετέρῳ\FnParseFormGloss{b}{dative neuter singular}{ὑμέτερος}{your} ἐλέει\FnParseFormGloss{c}{dative neuter singular}{ἔλεος}{mercy} ἵνα καὶ αὐτοὶ ⸀νῦν ἐλεηθῶσιν·\FnParse{d}{aorist passive subjunctive 3rd plural} 
\MainTextVerseMark{11}{32}συνέκλεισεν\FnParseFormGloss{a}{aorist active indicative 3rd singular}{συγκλείω}{to catch} γὰρ ὁ θεὸς τοὺς πάντας εἰς ἀπείθειαν\FnParseFormGloss{b}{accusative feminine singular}{ἀπείθεια}{disobedience} ἵνα τοὺς πάντας ἐλεήσῃ.\FnParse{c}{aorist active subjunctive 3rd singular} 
\MainTextVerseMark{11}{33}Ὦ\FnFormGloss{a}{ὦ}{O!} βάθος\FnParseFormGloss{b}{nominative neuter singular}{βάθος}{depth} πλούτου\FnParseFormGloss{c}{genitive masculine singular}{πλοῦτος}{riches} καὶ σοφίας καὶ γνώσεως\FnParseFormGloss{d}{genitive feminine singular}{γνῶσις}{knowledge} θεοῦ· ὡς ἀνεξεραύνητα\FnParseFormGloss{e}{nominative neuter plural}{ἀνεξεραύνητος}{unsearchable} τὰ κρίματα\FnParseFormGloss{f}{nominative neuter plural}{κρίμα}{judgment} αὐτοῦ καὶ ἀνεξιχνίαστοι\FnParseFormGloss{g}{nominative feminine plural}{ἀνεξιχνίαστος}{unsearchable} αἱ ὁδοὶ αὐτοῦ. 
\MainTextVerseMark{11}{34}Τίς γὰρ ἔγνω\FnParse{a}{aorist active indicative 3rd singular} νοῦν\FnParseFormGloss{b}{accusative masculine singular}{νοῦς}{mind} κυρίου; ἢ τίς σύμβουλος\FnParseFormGloss{c}{nominative masculine singular}{σύμβουλος}{counselor} αὐτοῦ ἐγένετο;\FnParse{d}{aorist middle indicative 3rd singular} 
\MainTextVerseMark{11}{35}ἢ τίς προέδωκεν\FnParseFormGloss{a}{aorist active indicative 3rd singular}{προδίδωμι}{to give beforehand} αὐτῷ, καὶ ἀνταποδοθήσεται\FnParseFormGloss{b}{future passive indicative 3rd singular}{ἀνταποδίδωμι}{to repay} αὐτῷ; 
\MainTextVerseMark{11}{36}ὅτι ἐξ αὐτοῦ καὶ δι’ αὐτοῦ καὶ εἰς αὐτὸν τὰ πάντα· αὐτῷ ἡ δόξα εἰς τοὺς αἰῶνας, ἀμήν. 
\ChapterHeader{12}{Chapter 12}

\MainTextVerseMark{12}{1}Παρακαλῶ\FnParse{a}{present active indicative 1st singular} οὖν ὑμᾶς, ἀδελφοί, διὰ τῶν οἰκτιρμῶν\FnParseFormGloss{b}{genitive masculine plural}{οἰκτιρμός}{compassion} τοῦ θεοῦ παραστῆσαι\FnParse{c}{aorist active infinitive} τὰ σώματα ὑμῶν θυσίαν\FnParseFormGloss{d}{accusative feminine singular}{θυσία}{sacrifice} ζῶσαν\FnParse{e}{present active participle accusative feminine singular} ἁγίαν ⸂εὐάρεστον\FnParseFormGloss{f}{accusative feminine singular}{εὐάρεστος}{pleasing} τῷ θεῷ⸃, τὴν λογικὴν\FnParseFormGloss{g}{accusative feminine singular}{λογικός}{spiritual} λατρείαν\FnParseFormGloss{h}{accusative feminine singular}{λατρεία}{worship} ὑμῶν· 
\MainTextVerseMark{12}{2}καὶ μὴ ⸀συσχηματίζεσθε\FnParseFormGloss{a}{present middle imperative 2nd plural}{συσχηματίζομαι}{to conform to a pattern} τῷ αἰῶνι τούτῳ, ἀλλὰ ⸀μεταμορφοῦσθε\FnParseFormGloss{b}{present passive imperative 2nd plural}{μεταμορφόομαι}{to be transformed} τῇ ἀνακαινώσει\FnParseFormGloss{c}{dative feminine singular}{ἀνακαίνωσις}{renewal} τοῦ ⸀νοός,\FnParseFormGloss{d}{genitive masculine singular}{νοῦς}{mind} εἰς τὸ δοκιμάζειν\FnParseFormGloss{e}{present active infinitive}{δοκιμάζω}{to test} ὑμᾶς τί τὸ θέλημα τοῦ θεοῦ, τὸ ἀγαθὸν καὶ εὐάρεστον\FnParseFormGloss{f}{nominative neuter singular}{εὐάρεστος}{pleasing} καὶ τέλειον.\FnParseFormGloss{g}{nominative neuter singular}{τέλειος}{perfect} 
\MainTextVerseMark{12}{3}Λέγω\FnParse{a}{present active indicative 1st singular} γὰρ διὰ τῆς χάριτος τῆς δοθείσης\FnParse{b}{aorist passive participle genitive feminine singular} μοι παντὶ τῷ ὄντι\FnParse{c}{present active participle dative masculine singular} ἐν ὑμῖν μὴ ὑπερφρονεῖν\FnParseFormGloss{d}{present active infinitive}{ὑπερφρονέω}{to have lofty thoughts} παρ’ ὃ δεῖ\FnParse{e}{present active indicative 3rd singular} φρονεῖν,\FnParseFormGloss{f}{present active infinitive}{φρονέω}{to think} ἀλλὰ φρονεῖν\FnParseFormGloss{g}{present active infinitive}{φρονέω}{to think} εἰς τὸ σωφρονεῖν,\FnParseFormGloss{h}{present active infinitive}{σωφρονέω}{to be of a sound mind} ἑκάστῳ ὡς ὁ θεὸς ἐμέρισεν\FnParseFormGloss{i}{aorist active indicative 3rd singular}{μερίζω}{to give} μέτρον\FnParseFormGloss{j}{accusative neuter singular}{μέτρον}{measure} πίστεως. 
\MainTextVerseMark{12}{4}καθάπερ\FnFormGloss{a}{καθάπερ}{as} γὰρ ἐν ἑνὶ σώματι ⸂πολλὰ μέλη⸃ ἔχομεν,\FnParse{b}{present active indicative 1st plural} τὰ δὲ μέλη πάντα οὐ τὴν αὐτὴν ἔχει\FnParse{c}{present active indicative 3rd singular} πρᾶξιν,\FnParseFormGloss{d}{accusative feminine singular}{πρᾶξις}{deed} 
\MainTextVerseMark{12}{5}οὕτως οἱ πολλοὶ ἓν σῶμά ἐσμεν\FnParse{a}{present active indicative 1st plural} ἐν Χριστῷ, ⸀τὸ δὲ καθ’ εἷς ἀλλήλων μέλη. 
\MainTextVerseMark{12}{6}ἔχοντες\FnParse{a}{present active participle nominative masculine plural} δὲ χαρίσματα\FnParseFormGloss{b}{accusative neuter plural}{χάρισμα}{gracious gift} κατὰ τὴν χάριν τὴν δοθεῖσαν\FnParse{c}{aorist passive participle accusative feminine singular} ἡμῖν διάφορα,\FnParseFormGloss{d}{accusative neuter plural}{διάφορος}{different} εἴτε προφητείαν\FnParseFormGloss{e}{accusative feminine singular}{προφητεία}{prophecy} κατὰ τὴν ἀναλογίαν\FnParseFormGloss{f}{accusative feminine singular}{ἀναλογία}{proportion} τῆς πίστεως, 
\MainTextVerseMark{12}{7}εἴτε διακονίαν ἐν τῇ διακονίᾳ, εἴτε ὁ διδάσκων\FnParse{a}{present active participle nominative masculine singular} ἐν τῇ διδασκαλίᾳ,\FnParseFormGloss{b}{dative feminine singular}{διδασκαλία}{teaching} 
\MainTextVerseMark{12}{8}εἴτε ὁ παρακαλῶν\FnParse{a}{present active participle nominative masculine singular} ἐν τῇ παρακλήσει,\FnParseFormGloss{b}{dative feminine singular}{παράκλησις}{encouragement} ὁ μεταδιδοὺς\FnParseFormGloss{c}{present active participle nominative masculine singular}{μεταδίδωμι}{to impart} ἐν ἁπλότητι,\FnParseFormGloss{d}{dative feminine singular}{ἁπλότης}{generosity; sincerity} ὁ προϊστάμενος\FnParseFormGloss{e}{present middle participle nominative masculine singular}{προΐστημι}{to manage} ἐν σπουδῇ,\FnParseFormGloss{f}{dative feminine singular}{σπουδή}{hurry} ὁ ἐλεῶν\FnParse{g}{present active participle nominative masculine singular} ἐν ἱλαρότητι.\FnParseFormGloss{h}{dative feminine singular}{ἱλαρότης}{cheerfully} 
\MainTextVerseMark{12}{9}Ἡ ἀγάπη ἀνυπόκριτος.\FnParseFormGloss{a}{nominative feminine singular}{ἀνυπόκριτος}{sincere} ἀποστυγοῦντες\FnParseFormGloss{b}{present active participle nominative masculine plural}{ἀποστυγέω}{to hate} τὸ πονηρόν, κολλώμενοι\FnParseFormGloss{c}{present passive participle nominative masculine plural}{κολλάομαι}{to join} τῷ ἀγαθῷ· 
\MainTextVerseMark{12}{10}τῇ φιλαδελφίᾳ\FnParseFormGloss{a}{dative feminine singular}{φιλαδελφία}{brotherly love} εἰς ἀλλήλους φιλόστοργοι,\FnParseFormGloss{b}{nominative masculine plural}{φιλόστοργος}{devoted} τῇ τιμῇ ἀλλήλους προηγούμενοι,\FnParseFormGloss{c}{present middle participle nominative masculine plural}{προηγέομαι}{to put above} 
\MainTextVerseMark{12}{11}τῇ σπουδῇ\FnParseFormGloss{a}{dative feminine singular}{σπουδή}{hurry} μὴ ὀκνηροί,\FnParseFormGloss{b}{nominative masculine plural}{ὀκνηρός}{lazy} τῷ πνεύματι ζέοντες,\FnParseFormGloss{c}{present active participle nominative masculine plural}{ζέω}{to have great fervor} τῷ κυρίῳ δουλεύοντες,\FnParseFormGloss{d}{present active participle nominative masculine plural}{δουλεύω}{to serve} 
\MainTextVerseMark{12}{12}τῇ ἐλπίδι χαίροντες,\FnParse{a}{present active participle nominative masculine plural} τῇ θλίψει ὑπομένοντες,\FnParseFormGloss{b}{present active participle nominative masculine plural}{ὑπομένω}{to stay behind} τῇ προσευχῇ προσκαρτεροῦντες,\FnParseFormGloss{c}{present active participle nominative masculine plural}{προσκαρτερέω}{to join} 
\MainTextVerseMark{12}{13}ταῖς χρείαις τῶν ἁγίων κοινωνοῦντες,\FnParseFormGloss{a}{present active participle nominative masculine plural}{κοινωνέω}{to share in} τὴν φιλοξενίαν\FnParseFormGloss{b}{accusative feminine singular}{φιλοξενία}{hospitality} διώκοντες.\FnParse{c}{present active participle nominative masculine plural} 
\MainTextVerseMark{12}{14}εὐλογεῖτε\FnParse{a}{present active imperative 2nd plural} τοὺς ⸀διώκοντας,\FnParse{b}{present active participle accusative masculine plural} εὐλογεῖτε\FnParse{c}{present active imperative 2nd plural} καὶ μὴ καταρᾶσθε.\FnParseFormGloss{d}{present middle imperative 2nd plural}{καταράομαι}{to curse} 
\MainTextVerseMark{12}{15}χαίρειν\FnParse{a}{present active infinitive} μετὰ ⸀χαιρόντων,\FnParse{b}{present active participle genitive masculine plural} κλαίειν\FnParse{c}{present active infinitive} μετὰ κλαιόντων.\FnParse{d}{present active participle genitive masculine plural} 
\MainTextVerseMark{12}{16}τὸ αὐτὸ εἰς ἀλλήλους φρονοῦντες,\FnParseFormGloss{a}{present active participle nominative masculine plural}{φρονέω}{to think} μὴ τὰ ὑψηλὰ\FnParseFormGloss{b}{accusative neuter plural}{ὑψηλός}{high} φρονοῦντες\FnParseFormGloss{c}{present active participle nominative masculine plural}{φρονέω}{to think} ἀλλὰ τοῖς ταπεινοῖς\FnParseFormGloss{d}{dative masculine plural}{ταπεινός}{humble} συναπαγόμενοι.\FnParseFormGloss{e}{present passive participle nominative masculine plural}{συναπάγομαι}{to be led away} μὴ γίνεσθε\FnParse{f}{present middle imperative 2nd plural} φρόνιμοι\FnParseFormGloss{g}{nominative masculine plural}{φρόνιμος}{wise} παρ’ ἑαυτοῖς. 
\MainTextVerseMark{12}{17}μηδενὶ κακὸν ἀντὶ\FnFormGloss{a}{ἀντί}{in exchange for} κακοῦ ἀποδιδόντες·\FnParse{b}{present active participle nominative masculine plural} προνοούμενοι\FnParseFormGloss{c}{present middle participle nominative masculine plural}{προνοέω}{to provide for} καλὰ ἐνώπιον πάντων ἀνθρώπων· 
\MainTextVerseMark{12}{18}εἰ δυνατόν, τὸ ἐξ ὑμῶν μετὰ πάντων ἀνθρώπων εἰρηνεύοντες·\FnParseFormGloss{a}{present active participle nominative masculine plural}{εἰρηνεύω}{to live in peace} 
\MainTextVerseMark{12}{19}μὴ ἑαυτοὺς ἐκδικοῦντες,\FnParseFormGloss{a}{present active participle nominative masculine plural}{ἐκδικέω}{to avenge} ἀγαπητοί, ἀλλὰ δότε\FnParse{b}{aorist active imperative 2nd plural} τόπον τῇ ὀργῇ, γέγραπται\FnParse{c}{perfect passive indicative 3rd singular} γάρ· Ἐμοὶ ἐκδίκησις,\FnParseFormGloss{d}{nominative feminine singular}{ἐκδίκησις}{justice} ἐγὼ ἀνταποδώσω,\FnParseFormGloss{e}{future active indicative 1st singular}{ἀνταποδίδωμι}{to repay} λέγει\FnParse{f}{present active indicative 3rd singular} κύριος. 
\MainTextVerseMark{12}{20}⸂ἀλλὰ ἐὰν⸃ πεινᾷ\FnParseFormGloss{a}{present active subjunctive 3rd singular}{πεινάω}{to be hungry} ὁ ἐχθρός σου, ψώμιζε\FnParseFormGloss{b}{present active imperative 2nd singular}{ψωμίζω}{to feed} αὐτόν· ἐὰν διψᾷ,\FnParseFormGloss{c}{present active subjunctive 3rd singular}{διψάω}{to be thirsty} πότιζε\FnParseFormGloss{d}{present active imperative 2nd singular}{ποτίζω}{to give a drink} αὐτόν· τοῦτο γὰρ ποιῶν\FnParse{e}{present active participle nominative masculine singular} ἄνθρακας\FnParseFormGloss{f}{accusative masculine plural}{ἄνθραξ}{coal} πυρὸς σωρεύσεις\FnParseFormGloss{g}{future active indicative 2nd singular}{σωρεύω}{to heap up} ἐπὶ τὴν κεφαλὴν αὐτοῦ. 
\MainTextVerseMark{12}{21}μὴ νικῶ\FnParseFormGloss{a}{present passive imperative 2nd singular}{νικάω}{to overcome} ὑπὸ τοῦ κακοῦ, ἀλλὰ νίκα\FnParseFormGloss{b}{present active imperative 2nd singular}{νικάω}{to overcome} ἐν τῷ ἀγαθῷ τὸ κακόν. 
\ChapterHeader{13}{Chapter 13}

\MainTextVerseMark{13}{1}Πᾶσα ψυχὴ ἐξουσίαις ὑπερεχούσαις\FnParseFormGloss{a}{present active participle dative feminine plural}{ὑπερέχω}{to govern} ὑποτασσέσθω,\FnParse{b}{present passive imperative 3rd singular} οὐ γὰρ ἔστιν\FnParse{c}{present active indicative 3rd singular} ἐξουσία εἰ μὴ ὑπὸ θεοῦ, αἱ δὲ ⸀οὖσαι\FnParse{d}{present active participle nominative feminine plural} ⸀ὑπὸ θεοῦ τεταγμέναι\FnParseFormGloss{e}{perfect passive participle nominative feminine plural}{τάσσω}{to appoint} εἰσίν.\FnParse{f}{present active indicative 3rd plural} 
\MainTextVerseMark{13}{2}ὥστε ὁ ἀντιτασσόμενος\FnParseFormGloss{a}{present middle participle nominative masculine singular}{ἀντιτάσσομαι}{to oppose} τῇ ἐξουσίᾳ τῇ τοῦ θεοῦ διαταγῇ\FnParseFormGloss{b}{dative feminine singular}{διαταγή}{putting into effect} ἀνθέστηκεν,\FnParseFormGloss{c}{perfect active indicative 3rd singular}{ἀνθίστημι}{to resist} οἱ δὲ ἀνθεστηκότες\FnParseFormGloss{d}{perfect active participle nominative masculine plural}{ἀνθίστημι}{to resist} ἑαυτοῖς κρίμα\FnParseFormGloss{e}{accusative neuter singular}{κρίμα}{judgment} λήμψονται.\FnParse{f}{future middle indicative 3rd plural} 
\MainTextVerseMark{13}{3}οἱ γὰρ ἄρχοντες οὐκ εἰσὶν\FnParse{a}{present active indicative 3rd plural} φόβος ⸂τῷ ἀγαθῷ ἔργῳ⸃ ἀλλὰ ⸂τῷ κακῷ⸃. θέλεις\FnParse{b}{present active indicative 2nd singular} δὲ μὴ φοβεῖσθαι\FnParse{c}{present middle infinitive} τὴν ἐξουσίαν; τὸ ἀγαθὸν ποίει,\FnParse{d}{present active imperative 2nd singular} καὶ ἕξεις\FnParse{e}{future active indicative 2nd singular} ἔπαινον\FnParseFormGloss{f}{accusative masculine singular}{ἔπαινος}{praise} ἐξ αὐτῆς· 
\MainTextVerseMark{13}{4}θεοῦ γὰρ διάκονός\FnParseFormGloss{a}{nominative masculine singular}{διάκονος}{servant} ἐστιν\FnParse{b}{present active indicative 3rd singular} σοὶ εἰς τὸ ἀγαθόν. ἐὰν δὲ τὸ κακὸν ποιῇς,\FnParse{c}{present active subjunctive 2nd singular} φοβοῦ·\FnParse{d}{present middle imperative 2nd singular} οὐ γὰρ εἰκῇ\FnFormGloss{e}{εἰκῇ}{in vain} τὴν μάχαιραν\FnParseFormGloss{f}{accusative feminine singular}{μάχαιρα}{sword} φορεῖ·\FnParseFormGloss{g}{present active indicative 3rd singular}{φορέω}{to wear} θεοῦ γὰρ διάκονός\FnParseFormGloss{h}{nominative masculine singular}{διάκονος}{servant} ἐστιν,\FnParse{i}{present active indicative 3rd singular} ἔκδικος\FnParseFormGloss{j}{nominative masculine singular}{ἔκδικος}{punishing} εἰς ὀργὴν τῷ τὸ κακὸν πράσσοντι.\FnParse{k}{present active participle dative masculine singular} 
\MainTextVerseMark{13}{5}διὸ ἀνάγκη\FnParseFormGloss{a}{nominative feminine singular}{ἀνάγκη}{necessity} ὑποτάσσεσθαι,\FnParse{b}{present passive infinitive} οὐ μόνον διὰ τὴν ὀργὴν ἀλλὰ καὶ διὰ τὴν συνείδησιν, 
\MainTextVerseMark{13}{6}διὰ τοῦτο γὰρ καὶ φόρους\FnParseFormGloss{a}{accusative masculine plural}{φόρος}{tax} τελεῖτε,\FnParseFormGloss{b}{present active indicative 2nd plural}{τελέω}{to finish} λειτουργοὶ\FnParseFormGloss{c}{nominative masculine plural}{λειτουργός}{servant} γὰρ θεοῦ εἰσιν\FnParse{d}{present active indicative 3rd plural} εἰς αὐτὸ τοῦτο προσκαρτεροῦντες.\FnParseFormGloss{e}{present active participle nominative masculine plural}{προσκαρτερέω}{to join} 
\MainTextVerseMark{13}{7}⸀ἀπόδοτε\FnParse{a}{aorist active imperative 2nd plural} πᾶσι τὰς ὀφειλάς,\FnParseFormGloss{b}{accusative feminine plural}{ὀφειλή}{debt} τῷ τὸν φόρον\FnParseFormGloss{c}{accusative masculine singular}{φόρος}{tax} τὸν φόρον,\FnParseFormGloss{d}{accusative masculine singular}{φόρος}{tax} τῷ τὸ τέλος τὸ τέλος, τῷ τὸν φόβον τὸν φόβον, τῷ τὴν τιμὴν τὴν τιμήν. 
\MainTextVerseMark{13}{8}Μηδενὶ μηδὲν ὀφείλετε,\FnParse{a}{present active imperative 2nd plural} εἰ μὴ τὸ ⸂ἀλλήλους ἀγαπᾶν⸃·\FnParse{b}{present active infinitive} ὁ γὰρ ἀγαπῶν\FnParse{c}{present active participle nominative masculine singular} τὸν ἕτερον νόμον πεπλήρωκεν.\FnParse{d}{perfect active indicative 3rd singular} 
\MainTextVerseMark{13}{9}τὸ γάρ· Οὐ μοιχεύσεις,\FnParseFormGloss{a}{future active indicative 2nd singular}{μοιχεύω}{to commit adultery} Οὐ φονεύσεις,\FnParseFormGloss{b}{future active indicative 2nd singular}{φονεύω}{to commit murder} Οὐ κλέψεις,\FnParseFormGloss{c}{future active indicative 2nd singular}{κλέπτω}{steal} Οὐκ ἐπιθυμήσεις,\FnParseFormGloss{d}{future active indicative 2nd singular}{ἐπιθυμέω}{to long for} καὶ εἴ τις ἑτέρα ἐντολή, ἐν ⸂τῷ λόγῳ τούτῳ⸃ ἀνακεφαλαιοῦται,\FnParseFormGloss{e}{present passive indicative 3rd singular}{ἀνακεφαλαιόω}{to bring together under one head} ⸂ἐν τῷ⸃· Ἀγαπήσεις\FnParse{f}{future active indicative 2nd singular} τὸν πλησίον\FnFormGloss{g}{πλησίον}{near} σου ὡς σεαυτόν. 
\MainTextVerseMark{13}{10}ἡ ἀγάπη τῷ πλησίον\FnFormGloss{a}{πλησίον}{near} κακὸν οὐκ ἐργάζεται·\FnParse{b}{present middle indicative 3rd singular} πλήρωμα\FnParseFormGloss{c}{nominative neuter singular}{πλήρωμα}{fullness} οὖν νόμου ἡ ἀγάπη. 
\MainTextVerseMark{13}{11}Καὶ τοῦτο εἰδότες\FnParse{a}{perfect active participle nominative masculine plural} τὸν καιρόν, ὅτι ὥρα ⸂ἤδη ὑμᾶς⸃ ἐξ ὕπνου\FnParseFormGloss{b}{genitive masculine singular}{ὕπνος}{sleep} ἐγερθῆναι,\FnParse{c}{aorist passive infinitive} νῦν γὰρ ἐγγύτερον\FnFormGloss{d}{ἐγγύτερον}{nearer} ἡμῶν ἡ σωτηρία ἢ ὅτε ἐπιστεύσαμεν.\FnParse{e}{aorist active indicative 1st plural} 
\MainTextVerseMark{13}{12}ἡ νὺξ προέκοψεν,\FnParseFormGloss{a}{aorist active indicative 3rd singular}{προκόπτω}{to go ahead} ἡ δὲ ἡμέρα ἤγγικεν.\FnParse{b}{perfect active indicative 3rd singular} ⸀ἀποβαλώμεθα\FnParseFormGloss{c}{aorist middle subjunctive 1st plural}{ἀποβάλλω}{to throw away} οὖν τὰ ἔργα τοῦ σκότους, ⸂ἐνδυσώμεθα\FnParseFormGloss{d}{aorist middle subjunctive 1st plural}{ἐνδύω}{to clothe} δὲ⸃ τὰ ὅπλα\FnParseFormGloss{e}{accusative neuter plural}{ὅπλον}{instrument} τοῦ φωτός. 
\MainTextVerseMark{13}{13}ὡς ἐν ἡμέρᾳ εὐσχημόνως\FnFormGloss{a}{εὐσχημόνως}{decently} περιπατήσωμεν,\FnParse{b}{aorist active subjunctive 1st plural} μὴ κώμοις\FnParseFormGloss{c}{dative masculine plural}{κῶμος}{orgy} καὶ μέθαις,\FnParseFormGloss{d}{dative feminine plural}{μέθη}{drunkenness} μὴ κοίταις\FnParseFormGloss{e}{dative feminine plural}{κοίτη}{bed} καὶ ἀσελγείαις,\FnParseFormGloss{f}{dative feminine plural}{ἀσέλγεια}{debauchery} μὴ ἔριδι\FnParseFormGloss{g}{dative feminine singular}{ἔρις}{quarrel} καὶ ζήλῳ,\FnParseFormGloss{h}{dative masculine singular}{ζῆλος}{zeal} 
\MainTextVerseMark{13}{14}ἀλλὰ ἐνδύσασθε\FnParseFormGloss{a}{aorist middle imperative 2nd plural}{ἐνδύω}{to clothe} τὸν κύριον Ἰησοῦν Χριστόν, καὶ τῆς σαρκὸς πρόνοιαν\FnParseFormGloss{b}{accusative feminine singular}{πρόνοια}{foresight} μὴ ποιεῖσθε\FnParse{c}{present middle imperative 2nd plural} εἰς ἐπιθυμίας. 
\ChapterHeader{14}{Chapter 14}

\MainTextVerseMark{14}{1}Τὸν δὲ ἀσθενοῦντα\FnParse{a}{present active participle accusative masculine singular} τῇ πίστει προσλαμβάνεσθε,\FnParseFormGloss{b}{present middle imperative 2nd plural}{προσλαμβάνομαι}{to take aside} μὴ εἰς διακρίσεις\FnParseFormGloss{c}{accusative feminine plural}{διάκρισις}{distinguishing} διαλογισμῶν.\FnParseFormGloss{d}{genitive masculine plural}{διαλογισμός}{thought} 
\MainTextVerseMark{14}{2}ὃς μὲν πιστεύει\FnParse{a}{present active indicative 3rd singular} φαγεῖν\FnParse{b}{aorist active infinitive} πάντα, ὁ δὲ ἀσθενῶν\FnParse{c}{present active participle nominative masculine singular} λάχανα\FnParseFormGloss{d}{accusative neuter plural}{λάχανον}{plant} ἐσθίει.\FnParse{e}{present active indicative 3rd singular} 
\MainTextVerseMark{14}{3}ὁ ἐσθίων\FnParse{a}{present active participle nominative masculine singular} τὸν μὴ ἐσθίοντα\FnParse{b}{present active participle accusative masculine singular} μὴ ἐξουθενείτω,\FnParseFormGloss{c}{present active imperative 3rd singular}{ἐξουθενέω}{to treat with contempt} ⸂ὁ δὲ⸃ μὴ ἐσθίων\FnParse{d}{present active participle nominative masculine singular} τὸν ἐσθίοντα\FnParse{e}{present active participle accusative masculine singular} μὴ κρινέτω,\FnParse{f}{present active imperative 3rd singular} ὁ θεὸς γὰρ αὐτὸν προσελάβετο.\FnParseFormGloss{g}{aorist middle indicative 3rd singular}{προσλαμβάνομαι}{to take aside} 
\MainTextVerseMark{14}{4}σὺ τίς εἶ\FnParse{a}{present active indicative 2nd singular} ὁ κρίνων\FnParse{b}{present active participle nominative masculine singular} ἀλλότριον\FnParseFormGloss{c}{accusative masculine singular}{ἀλλότριος}{belonging to another} οἰκέτην;\FnParseFormGloss{d}{accusative masculine singular}{οἰκέτης}{house servant} τῷ ἰδίῳ κυρίῳ στήκει\FnParseFormGloss{e}{present active indicative 3rd singular}{στήκω}{to stand} ἢ πίπτει·\FnParse{f}{present active indicative 3rd singular} σταθήσεται\FnParse{g}{future passive indicative 3rd singular} δέ, ⸂δυνατεῖ\FnParseFormGloss{h}{present active indicative 3rd singular}{δυνατέω}{to be able} γὰρ⸃ ὁ ⸀κύριος στῆσαι\FnParse{i}{aorist active infinitive} αὐτόν. 
\MainTextVerseMark{14}{5}Ὃς ⸀μὲν κρίνει\FnParse{a}{present active indicative 3rd singular} ἡμέραν παρ’ ἡμέραν, ὃς δὲ κρίνει\FnParse{b}{present active indicative 3rd singular} πᾶσαν ἡμέραν· ἕκαστος ἐν τῷ ἰδίῳ νοῒ\FnParseFormGloss{c}{dative masculine singular}{νοῦς}{mind} πληροφορείσθω·\FnParseFormGloss{d}{present passive imperative 3rd singular}{πληροφορέω}{to fulfill} 
\MainTextVerseMark{14}{6}ὁ φρονῶν\FnParseFormGloss{a}{present active participle nominative masculine singular}{φρονέω}{to think} τὴν ἡμέραν κυρίῳ ⸀φρονεῖ.\FnParseFormGloss{b}{present active indicative 3rd singular}{φρονέω}{to think} καὶ ὁ ἐσθίων\FnParse{c}{present active participle nominative masculine singular} κυρίῳ ἐσθίει,\FnParse{d}{present active indicative 3rd singular} εὐχαριστεῖ\FnParse{e}{present active indicative 3rd singular} γὰρ τῷ θεῷ· καὶ ὁ μὴ ἐσθίων\FnParse{f}{present active participle nominative masculine singular} κυρίῳ οὐκ ἐσθίει,\FnParse{g}{present active indicative 3rd singular} καὶ εὐχαριστεῖ\FnParse{h}{present active indicative 3rd singular} τῷ θεῷ. 
\MainTextVerseMark{14}{7}Οὐδεὶς γὰρ ἡμῶν ἑαυτῷ ζῇ,\FnParse{a}{present active indicative 3rd singular} καὶ οὐδεὶς ἑαυτῷ ἀποθνῄσκει·\FnParse{b}{present active indicative 3rd singular} 
\MainTextVerseMark{14}{8}ἐάν τε γὰρ ζῶμεν,\FnParse{a}{present active subjunctive 1st plural} τῷ κυρίῳ ζῶμεν,\FnParse{b}{present active indicative 1st plural} ἐάν τε ἀποθνῄσκωμεν,\FnParse{c}{present active subjunctive 1st plural} τῷ κυρίῳ ἀποθνῄσκομεν.\FnParse{d}{present active indicative 1st plural} ἐάν τε οὖν ζῶμεν\FnParse{e}{present active subjunctive 1st plural} ἐάν τε ἀποθνῄσκωμεν,\FnParse{f}{present active subjunctive 1st plural} τοῦ κυρίου ἐσμέν.\FnParse{g}{present active indicative 1st plural} 
\MainTextVerseMark{14}{9}εἰς τοῦτο γὰρ ⸀Χριστὸς ⸀ἀπέθανεν\FnParse{a}{aorist active indicative 3rd singular} καὶ ἔζησεν\FnParse{b}{aorist active indicative 3rd singular} ἵνα καὶ νεκρῶν καὶ ζώντων\FnParse{c}{present active participle genitive masculine plural} κυριεύσῃ.\FnParseFormGloss{d}{aorist active subjunctive 3rd singular}{κυριεύω}{to lord over} 
\MainTextVerseMark{14}{10}Σὺ δὲ τί κρίνεις\FnParse{a}{present active indicative 2nd singular} τὸν ἀδελφόν σου; ἢ καὶ σὺ τί ἐξουθενεῖς\FnParseFormGloss{b}{present active indicative 2nd singular}{ἐξουθενέω}{to treat with contempt} τὸν ἀδελφόν σου; πάντες γὰρ παραστησόμεθα\FnParse{c}{future middle indicative 1st plural} τῷ βήματι\FnParseFormGloss{d}{dative neuter singular}{βῆμα}{judicial court} τοῦ ⸀θεοῦ, 
\MainTextVerseMark{14}{11}γέγραπται\FnParse{a}{perfect passive indicative 3rd singular} γάρ· Ζῶ\FnParse{b}{present active indicative 1st singular} ἐγώ, λέγει\FnParse{c}{present active indicative 3rd singular} κύριος, ὅτι ἐμοὶ κάμψει\FnParseFormGloss{d}{future active indicative 3rd singular}{κάμπτω}{to bend} πᾶν γόνυ,\FnParseFormGloss{e}{nominative neuter singular}{γόνυ}{knee} καὶ ⸂πᾶσα γλῶσσα ἐξομολογήσεται⸃\FnParseFormGloss{f}{future middle indicative 3rd singular}{ἐξομολογέω}{to consent} τῷ θεῷ. 
\MainTextVerseMark{14}{12}⸀ἄρα ἕκαστος ἡμῶν περὶ ἑαυτοῦ λόγον ⸀δώσει.\FnParse{a}{future active indicative 3rd singular} 
\MainTextVerseMark{14}{13}Μηκέτι\FnFormGloss{a}{μηκέτι}{no longer} οὖν ἀλλήλους κρίνωμεν·\FnParse{b}{present active subjunctive 1st plural} ἀλλὰ τοῦτο κρίνατε\FnParse{c}{aorist active imperative 2nd plural} μᾶλλον, τὸ μὴ τιθέναι\FnParse{d}{present active infinitive} πρόσκομμα\FnParseFormGloss{e}{accusative neuter singular}{πρόσκομμα}{stumbling block} τῷ ἀδελφῷ ἢ σκάνδαλον.\FnParseFormGloss{f}{accusative neuter singular}{σκάνδαλον}{stumbling block} 
\MainTextVerseMark{14}{14}οἶδα\FnParse{a}{perfect active indicative 1st singular} καὶ πέπεισμαι\FnParse{b}{perfect passive indicative 1st singular} ἐν κυρίῳ Ἰησοῦ ὅτι οὐδὲν κοινὸν\FnParseFormGloss{c}{nominative neuter singular}{κοινός}{common} δι’ ⸀ἑαυτοῦ· εἰ μὴ τῷ λογιζομένῳ\FnParse{d}{present middle participle dative masculine singular} τι κοινὸν\FnParseFormGloss{e}{nominative neuter singular}{κοινός}{common} εἶναι,\FnParse{f}{present active infinitive} ἐκείνῳ κοινόν.\FnParseFormGloss{g}{nominative neuter singular}{κοινός}{common} 
\MainTextVerseMark{14}{15}εἰ ⸀γὰρ διὰ βρῶμα\FnParseFormGloss{a}{accusative neuter singular}{βρῶμα}{food} ὁ ἀδελφός σου λυπεῖται,\FnParseFormGloss{b}{present passive indicative 3rd singular}{λυπέω}{to cause sorrow} οὐκέτι κατὰ ἀγάπην περιπατεῖς.\FnParse{c}{present active indicative 2nd singular} μὴ τῷ βρώματί\FnParseFormGloss{d}{dative neuter singular}{βρῶμα}{food} σου ἐκεῖνον ἀπόλλυε\FnParse{e}{present active imperative 2nd singular} ὑπὲρ οὗ Χριστὸς ἀπέθανεν.\FnParse{f}{aorist active indicative 3rd singular} 
\MainTextVerseMark{14}{16}μὴ βλασφημείσθω\FnParse{a}{present passive imperative 3rd singular} οὖν ὑμῶν τὸ ἀγαθόν. 
\MainTextVerseMark{14}{17}οὐ γάρ ἐστιν\FnParse{a}{present active indicative 3rd singular} ἡ βασιλεία τοῦ θεοῦ βρῶσις\FnParseFormGloss{b}{nominative feminine singular}{βρῶσις}{consumable} καὶ πόσις,\FnParseFormGloss{c}{nominative feminine singular}{πόσις}{drinking} ἀλλὰ δικαιοσύνη καὶ εἰρήνη καὶ χαρὰ ἐν πνεύματι ἁγίῳ· 
\MainTextVerseMark{14}{18}ὁ γὰρ ἐν ⸀τούτῳ δουλεύων\FnParseFormGloss{a}{present active participle nominative masculine singular}{δουλεύω}{to serve} τῷ Χριστῷ εὐάρεστος\FnParseFormGloss{b}{nominative masculine singular}{εὐάρεστος}{pleasing} τῷ θεῷ καὶ δόκιμος\FnParseFormGloss{c}{nominative masculine singular}{δόκιμος}{approved by testing} τοῖς ἀνθρώποις. 
\MainTextVerseMark{14}{19}ἄρα οὖν τὰ τῆς εἰρήνης διώκωμεν\FnParse{a}{present active subjunctive 1st plural} καὶ τὰ τῆς οἰκοδομῆς\FnParseFormGloss{b}{genitive feminine singular}{οἰκοδομή}{building} τῆς εἰς ἀλλήλους. 
\MainTextVerseMark{14}{20}μὴ ἕνεκεν\FnFormGloss{a}{ἕνεκεν}{for} βρώματος\FnParseFormGloss{b}{genitive neuter singular}{βρῶμα}{food} κατάλυε\FnParseFormGloss{c}{present active imperative 2nd singular}{καταλύω}{throw down} τὸ ἔργον τοῦ θεοῦ. πάντα μὲν καθαρά,\FnParseFormGloss{d}{nominative neuter plural}{καθαρός}{clean} ἀλλὰ κακὸν τῷ ἀνθρώπῳ τῷ διὰ προσκόμματος\FnParseFormGloss{e}{genitive neuter singular}{πρόσκομμα}{stumbling block} ἐσθίοντι.\FnParse{f}{present active participle dative masculine singular} 
\MainTextVerseMark{14}{21}καλὸν τὸ μὴ φαγεῖν\FnParse{a}{aorist active infinitive} κρέα\FnParseFormGloss{b}{accusative neuter plural}{κρέας}{meat} μηδὲ πιεῖν\FnParse{c}{aorist active infinitive} οἶνον μηδὲ ἐν ᾧ ὁ ἀδελφός σου προσκόπτει\FnParseFormGloss{d}{present active indicative 3rd singular}{προσκόπτω}{to strike} ⸂ἢ σκανδαλίζεται\FnParse{e}{present passive indicative 3rd singular} ἢ ἀσθενεῖ⸃·\FnParse{f}{present active indicative 3rd singular} 
\MainTextVerseMark{14}{22}σὺ πίστιν ⸀ἣν ἔχεις\FnParse{a}{present active indicative 2nd singular} κατὰ σεαυτὸν ἔχε\FnParse{b}{present active imperative 2nd singular} ἐνώπιον τοῦ θεοῦ. μακάριος ὁ μὴ κρίνων\FnParse{c}{present active participle nominative masculine singular} ἑαυτὸν ἐν ᾧ δοκιμάζει·\FnParseFormGloss{d}{present active indicative 3rd singular}{δοκιμάζω}{to test} 
\MainTextVerseMark{14}{23}ὁ δὲ διακρινόμενος\FnParseFormGloss{a}{present middle participle nominative masculine singular}{διακρίνω}{to make a distinction} ἐὰν φάγῃ\FnParse{b}{aorist active subjunctive 3rd singular} κατακέκριται,\FnParseFormGloss{c}{perfect passive indicative 3rd singular}{κατακρίνω}{to condemn} ὅτι οὐκ ἐκ πίστεως· πᾶν δὲ ὃ οὐκ ἐκ πίστεως ἁμαρτία ⸀ἐστίν.\FnParse{d}{present active indicative 3rd singular} 
\ChapterHeader{15}{Chapter 15}

\MainTextVerseMark{15}{1}Ὀφείλομεν\FnParse{a}{present active indicative 1st plural} δὲ ἡμεῖς οἱ δυνατοὶ τὰ ἀσθενήματα\FnParseFormGloss{b}{accusative neuter plural}{ἀσθένημα}{failing} τῶν ἀδυνάτων\FnParseFormGloss{c}{genitive masculine plural}{ἀδύνατος}{impossible} βαστάζειν,\FnParseFormGloss{d}{present active infinitive}{βαστάζω}{to carry} καὶ μὴ ἑαυτοῖς ἀρέσκειν.\FnParseFormGloss{e}{present active infinitive}{ἀρέσκω}{to please} 
\MainTextVerseMark{15}{2}ἕκαστος ἡμῶν τῷ πλησίον\FnFormGloss{a}{πλησίον}{near} ἀρεσκέτω\FnParseFormGloss{b}{present active imperative 3rd singular}{ἀρέσκω}{to please} εἰς τὸ ἀγαθὸν πρὸς οἰκοδομήν·\FnParseFormGloss{c}{accusative feminine singular}{οἰκοδομή}{building} 
\MainTextVerseMark{15}{3}καὶ γὰρ ὁ Χριστὸς οὐχ ἑαυτῷ ἤρεσεν·\FnParseFormGloss{a}{aorist active indicative 3rd singular}{ἀρέσκω}{to please} ἀλλὰ καθὼς γέγραπται·\FnParse{b}{perfect passive indicative 3rd singular} Οἱ ὀνειδισμοὶ\FnParseFormGloss{c}{nominative masculine plural}{ὀνειδισμός}{disgrace} τῶν ὀνειδιζόντων\FnParseFormGloss{d}{present active participle genitive masculine plural}{ὀνειδίζω}{to heap insults on} σε ἐπέπεσαν\FnParseFormGloss{e}{aorist active indicative 3rd plural}{ἐπιπίπτω}{to fall upon} ἐπ’ ἐμέ. 
\MainTextVerseMark{15}{4}ὅσα γὰρ προεγράφη,\FnParseFormGloss{a}{aorist passive indicative 3rd singular}{προγράφω}{to write beforehand} ⸀εἰς τὴν ἡμετέραν\FnParseFormGloss{b}{accusative feminine singular}{ἡμέτερος}{our} διδασκαλίαν\FnParseFormGloss{c}{accusative feminine singular}{διδασκαλία}{teaching} ⸀ἐγράφη,\FnParse{d}{aorist passive indicative 3rd singular} ἵνα διὰ τῆς ὑπομονῆς καὶ διὰ τῆς παρακλήσεως\FnParseFormGloss{e}{genitive feminine singular}{παράκλησις}{encouragement} τῶν γραφῶν τὴν ἐλπίδα ἔχωμεν.\FnParse{f}{present active subjunctive 1st plural} 
\MainTextVerseMark{15}{5}ὁ δὲ θεὸς τῆς ὑπομονῆς καὶ τῆς παρακλήσεως\FnParseFormGloss{a}{genitive feminine singular}{παράκλησις}{encouragement} δῴη\FnParse{b}{aorist active optative 3rd singular} ὑμῖν τὸ αὐτὸ φρονεῖν\FnParseFormGloss{c}{present active infinitive}{φρονέω}{to think} ἐν ἀλλήλοις κατὰ ⸂Χριστὸν Ἰησοῦν⸃, 
\MainTextVerseMark{15}{6}ἵνα ὁμοθυμαδὸν\FnFormGloss{a}{ὁμοθυμαδόν}{united} ἐν ἑνὶ στόματι δοξάζητε\FnParse{b}{present active subjunctive 2nd plural} τὸν θεὸν καὶ πατέρα τοῦ κυρίου ἡμῶν Ἰησοῦ Χριστοῦ. 
\MainTextVerseMark{15}{7}Διὸ προσλαμβάνεσθε\FnParseFormGloss{a}{present middle imperative 2nd plural}{προσλαμβάνομαι}{to take aside} ἀλλήλους, καθὼς καὶ ὁ Χριστὸς προσελάβετο\FnParseFormGloss{b}{aorist middle indicative 3rd singular}{προσλαμβάνομαι}{to take aside} ⸀ὑμᾶς, εἰς δόξαν ⸀τοῦ θεοῦ. 
\MainTextVerseMark{15}{8}λέγω\FnParse{a}{present active indicative 1st singular} ⸀γὰρ ⸀Χριστὸν διάκονον\FnParseFormGloss{b}{accusative masculine singular}{διάκονος}{servant} ⸀γεγενῆσθαι\FnParse{c}{perfect passive infinitive} περιτομῆς ὑπὲρ ἀληθείας θεοῦ, εἰς τὸ βεβαιῶσαι\FnParseFormGloss{d}{aorist active infinitive}{βεβαιόω}{to confirm} τὰς ἐπαγγελίας τῶν πατέρων, 
\MainTextVerseMark{15}{9}τὰ δὲ ἔθνη ὑπὲρ ἐλέους\FnParseFormGloss{a}{genitive neuter singular}{ἔλεος}{mercy} δοξάσαι\FnParse{b}{aorist active infinitive} τὸν θεόν· καθὼς γέγραπται·\FnParse{c}{perfect passive indicative 3rd singular} Διὰ τοῦτο ἐξομολογήσομαί\FnParseFormGloss{d}{future middle indicative 1st singular}{ἐξομολογέω}{to consent} σοι ἐν ἔθνεσι, καὶ τῷ ὀνόματί σου ψαλῶ.\FnParseFormGloss{e}{future active indicative 1st singular}{ψάλλω}{to sing hymns} 
\MainTextVerseMark{15}{10}καὶ πάλιν λέγει·\FnParse{a}{present active indicative 3rd singular} Εὐφράνθητε,\FnParseFormGloss{b}{aorist passive imperative 2nd plural}{εὐφραίνω}{to cause celebration} ἔθνη, μετὰ τοῦ λαοῦ αὐτοῦ. 
\MainTextVerseMark{15}{11}καὶ πάλιν· Αἰνεῖτε,\FnParseFormGloss{a}{present active imperative 2nd plural}{αἰνέω}{to praise} ⸂πάντα τὰ ἔθνη, τὸν κύριον⸃, καὶ ⸀ἐπαινεσάτωσαν\FnParseFormGloss{b}{aorist active imperative 3rd plural}{ἐπαινέω}{to praise} αὐτὸν πάντες οἱ λαοί. 
\MainTextVerseMark{15}{12}καὶ πάλιν Ἠσαΐας\FnParseFormGloss{a}{nominative masculine singular}{Ἠσαΐας}{Isaiah} λέγει·\FnParse{b}{present active indicative 3rd singular} Ἔσται\FnParse{c}{future middle indicative 3rd singular} ἡ ῥίζα\FnParseFormGloss{d}{nominative feminine singular}{ῥίζα}{root} τοῦ Ἰεσσαί,\FnParseFormGloss{e}{genitive masculine singular}{Ἰεσσαί}{Jesse} καὶ ὁ ἀνιστάμενος\FnParse{f}{present middle participle nominative masculine singular} ἄρχειν\FnParse{g}{present active infinitive} ἐθνῶν· ἐπ’ αὐτῷ ἔθνη ἐλπιοῦσιν.\FnParse{h}{future active indicative 3rd plural} 
\MainTextVerseMark{15}{13}ὁ δὲ θεὸς τῆς ἐλπίδος πληρώσαι\FnParse{a}{aorist active optative 3rd singular} ὑμᾶς πάσης χαρᾶς καὶ εἰρήνης ἐν τῷ πιστεύειν,\FnParse{b}{present active infinitive} εἰς τὸ περισσεύειν\FnParse{c}{present active infinitive} ὑμᾶς ἐν τῇ ἐλπίδι ἐν δυνάμει πνεύματος ἁγίου. 
\MainTextVerseMark{15}{14}Πέπεισμαι\FnParse{a}{perfect passive indicative 1st singular} δέ, ἀδελφοί μου, καὶ αὐτὸς ἐγὼ περὶ ὑμῶν, ὅτι καὶ αὐτοὶ μεστοί\FnParseFormGloss{b}{nominative masculine plural}{μεστός}{full} ἐστε\FnParse{c}{present active indicative 2nd plural} ἀγαθωσύνης,\FnParseFormGloss{d}{genitive feminine singular}{ἀγαθωσύνη}{goodness} πεπληρωμένοι\FnParse{e}{perfect passive participle nominative masculine plural} ⸀πάσης γνώσεως,\FnParseFormGloss{f}{genitive feminine singular}{γνῶσις}{knowledge} δυνάμενοι\FnParse{g}{present middle participle nominative masculine plural} καὶ ⸀ἀλλήλους νουθετεῖν.\FnParseFormGloss{h}{present active infinitive}{νουθετέω}{to warn} 
\MainTextVerseMark{15}{15}⸀τολμηρότερον\FnFormGloss{a}{τολμηρότερον}{rather boldly} δὲ ἔγραψα\FnParse{b}{aorist active indicative 1st singular} ⸀ὑμῖν ἀπὸ μέρους, ὡς ἐπαναμιμνῄσκων\FnParseFormGloss{c}{present active participle nominative masculine singular}{ἐπαναμιμνῄσκω}{to remind again} ὑμᾶς, διὰ τὴν χάριν τὴν δοθεῖσάν\FnParse{d}{aorist passive participle accusative feminine singular} μοι ⸀ὑπὸ τοῦ θεοῦ 
\MainTextVerseMark{15}{16}εἰς τὸ εἶναί\FnParse{a}{present active infinitive} με λειτουργὸν\FnParseFormGloss{b}{accusative masculine singular}{λειτουργός}{servant} ⸂Χριστοῦ Ἰησοῦ⸃ εἰς τὰ ἔθνη, ἱερουργοῦντα\FnParseFormGloss{c}{present active participle accusative masculine singular}{ἱερουργέω}{to perform priestly duty} τὸ εὐαγγέλιον τοῦ θεοῦ, ἵνα γένηται\FnParse{d}{aorist middle subjunctive 3rd singular} ἡ προσφορὰ\FnParseFormGloss{e}{nominative feminine singular}{προσφορά}{offering} τῶν ἐθνῶν εὐπρόσδεκτος,\FnParseFormGloss{f}{nominative feminine singular}{εὐπρόσδεκτος}{acceptable} ἡγιασμένη\FnParseFormGloss{g}{perfect passive participle nominative feminine singular}{ἁγιάζω}{to sanctify} ἐν πνεύματι ἁγίῳ. 
\MainTextVerseMark{15}{17}ἔχω\FnParse{a}{present active indicative 1st singular} οὖν ⸀τὴν καύχησιν\FnParseFormGloss{b}{accusative feminine singular}{καύχησις}{boasting} ἐν Χριστῷ Ἰησοῦ τὰ πρὸς τὸν θεόν· 
\MainTextVerseMark{15}{18}οὐ γὰρ τολμήσω\FnParseFormGloss{a}{future active indicative 1st singular}{τολμάω}{to dare} ⸂τι λαλεῖν⸃\FnParse{b}{present active infinitive} ὧν οὐ κατειργάσατο\FnParseFormGloss{c}{aorist middle indicative 3rd singular}{κατεργάζομαι}{to produce} Χριστὸς δι’ ἐμοῦ εἰς ὑπακοὴν\FnParseFormGloss{d}{accusative feminine singular}{ὑπακοή}{obedience} ἐθνῶν, λόγῳ καὶ ἔργῳ, 
\MainTextVerseMark{15}{19}ἐν δυνάμει σημείων καὶ τεράτων,\FnParseFormGloss{a}{genitive neuter plural}{τέρας}{wonder} ἐν δυνάμει ⸀πνεύματος· ὥστε με ἀπὸ Ἰερουσαλὴμ καὶ κύκλῳ\FnFormGloss{b}{κύκλῳ}{all around} μέχρι\FnFormGloss{c}{μέχρι(ς)}{until} τοῦ Ἰλλυρικοῦ\FnParseFormGloss{d}{genitive neuter singular}{Ἰλλυρικόν}{Illyricum} πεπληρωκέναι\FnParse{e}{perfect active infinitive} τὸ εὐαγγέλιον τοῦ Χριστοῦ, 
\MainTextVerseMark{15}{20}οὕτως δὲ ⸀φιλοτιμούμενον\FnParseFormGloss{a}{present middle participle accusative masculine singular}{φιλοτιμέομαι}{to have an ambition} εὐαγγελίζεσθαι\FnParse{b}{present middle infinitive} οὐχ ὅπου ὠνομάσθη\FnParseFormGloss{c}{aorist passive indicative 3rd singular}{ὀνομάζω}{to give a name} Χριστός, ἵνα μὴ ἐπ’ ἀλλότριον\FnParseFormGloss{d}{accusative masculine singular}{ἀλλότριος}{belonging to another} θεμέλιον\FnParseFormGloss{e}{accusative masculine singular}{θεμέλιος}{foundation} οἰκοδομῶ,\FnParse{f}{present active subjunctive 1st singular} 
\MainTextVerseMark{15}{21}ἀλλὰ καθὼς γέγραπται·\FnParse{a}{perfect passive indicative 3rd singular} ⸂Οἷς οὐκ ἀνηγγέλη\FnParseFormGloss{b}{aorist passive indicative 3rd singular}{ἀναγγέλλω}{to tell} περὶ αὐτοῦ ὄψονται⸃,\FnParse{c}{future middle indicative 3rd plural} καὶ οἳ οὐκ ἀκηκόασιν\FnParse{d}{perfect active indicative 3rd plural} συνήσουσιν.\FnParseFormGloss{e}{future active indicative 3rd plural}{συνίημι}{to understand} 
\MainTextVerseMark{15}{22}Διὸ καὶ ἐνεκοπτόμην\FnParseFormGloss{a}{imperfect passive indicative 1st singular}{ἐγκόπτω}{to hinder} τὰ πολλὰ τοῦ ἐλθεῖν\FnParse{b}{aorist active infinitive} πρὸς ὑμᾶς· 
\MainTextVerseMark{15}{23}νυνὶ\FnFormGloss{a}{νυνί}{now} δὲ μηκέτι\FnFormGloss{b}{μηκέτι}{no longer} τόπον ἔχων\FnParse{c}{present active participle nominative masculine singular} ἐν τοῖς κλίμασι\FnParseFormGloss{d}{dative neuter plural}{κλίμα}{region} τούτοις, ἐπιποθίαν\FnParseFormGloss{e}{accusative feminine singular}{ἐπιποθία}{longing} δὲ ἔχων\FnParse{f}{present active participle nominative masculine singular} τοῦ ἐλθεῖν\FnParse{g}{aorist active infinitive} πρὸς ὑμᾶς ἀπὸ ⸀ἱκανῶν ἐτῶν, 
\MainTextVerseMark{15}{24}ὡς ⸀ἂν πορεύωμαι\FnParse{a}{present middle subjunctive 1st singular} εἰς τὴν ⸀Σπανίαν,\FnParseFormGloss{b}{accusative feminine singular}{Σπανία}{Spain} ἐλπίζω\FnParse{c}{present active indicative 1st singular} γὰρ διαπορευόμενος\FnParseFormGloss{d}{present middle participle nominative masculine singular}{διαπορεύομαι}{to go through} θεάσασθαι\FnParseFormGloss{e}{aorist middle infinitive}{θεάομαι}{to see} ὑμᾶς καὶ ὑφ’ ὑμῶν προπεμφθῆναι\FnParseFormGloss{f}{aorist passive infinitive}{προπέμπω}{to accompany} ἐκεῖ ἐὰν ὑμῶν πρῶτον ἀπὸ μέρους ἐμπλησθῶ—\FnParseFormGloss{g}{aorist passive subjunctive 1st singular}{ἐμπίμπλημι}{to fill} 
\MainTextVerseMark{15}{25}νυνὶ\FnFormGloss{a}{νυνί}{now} δὲ πορεύομαι\FnParse{b}{present middle indicative 1st singular} εἰς Ἰερουσαλὴμ διακονῶν\FnParse{c}{present active participle nominative masculine singular} τοῖς ἁγίοις. 
\MainTextVerseMark{15}{26}εὐδόκησαν\FnParseFormGloss{a}{aorist active indicative 3rd plural}{εὐδοκέω}{to be well pleased} γὰρ Μακεδονία\FnParseFormGloss{b}{nominative feminine singular}{Μακεδονία}{Macedonia} καὶ Ἀχαΐα\FnParseFormGloss{c}{nominative feminine singular}{Ἀχαΐα}{Achaia} κοινωνίαν\FnParseFormGloss{d}{accusative feminine singular}{κοινωνία}{fellowship} τινὰ ποιήσασθαι\FnParse{e}{aorist middle infinitive} εἰς τοὺς πτωχοὺς τῶν ἁγίων τῶν ἐν Ἰερουσαλήμ. 
\MainTextVerseMark{15}{27}εὐδόκησαν\FnParseFormGloss{a}{aorist active indicative 3rd plural}{εὐδοκέω}{to be well pleased} γάρ, καὶ ὀφειλέται\FnParseFormGloss{b}{nominative masculine plural}{ὀφειλέτης}{debtor} ⸂εἰσὶν\FnParse{c}{present active indicative 3rd plural} αὐτῶν⸃· εἰ γὰρ τοῖς πνευματικοῖς\FnParseFormGloss{d}{dative neuter plural}{πνευματικός}{spiritual} αὐτῶν ἐκοινώνησαν\FnParseFormGloss{e}{aorist active indicative 3rd plural}{κοινωνέω}{to share in} τὰ ἔθνη, ὀφείλουσιν\FnParse{f}{present active indicative 3rd plural} καὶ ἐν τοῖς σαρκικοῖς\FnParseFormGloss{g}{dative neuter plural}{σαρκικός}{material} λειτουργῆσαι\FnParseFormGloss{h}{aorist active infinitive}{λειτουργέω}{to perform religious duties} αὐτοῖς. 
\MainTextVerseMark{15}{28}τοῦτο οὖν ἐπιτελέσας,\FnParseFormGloss{a}{aorist active participle nominative masculine singular}{ἐπιτελέω}{to finish} καὶ σφραγισάμενος\FnParseFormGloss{b}{aorist middle participle nominative masculine singular}{σφραγίζω}{to seal} αὐτοῖς τὸν καρπὸν τοῦτον, ἀπελεύσομαι\FnParse{c}{future middle indicative 1st singular} δι’ ὑμῶν ⸀εἰς Σπανίαν·\FnParseFormGloss{d}{accusative feminine singular}{Σπανία}{Spain} 
\MainTextVerseMark{15}{29}οἶδα\FnParse{a}{perfect active indicative 1st singular} δὲ ὅτι ἐρχόμενος\FnParse{b}{present middle participle nominative masculine singular} πρὸς ὑμᾶς ἐν πληρώματι\FnParseFormGloss{c}{dative neuter singular}{πλήρωμα}{fullness} ⸀εὐλογίας\FnParseFormGloss{d}{genitive feminine singular}{εὐλογία}{blessing} Χριστοῦ ἐλεύσομαι.\FnParse{e}{future middle indicative 1st singular} 
\MainTextVerseMark{15}{30}Παρακαλῶ\FnParse{a}{present active indicative 1st singular} δὲ ὑμᾶς, ἀδελφοί, διὰ τοῦ κυρίου ἡμῶν Ἰησοῦ Χριστοῦ καὶ διὰ τῆς ἀγάπης τοῦ πνεύματος συναγωνίσασθαί\FnParseFormGloss{b}{aorist middle infinitive}{συναγωνίζομαι}{to join in a struggle} μοι ἐν ταῖς προσευχαῖς ὑπὲρ ἐμοῦ πρὸς τὸν θεόν, 
\MainTextVerseMark{15}{31}ἵνα ῥυσθῶ\FnParseFormGloss{a}{aorist passive subjunctive 1st singular}{ῥύομαι}{to rescue} ἀπὸ τῶν ἀπειθούντων\FnParseFormGloss{b}{present active participle genitive masculine plural}{ἀπειθέω}{to disobey} ἐν τῇ Ἰουδαίᾳ ⸀καὶ ἡ διακονία μου ἡ εἰς Ἰερουσαλὴμ εὐπρόσδεκτος\FnParseFormGloss{c}{nominative feminine singular}{εὐπρόσδεκτος}{acceptable} ⸂τοῖς ἁγίοις γένηται⸃,\FnParse{d}{aorist middle subjunctive 3rd singular} 
\MainTextVerseMark{15}{32}ἵνα ἐν χαρᾷ ⸀ἐλθὼν\FnParse{a}{aorist active participle nominative masculine singular} πρὸς ὑμᾶς διὰ θελήματος ⸀θεοῦ συναναπαύσωμαι\FnParseFormGloss{b}{aorist middle subjunctive 1st singular}{συναναπαύομαι}{to find rest together} ὑμῖν. 
\MainTextVerseMark{15}{33}ὁ δὲ θεὸς τῆς εἰρήνης μετὰ πάντων ὑμῶν· ἀμήν. 
\ChapterHeader{16}{Chapter 16}

\MainTextVerseMark{16}{1}Συνίστημι\FnParseFormGloss{a}{present active indicative 1st singular}{συνίστημι}{to commend} δὲ ὑμῖν Φοίβην\FnParseFormGloss{b}{accusative feminine singular}{Φοίβη}{Phoebe} τὴν ἀδελφὴν\FnParseFormGloss{c}{accusative feminine singular}{ἀδελφή}{sister} ἡμῶν, οὖσαν\FnParse{d}{present active participle accusative feminine singular} ⸀καὶ διάκονον\FnParseFormGloss{e}{accusative feminine singular}{διάκονος}{servant} τῆς ἐκκλησίας τῆς ἐν Κεγχρεαῖς,\FnParseFormGloss{f}{dative feminine plural}{Κεγχρεαί}{Cenchrea} 
\MainTextVerseMark{16}{2}ἵνα ⸂αὐτὴν προσδέξησθε⸃\FnParseFormGloss{a}{aorist middle subjunctive 2nd plural}{προσδέχομαι}{to receive} ἐν κυρίῳ ἀξίως\FnFormGloss{b}{ἀξίως}{worthily} τῶν ἁγίων, καὶ παραστῆτε\FnParse{c}{aorist active subjunctive 2nd plural} αὐτῇ ἐν ᾧ ἂν ὑμῶν χρῄζῃ\FnParseFormGloss{d}{present active subjunctive 3rd singular}{χρῄζω}{to need} πράγματι,\FnParseFormGloss{e}{dative neuter singular}{πρᾶγμα}{thing} καὶ γὰρ αὐτὴ προστάτις\FnParseFormGloss{f}{nominative feminine singular}{προστάτις}{helper} πολλῶν ἐγενήθη\FnParse{g}{aorist passive indicative 3rd singular} καὶ ⸂ἐμοῦ αὐτοῦ⸃. 
\MainTextVerseMark{16}{3}Ἀσπάσασθε\FnParse{a}{aorist middle imperative 2nd plural} Πρίσκαν\FnParseFormGloss{b}{accusative feminine singular}{Πρίσκα}{Prisca} καὶ Ἀκύλαν\FnParseFormGloss{c}{accusative masculine singular}{Ἀκύλας}{Aquila} τοὺς συνεργούς\FnParseFormGloss{d}{accusative masculine plural}{συνεργός}{fellow worker} μου ἐν Χριστῷ Ἰησοῦ, 
\MainTextVerseMark{16}{4}οἵτινες ὑπὲρ τῆς ψυχῆς μου τὸν ἑαυτῶν τράχηλον\FnParseFormGloss{a}{accusative masculine singular}{τράχηλος}{neck} ὑπέθηκαν,\FnParseFormGloss{b}{aorist active indicative 3rd plural}{ὑποτίθημι}{to risk} οἷς οὐκ ἐγὼ μόνος εὐχαριστῶ\FnParse{c}{present active indicative 1st singular} ἀλλὰ καὶ πᾶσαι αἱ ἐκκλησίαι τῶν ἐθνῶν, 
\MainTextVerseMark{16}{5}καὶ τὴν κατ’ οἶκον αὐτῶν ἐκκλησίαν. ἀσπάσασθε\FnParse{a}{aorist middle imperative 2nd plural} Ἐπαίνετον\FnParseFormGloss{b}{accusative masculine singular}{Ἐπαίνετος}{Epenetus} τὸν ἀγαπητόν μου, ὅς ἐστιν\FnParse{c}{present active indicative 3rd singular} ἀπαρχὴ\FnParseFormGloss{d}{nominative feminine singular}{ἀπαρχή}{firstfruits} τῆς ⸀Ἀσίας\FnParseFormGloss{e}{genitive feminine singular}{Ἀσία}{Asia} εἰς Χριστόν. 
\MainTextVerseMark{16}{6}ἀσπάσασθε\FnParse{a}{aorist middle imperative 2nd plural} ⸀Μαριάμ,\FnParseFormGloss{b}{accusative feminine singular}{Μαριάμ}{Mary} ἥτις πολλὰ ἐκοπίασεν\FnParseFormGloss{c}{aorist active indicative 3rd singular}{κοπιάω}{to work} εἰς ⸀ὑμᾶς. 
\MainTextVerseMark{16}{7}ἀσπάσασθε\FnParse{a}{aorist middle imperative 2nd plural} Ἀνδρόνικον\FnParseFormGloss{b}{accusative masculine singular}{Ἀνδρόνικος}{Andronicus} καὶ ⸀Ἰουνίαν\FnParseFormGloss{c}{accusative feminine singular}{Ἰουνία}{Junia} τοὺς συγγενεῖς\FnParseFormGloss{d}{accusative masculine plural}{συγγενής}{family} μου καὶ συναιχμαλώτους\FnParseFormGloss{e}{accusative masculine plural}{συναιχμάλωτος}{fellow prisoner} μου, οἵτινές εἰσιν\FnParse{f}{present active indicative 3rd plural} ἐπίσημοι\FnParseFormGloss{g}{nominative masculine plural}{ἐπίσημος}{notorious} ἐν τοῖς ἀποστόλοις, οἳ καὶ πρὸ ἐμοῦ ⸀γέγοναν\FnParse{h}{perfect active indicative 3rd plural} ἐν Χριστῷ. 
\MainTextVerseMark{16}{8}ἀσπάσασθε\FnParse{a}{aorist middle imperative 2nd plural} ⸀Ἀμπλιᾶτον\FnParseFormGloss{b}{accusative masculine singular}{Ἀμπλιᾶτος}{Ampliatus} τὸν ἀγαπητόν μου ἐν κυρίῳ. 
\MainTextVerseMark{16}{9}ἀσπάσασθε\FnParse{a}{aorist middle imperative 2nd plural} Οὐρβανὸν\FnParseFormGloss{b}{accusative masculine singular}{Οὐρβανός}{Urbanus} τὸν συνεργὸν\FnParseFormGloss{c}{accusative masculine singular}{συνεργός}{fellow worker} ἡμῶν ἐν Χριστῷ καὶ Στάχυν\FnParseFormGloss{d}{accusative masculine singular}{Στάχυς}{Stachys} τὸν ἀγαπητόν μου. 
\MainTextVerseMark{16}{10}ἀσπάσασθε\FnParse{a}{aorist middle imperative 2nd plural} Ἀπελλῆν\FnParseFormGloss{b}{accusative masculine singular}{Ἀπελλῆς}{Apelles} τὸν δόκιμον\FnParseFormGloss{c}{accusative masculine singular}{δόκιμος}{approved by testing} ἐν Χριστῷ. ἀσπάσασθε\FnParse{d}{aorist middle imperative 2nd plural} τοὺς ἐκ τῶν Ἀριστοβούλου.\FnParseFormGloss{e}{genitive masculine singular}{Ἀριστόβουλος}{Aristobulus} 
\MainTextVerseMark{16}{11}ἀσπάσασθε\FnParse{a}{aorist middle imperative 2nd plural} Ἡρῳδίωνα\FnParseFormGloss{b}{accusative masculine singular}{Ἡρῳδίων}{Herodion} τὸν ⸀συγγενῆ\FnParseFormGloss{c}{accusative masculine singular}{συγγενής}{family} μου. ἀσπάσασθε\FnParse{d}{aorist middle imperative 2nd plural} τοὺς ἐκ τῶν Ναρκίσσου\FnParseFormGloss{e}{genitive masculine singular}{Νάρκισσος}{Narcissus} τοὺς ὄντας\FnParse{f}{present active participle accusative masculine plural} ἐν κυρίῳ. 
\MainTextVerseMark{16}{12}ἀσπάσασθε\FnParse{a}{aorist middle imperative 2nd plural} Τρύφαιναν\FnParseFormGloss{b}{accusative feminine singular}{Τρύφαινα}{Tryphena} καὶ Τρυφῶσαν\FnParseFormGloss{c}{accusative feminine singular}{Τρυφῶσα}{Tryphosa} τὰς κοπιώσας\FnParseFormGloss{d}{present active participle accusative feminine plural}{κοπιάω}{to work} ἐν κυρίῳ. ἀσπάσασθε\FnParse{e}{aorist middle imperative 2nd plural} Περσίδα\FnParseFormGloss{f}{accusative feminine singular}{Περσίς}{Persis} τὴν ἀγαπητήν, ἥτις πολλὰ ἐκοπίασεν\FnParseFormGloss{g}{aorist active indicative 3rd singular}{κοπιάω}{to work} ἐν κυρίῳ. 
\MainTextVerseMark{16}{13}ἀσπάσασθε\FnParse{a}{aorist middle imperative 2nd plural} Ῥοῦφον\FnParseFormGloss{b}{accusative masculine singular}{Ῥοῦφος}{Rufus} τὸν ἐκλεκτὸν\FnParseFormGloss{c}{accusative masculine singular}{ἐκλεκτός}{elect} ἐν κυρίῳ καὶ τὴν μητέρα αὐτοῦ καὶ ἐμοῦ. 
\MainTextVerseMark{16}{14}ἀσπάσασθε\FnParse{a}{aorist middle imperative 2nd plural} Ἀσύγκριτον,\FnParseFormGloss{b}{accusative masculine singular}{Ἀσύγκριτος}{Asyncritus} Φλέγοντα,\FnParseFormGloss{c}{accusative masculine singular}{Φλέγων}{Phlegon} ⸂Ἑρμῆν,\FnParseFormGloss{d}{accusative masculine singular}{Ἑρμῆς}{Hermes} Πατροβᾶν,\FnParseFormGloss{e}{accusative masculine singular}{Πατροβᾶς}{Patrobas} Ἑρμᾶν⸃\FnParseFormGloss{f}{accusative masculine singular}{Ἑρμᾶς}{Hermas} καὶ τοὺς σὺν αὐτοῖς ἀδελφούς. 
\MainTextVerseMark{16}{15}ἀσπάσασθε\FnParse{a}{aorist middle imperative 2nd plural} Φιλόλογον\FnParseFormGloss{b}{accusative masculine singular}{Φιλόλογος}{Philologus} καὶ Ἰουλίαν,\FnParseFormGloss{c}{accusative feminine singular}{Ἰουλία}{Julia} Νηρέα\FnParseFormGloss{d}{accusative masculine singular}{Νηρεύς}{Nereus} καὶ τὴν ἀδελφὴν\FnParseFormGloss{e}{accusative feminine singular}{ἀδελφή}{sister} αὐτοῦ, καὶ Ὀλυμπᾶν\FnParseFormGloss{f}{accusative masculine singular}{Ὀλυμπᾶς}{Olympas} καὶ τοὺς σὺν αὐτοῖς πάντας ἁγίους. 
\MainTextVerseMark{16}{16}Ἀσπάσασθε\FnParse{a}{aorist middle imperative 2nd plural} ἀλλήλους ἐν φιλήματι\FnParseFormGloss{b}{dative neuter singular}{φίλημα}{kiss} ἁγίῳ. Ἀσπάζονται\FnParse{c}{present middle indicative 3rd plural} ὑμᾶς αἱ ἐκκλησίαι ⸀πᾶσαι τοῦ Χριστοῦ. 
\MainTextVerseMark{16}{17}Παρακαλῶ\FnParse{a}{present active indicative 1st singular} δὲ ὑμᾶς, ἀδελφοί, σκοπεῖν\FnParseFormGloss{b}{present active infinitive}{σκοπέω}{to watch out for} τοὺς τὰς διχοστασίας\FnParseFormGloss{c}{accusative feminine plural}{διχοστασία}{division} καὶ τὰ σκάνδαλα\FnParseFormGloss{d}{accusative neuter plural}{σκάνδαλον}{stumbling block} παρὰ τὴν διδαχὴν ἣν ὑμεῖς ἐμάθετε\FnParseFormGloss{e}{aorist active indicative 2nd plural}{μανθάνω}{to learn} ποιοῦντας,\FnParse{f}{present active participle accusative masculine plural} καὶ ⸀ἐκκλίνετε\FnParseFormGloss{g}{present active imperative 2nd plural}{ἐκκλίνω}{to turn away} ἀπ’ αὐτῶν· 
\MainTextVerseMark{16}{18}οἱ γὰρ τοιοῦτοι τῷ κυρίῳ ⸀ἡμῶν Χριστῷ οὐ δουλεύουσιν\FnParseFormGloss{a}{present active indicative 3rd plural}{δουλεύω}{to serve} ἀλλὰ τῇ ἑαυτῶν κοιλίᾳ,\FnParseFormGloss{b}{dative feminine singular}{κοιλία}{any and all internal organs} καὶ διὰ τῆς χρηστολογίας\FnParseFormGloss{c}{genitive feminine singular}{χρηστολογία}{smooth talk} καὶ εὐλογίας\FnParseFormGloss{d}{genitive feminine singular}{εὐλογία}{blessing} ἐξαπατῶσι\FnParseFormGloss{e}{present active indicative 3rd plural}{ἐξαπατάω}{to deceive} τὰς καρδίας τῶν ἀκάκων.\FnParseFormGloss{f}{genitive masculine plural}{ἄκακος}{blameless} 
\MainTextVerseMark{16}{19}ἡ γὰρ ὑμῶν ὑπακοὴ\FnParseFormGloss{a}{nominative feminine singular}{ὑπακοή}{obedience} εἰς πάντας ἀφίκετο·\FnParseFormGloss{b}{aorist middle indicative 3rd singular}{ἀφικνέομαι}{to reach} ⸂ἐφ’ ὑμῖν οὖν χαίρω⸃,\FnParse{c}{present active indicative 1st singular} θέλω\FnParse{d}{present active indicative 1st singular} δὲ ὑμᾶς ⸀σοφοὺς\FnParseFormGloss{e}{accusative masculine plural}{σοφός}{wise} εἶναι\FnParse{f}{present active infinitive} εἰς τὸ ἀγαθόν, ἀκεραίους\FnParseFormGloss{g}{accusative masculine plural}{ἀκέραιος}{innocent} δὲ εἰς τὸ κακόν. 
\MainTextVerseMark{16}{20}ὁ δὲ θεὸς τῆς εἰρήνης συντρίψει\FnParseFormGloss{a}{future active indicative 3rd singular}{συντρίβω}{to break} τὸν Σατανᾶν ὑπὸ τοὺς πόδας ὑμῶν ἐν τάχει.\FnParseFormGloss{b}{dative neuter singular}{τάχος}{quickness} ἡ χάρις τοῦ κυρίου ἡμῶν Ἰησοῦ ⸀Χριστοῦ μεθ’ ὑμῶν. 
\MainTextVerseMark{16}{21}⸀Ἀσπάζεται\FnParse{a}{present middle indicative 3rd singular} ὑμᾶς Τιμόθεος\FnParseFormGloss{b}{nominative masculine singular}{Τιμόθεος}{Timothy} ὁ συνεργός\FnParseFormGloss{c}{nominative masculine singular}{συνεργός}{fellow worker} μου, καὶ Λούκιος\FnParseFormGloss{d}{nominative masculine singular}{Λούκιος}{Lucius} καὶ Ἰάσων\FnParseFormGloss{e}{nominative masculine singular}{Ἰάσων}{Jason} καὶ Σωσίπατρος\FnParseFormGloss{f}{nominative masculine singular}{Σωσίπατρος}{Sosipater} οἱ συγγενεῖς\FnParseFormGloss{g}{nominative masculine plural}{συγγενής}{family} μου. 
\MainTextVerseMark{16}{22}ἀσπάζομαι\FnParse{a}{present middle indicative 1st singular} ὑμᾶς ἐγὼ Τέρτιος\FnParseFormGloss{b}{nominative masculine singular}{Τέρτιος}{Tertius} ὁ γράψας\FnParse{c}{aorist active participle nominative masculine singular} τὴν ἐπιστολὴν\FnParseFormGloss{d}{accusative feminine singular}{ἐπιστολή}{letter} ἐν κυρίῳ. 
\MainTextVerseMark{16}{23}ἀσπάζεται\FnParse{a}{present middle indicative 3rd singular} ὑμᾶς Γάϊος\FnParseFormGloss{b}{nominative masculine singular}{Γάϊος}{Gaius} ὁ ξένος\FnParseFormGloss{c}{nominative masculine singular}{ξένος}{strange} μου καὶ ⸂ὅλης τῆς ἐκκλησίας⸃. ἀσπάζεται\FnParse{d}{present middle indicative 3rd singular} ὑμᾶς Ἔραστος\FnParseFormGloss{e}{nominative masculine singular}{Ἔραστος}{Erastus} ὁ οἰκονόμος\FnParseFormGloss{f}{nominative masculine singular}{οἰκονόμος}{manager} τῆς πόλεως καὶ Κούαρτος\FnParseFormGloss{g}{nominative masculine singular}{Κούαρτος}{Quartus} ὁ ἀδελφός. 
\MainTextVerseMark{16}{24}⸂Ἡ χάρις τοῦ κυρίου ἡμῶν Ἰησοῦ Χριστοῦ μετὰ πάντων ὑμῶν.⸃ ⸀Ἀμήν. 
\end{document}